% =====================================================================
%  PoP -> IEEE TDSC : 正文（英文）
%  由 paper/tdsc-skeleton.tex 扩写而来；写作规格见 docs/tdsc-writing-plan.md
%
%  用法：本文件是**在写的正文**，不是成稿。
%  所有 \need{} 与 \fixme{} 处都要填；提交前全部清掉。
%
%  编译校验：
%      pdflatex "\def\SKELQUIET{}% =====================================================================
%  PoP -> IEEE TDSC : 正文（英文）
%  由 paper/tdsc-skeleton.tex 扩写而来；写作规格见 docs/tdsc-writing-plan.md
%
%  用法：本文件是**在写的正文**，不是成稿。
%  所有 \need{} 与 \fixme{} 处都要填；提交前全部清掉。
%
%  编译校验：
%      pdflatex "\def\SKELQUIET{}% =====================================================================
%  PoP -> IEEE TDSC : 正文（英文）
%  由 paper/tdsc-skeleton.tex 扩写而来；写作规格见 docs/tdsc-writing-plan.md
%
%  用法：本文件是**在写的正文**，不是成稿。
%  所有 \need{} 与 \fixme{} 处都要填；提交前全部清掉。
%
%  编译校验：
%      pdflatex "\def\SKELQUIET{}% =====================================================================
%  PoP -> IEEE TDSC : 正文（英文）
%  由 paper/tdsc-skeleton.tex 扩写而来；写作规格见 docs/tdsc-writing-plan.md
%
%  用法：本文件是**在写的正文**，不是成稿。
%  所有 \need{} 与 \fixme{} 处都要填；提交前全部清掉。
%
%  编译校验：
%      pdflatex "\def\SKELQUIET{}\input{tdsc-paper.tex}"   % 两遍，第二遍收交叉引用
%  \SKELQUIET 让所有中文标记宏吞掉参数；不带它则需要能排版中文的引擎
%  （xelatex + CJK 字体）。正文本身全英文，标记里才有中文。
%
%  ⚠️ IEEEtran.cls 需要装在 TeX 发行版里。WSL 侧已装进
%     /home/dong/texlive/texmf-local/tex/latex/ieeetran/（2026-10-12）。
% =====================================================================
\documentclass[journal]{IEEEtran}

\usepackage{amsmath,amssymb,amsthm}
\usepackage{booktabs}
\usepackage{array}
\usepackage{graphicx}
\usepackage{xcolor}
\usepackage{hyperref}

% 正文用 Times（ptm，窄），\texttt 用 Courier（pcr，等宽且宽）。
% 一行里只要有一个长 Courier 标识符（\texttt{normalized\_keyword\_block}
% 这类）就几乎没有可拉伸的胶，TeX 只能拉出 badness 10000 的大空隙。
% emergencystretch 允许第三遍在这些行上额外拉伸，把最坏情况吸收掉；
% 代价是那些行略松，远好于一道 10000 的豁口。
% （2026-10-12：装上真正的 Times 后此类行从 3 涨到 12 才暴露。）
\setlength{\emergencystretch}{2em}

% microtype：字形微伸缩（pdftex 下 ±2%）+ 标点悬挂。
% emergencystretch 只能吸收「可拉伸胶不足」，但 badness 在 10000 处封顶，
% 段首被后面长片段（consequences---customer-facing、normalized\_keyword\_block）
% 挤空的行照样报 10000。microtype 从字形宽度上补，是这一类的标准解法。
\usepackage{microtype}

% 窄栏里用 p{} 会因两端对齐产生大片字距（编译时刷一堆 Underfull \hbox）。
% 表格正文列一律用 L{}：左对齐、允许自然参差。
\newcolumntype{L}[1]{>{\raggedright\arraybackslash}p{#1}}
\newcommand{\tabfont}{\footnotesize}

\newtheorem{theorem}{Theorem}
\newtheorem{lemma}[theorem]{Lemma}
\newtheorem{proposition}[theorem]{Proposition}
\newtheorem{corollary}[theorem]{Corollary}
\newtheorem{definition}[theorem]{Definition}
\newtheorem{assumption}[theorem]{Assumption}
\newtheorem{remark}[theorem]{Remark}

% --- 工作标记：提交前全部清掉 ---
% 注意：宏名不能叫 \wp —— \wp 已是 LaTeX/amssymb 的 Weierstrass ℘，
% \newcommand{\wp} 会直接报 "Command \wp already defined" 并使后文级联报错。
\ifdefined\SKELQUIET
  \newcommand{\wpmark}[1]{}
  \newcommand{\need}[1]{}
  \newcommand{\fixme}[1]{}
\else
  \newcommand{\wpmark}[1]{\textcolor{red}{\textsuperscript{[WP#1]}}}
  \newcommand{\need}[1]{\textcolor{blue}{[[ NEED: #1 ]]}}
  \newcommand{\fixme}[1]{\textcolor{orange}{[[ FIX: #1 ]]}}
\fi

\begin{document}

% ---------------------------------------------------------------------
% 标题：决策 D5 —— 保留 "Zero-Knowledge" 的前提是 groth16 成为交付档 (WP1)。
%       若 WP1 未打通 groth16，改用下面被注释掉的那个标题，
%       并同步删掉摘要与 §8 里的 zero-knowledge 措辞。
% ---------------------------------------------------------------------
\title{Proof-of-Policy: Zero-Knowledge Compliance Certificates\\ for LLM Agent Outputs}
% 备选标题（若 groth16 未打通）：
% \title{Proof-of-Policy: Succinct Compliance Certificates for LLM Agent Outputs}

\author{\need{authors}}
\maketitle

% ---------------------------------------------------------------------
\begin{abstract}
% 决策 E2：内容隐私上界必须出现在摘要，不能埋进 §Limitations。
% 决策 D1：贡献陈述按"分析 x 系统各半"写 —— 分析在前，系统是分析所需的那一实例。
Proof-of-Policy (PoP) lets an auditor check that an LLM agent's response satisfies
a published policy, without trusted hardware and without the operator disclosing
the response in the clear. Our contribution is analytical before it is
constructive. We give a quantitative bound on what such a certificate binds---a
guessing-entropy term that no additional computation can improve---together with
the structural assumption the bound requires, and a proof that the assumption is
necessary rather than convenient. The same accounting governs the request side,
where a certificate without receiver-supplied freshness is replayable without
limit. PoP is then presented as the instance these bounds require---a policy
compiler, a
% 决策 D5 条件措辞：若 groth16 未成为交付档，把下行的 "zero-knowledge" 换成 "succinct"，
% 并同步改标题与 §8。
zero-knowledge circuit, and a signed certificate layer---measured end to end
rather than asserted, with the leak profile of the certificate reported as a
result in its own right.

\smallskip
\noindent\textbf{Scope of privacy.} PoP provides \emph{policy}-zero-knowledge:
the certificate reveals neither the policy evaluation trace nor, in
policy-private mode, the policy itself. It does \emph{not} guarantee that $T$
is unrecoverable: as we show in \S\ref{sec:binding}, any publicly recomputable
binding to $T$ admits a re-execution attack whenever the adversary knows the
prompt and the model. The same bit-budget governs the request side: a
certificate is replayable without limit unless the receiver contributes fresh
randomness (Theorem \ref{thm:t3}).
\end{abstract}

\begin{IEEEkeywords}
zero-knowledge proofs, LLM agents, policy compliance, verifiable computation,
accountability
\end{IEEEkeywords}

% =====================================================================
\section{Introduction}
\label{sec:intro}

LLM agents increasingly take actions that carry con\-se\-quences---\allowbreak{}
customer-facing replies, tool invocations, financial and medical content. Regulation such as the
EU AI Act (Articles~12 and~13) and enterprise governance both require that
compliance be auditable. In practice it can only be taken on trust: logs can be
forged, guardrail claims cannot be checked, and content privacy forbids retaining
the plaintext that would settle the question.

Four lines of work each occupy one corner of this problem without spanning it.
\emph{TEE execution proofs}---Proof-of-Guardrail and its relatives---show that
guardrail code was executed and that inputs and outputs were signed; as that work
itself observes, \emph{execution is not compliance}, and the approach carries
roughly $34\%$ latency overhead together with a hardware dependency.
\emph{TEE governance stacks} (Agentrust, cMCP, cA2A) inherit both the hardware
and the vendor facility. \emph{Formal verification} (Lean-Agent and similar)
requires the entire policy to be formalised by hand, and does not apply to
black-box responses. \emph{ZK regex and PII primitives} (zk-regex~V2, verifiable
document redaction \cite{vdr}, ZK Email) are building blocks; none of them is a
policy-level system.

The nearest work, WITNESS, supports two policy primitives, has no peer-reviewed
paper, and---by its own project page---ships an Apache-2.0 reference
implementation, so ``open source'' is not a differentiator; it also offers no
content privacy. \texttt{@brivora/verify} is not zero-knowledge, and DeepProve
proves the \emph{model computation} layer rather than the \emph{policy
compliance} layer. The combination of a general policy DSL compiled to ZK
constraints, \emph{without} a TEE, with selective disclosure of response content
and a formal security model, is an explicit vacancy. This paper occupies it with
PoP---but the system is not the contribution we lead with. Leading with a system
invites the objection that a signature would do; leading with the analysis is
what makes the system necessary, and it is the analysis that this paper is
organised around.

\need{⚠️ 上面点名的 Proof-of-Guardrail / Agentrust,cMCP,cA2A / Lean-Agent /
zk-regex V2 / ZK Email / WITNESS / \texttt{@brivora/verify} / DeepProve，
以及 EU AI Act Art.~12/13，\emph{均未给 \cite{}}：条目尚未逐条核实
（docs/tdsc-fixlist.md \S2、WP8）。核实前不得凭空写 bibitem，故此处只留名。}

\smallskip
\noindent\textbf{What we do \emph{not} claim.} We do not claim that policy proofs
bind nothing: our own system already binds the response through a verifier-issued
nonce (Definition~\ref{def:system}). We do not claim to be the first to observe
that a proof can fail to bind the artifact it purports to describe: the
\emph{effort gap} for ZK inference proofs is formalised by Hollow-LLM
\cite{hollowllm}, and the \emph{source-integrity gap} for predicate proofs is
named by zk-attach \cite{zkattach}. And we do not claim the separation theorem
that our designated-verifier interface makes tempting (``public verification is
impossible when $k$ is small''); Remark~\ref{rem:designated} gives the reason,
and Corollary~\ref{cor:reexec} is the short version of it.

\smallskip
\noindent\textbf{Contributions.}
\begin{enumerate}
  \item We isolate the gap that survives in every design we surveyed---\emph{divergent
  delivery}---and show that neither hardware nor any proof closes it, while having
  the receiver contribute to the proven statement does
  (Proposition~\ref{thm:gap}).
  \item We give a quantitative bound on what a public binding reveals, together
  with the \emph{structural hypothesis} that bound needs, and a demonstration
  that the hypothesis is necessary: our own two public values sit on opposite
  sides of it (Theorem~\ref{thm:t1}, Remark~\ref{rem:betanecessary}). The
  governing quantity $k$ is measurable rather than assumed.
  \item We price the fix. A certificate is replayable without limit unless the
  receiver contributes fresh randomness, and with an empty nonce the adversary's
  advantage is $1$ regardless of every other parameter of the system
  (Theorem~\ref{thm:t3}).
  \item We describe PoP, a system that realises this design without a TEE
  (\S\ref{sec:system}), and we report what its certificates leak about the policy
  they certify (\S\ref{sec:private-policy}).
\end{enumerate}

\smallskip
\noindent\textbf{One spine, two sides.} Theorems~\ref{thm:t1} and~\ref{thm:t3}
state the same fact about two different inputs: \emph{a proof binds only as much
as the committed, adversary-non-selectable quantity has bits for}. The response
side is governed by entropy, $k = -\log_2 \mathsf{Guess}(T \mid Z)$
(Definition~\ref{def:k}); the request side by freshness, $f$. Neither can be
bought with more computation.

\fixme{全文禁止出现"比 Proof-of-Guardrail 更强"或任何强弱对比。
Proof-of-Guardrail 与本文正交：它证"计算发生过"，本文证"性质成立"。}

% =====================================================================
\section{System and Threat Model}
\label{sec:system}

The object of study is a \emph{policy-proof system}: a pipeline that compiles a
policy written in a small DSL into a circuit, evaluates an agent's response
against that circuit inside a zkVM, and emits a certificate binding the verdict
to the response. This section fixes the object and its adversary before any
theorem is stated.

\subsection{The system}
\label{sec:system-overview}

\paragraph{Pipeline and the cross-layer contract.} A policy package is compiled
to a \texttt{ConstraintSpec}---a JSON artefact carrying the compiled NFA---which
is the \emph{only} object crossing the layer boundary: the Python reference
implementation and the Rust/zkVM guest consume the same bytes, so cross-layer
drift is structurally impossible rather than merely tested against. The guest
runs a deterministic evaluator and commits its outcome. Of eight rule classes,
seven are decided \emph{inside} the circuit
(\texttt{keyword\_\allowbreak block},
\texttt{normalized\_\allowbreak keyword\_\allowbreak block},
\texttt{length\_\allowbreak bound}, \texttt{pattern\_\allowbreak block},
\texttt{format\_\allowbreak check}, \texttt{tool\_\allowbreak arg\_\allowbreak guard},
\texttt{budget\_\allowbreak bound}). The eighth,
\texttt{semantic\_bound}---a rule of the form ``harm probability $\le$
threshold''---is precisely the kind for which no formal proof can be learned, and
the circuit does not decide it. The circuit registers only \emph{that the rule
was delegated}, in the public value \texttt{delegated}, and a verifier must
\emph{conjunctively} check a companion \texttt{ezkl}/halo2 proof. Verifying the
zkVM side alone does not establish compliance, and we say so here because it is
the most likely misreading of the certificate.

\paragraph{What the verdict is bound to.} Rules that reason about the agent's
\emph{tool calls} cannot be judged against a list the prover fills in itself.
Receipts are therefore issued by a tool gateway after each call and chained
(\texttt{seq}, \texttt{prev}, signature); the circuit checks the chain's
\emph{structure}---entries consecutive, mutually consistent, each digest
recomputed from its own content---and fails closed to \texttt{trace\_unbound}
otherwise. Signature verification happens \emph{off circuit} (Ed25519 inside the
zkVM is prohibitively expensive), a boundary we state rather than gloss: the
circuit is blind to a forgery in the chain's final element, and two independent
off-circuit gates---digest comparison and signature verification---carry that
obligation, neither substituting for the other. Truncating the chain's tail is a
separate and weaker opening, closed by a session-end \texttt{ToolSeal}
commitment from the gateway, and it remains conditional on the verifier holding
the gateway's public key. \S\ref{sec:discussion} develops these residual
boundaries; they are not incidental to the design.

\paragraph{Dual privacy modes.} In \emph{transparent} mode the response and the
verdict are public. In \emph{private} mode the certificate publishes a response
commitment, a commitment per violation, and a redaction proof
(\texttt{redaction\_ok} together with \texttt{mask\_covered}, the latter
requiring the masked positions to fall inside a span the circuit verified to be a
genuine complete match). The guarantee is that \emph{public values contain no
plaintext}---not that $T$ cannot be recovered. A commitment is a public function
of $T$ against which anyone can guess and check, while the response binding must
be publicly recomputable for a verifier to check the string it was delivered.
Those two requirements are in direct tension, and \S\ref{sec:binding} quantifies
exactly when.

\paragraph{Certificates and anchoring.} A certificate is signed with a DSSE
envelope under Ed25519 with the key held by the issuing party, and its digest is
appended to a hash-chain anchor ledger. Two backends share one interface: a file
ledger, verifiable offline, and an on-chain registration storing only a 32-byte
digest. The transaction hash and on-chain timestamp are written back into the
local entry, so the offline chain and the public timestamp cross-check each
other and tampering with either is detectable.

\paragraph{Framework integration.} One gateway serves LangChain callbacks,
LangGraph (\texttt{guard\_node} plus full event streams), and MCP tool argument
and response checks, so trace certificates and proofs share a single chain of
evidence. The implementation passes 801 tests with 15 skipped; the skipped set is
enumerated in Appendix~\ref{app:repro}.

\begin{definition}[Policy-proof system]
\label{def:system}
A policy-proof system is a tuple $(\mathsf{Prove}, \mathsf{Ver}, \pi, \mathcal{M})$.
On challenge $n$ from a verifier, $\mathsf{Prove}(\pi, T, n)$ outputs a
certificate $c$ whose public values include
\begin{multline*}
b \;=\; H(\texttt{pop-bind-v1} \,\|\, \mathsf{u32be}(|n|) \,\|\, n \,\|\, T),\\
C = H(T).
\end{multline*}
$\mathsf{Ver}(c, n, T')$ accepts iff the policy relation holds \emph{and}
$b$ recomputed from the received $T'$ matches.
\end{definition}

\fixme{\textbf{H3/H4}：早先写的 ``the relation contains no term referring to the
delivered bytes'' 与仓库事实冲突——公开值里\emph{已有} $b$，验证时由送达的
$T'$ 现场重算（\texttt{circuits/types/src/lib.rs:444}、
\texttt{docs/security-model.md:232}）。命题已按上面的定义改写。见
\texttt{docs/tdsc-fixlist.md} \S H。}

\subsection{Threat model}
\label{sec:threat}

The adversary is the \emph{agent operator}: the party that holds the policy, runs
the model, and produces the certificate. The verifier is a regulator, auditor, or
third party holding only public values. The policy author issues $\pi$; in this
section we treat it as honest, and \S\ref{sec:wp4-tension} makes that assumption
explicit and scopes it, because in policy-private mode its honesty is doing real
work and cannot be waved through. We state the goals, and more importantly the
non-goals, before the theorems.

\begin{assumption}[Adversary side information]
\label{as:Z}
The adversary knows the prompt, the model (possibly its weights), and the
decoding parameters. \emph{This is the assumption that makes the theorem
non-trivial}: under greedy decoding, $k \approx 0$.
\end{assumption}

\paragraph{Goals.}
\begin{description}
  \item[(G1)] \emph{Soundness}---a response $T$ violating $\pi$ cannot yield an
  accepted certificate.
  \item[(G2)] \emph{Committed privacy}---in private mode the public values
  contain no plaintext of $T$. This is explicitly \emph{not} the claim that $T$
  is unrecoverable: once the guessing entropy $k$ of Definition~\ref{def:k} falls
  below $\log q + \log(1/\varepsilon)$, a publicly recomputable binding and
  content hiding are mutually exclusive (Theorem~\ref{thm:t1}).
  \item[(G3)] \emph{Traceable binding}---policy version, circuit, and proof
  artefact are bound to one another and recomputed at verification.
  \item[(G4)] \emph{Selective disclosure}---redaction and evidence opening are
  themselves provable.
  \item[(G5)] \emph{Record integrity}---tampering with, or reordering, the audit
  ledger is detectable.
\end{description}

\paragraph{Non-goals.} We do not prove the model's inference, and we do not cover
training data or model cards. (G1) is modulo the soundness of the zkVM and of the
hash function; we use those primitives and do not re-derive them.

\paragraph{Where the policy-private case is developed.} The three-party model,
the issuing prerequisite (P-PP), and the table of what each verifier may and may
not conclude are specific to policy-private mode, where the policy itself is
secret. They are stated once, in \S\ref{sec:wp4-tension}, and are deliberately
\emph{not} duplicated here---at a twelve-page limit they cannot be. The general
adversary of this section is the same operator, viewed without the
policy-secrecy refinement.

% =====================================================================
\section{Binding of Policy Proofs}
\label{sec:binding}

A policy proof is meant to bind a response: a verifier who accepts a
certificate should conclude that \emph{this} string was checked against the
policy. This section asks what that binding is actually worth. We first isolate
the gap that survives in every design we surveyed (\S\ref{sec:gap}), then give
a quantitative bound on what a public value reveals, together with the
structural hypothesis that bound needs (\S\ref{sec:tradeoff}), and finally
state the same fact on the request side, where the governing resource is
freshness rather than entropy (\S\ref{sec:t3}).

\subsection{The gap}
\label{sec:gap}

\begin{proposition}[Divergent delivery]
\label{thm:gap}
Definition \ref{def:system} binds $T'$ \emph{as received by the verifying party}. If the
provider delivers a different string $T'' \ne T'$ to the end user, verification
still succeeds on $T'$. Neither hardware nor a proof can prevent this; only the
end user holding $n$ and recomputing $b$ can.
\end{proposition}
\begin{proof}
$b$ is a function of $(n,T)$ only; the delivery channel to the end user carries
no authenticated term. One line.
\end{proof}

Two remarks scope this proposition. First, it describes a \emph{residual} gap,
not the claim that policy proofs bind nothing. Our own system already binds $T$
through the verifier-issued nonce $n$ of Definition \ref{def:system}, and a
verifier who receives $T'$ and recomputes $b$ does learn that $T'$ was the
proved string; the gap is that the binding is anchored to the
\emph{verifying party's} copy. Second, we do not claim to be the first to
observe that a proof can fail to bind the artifact it purports to describe. The
\emph{effort gap} for zero-knowledge inference proofs is formalized by
Hollow-LLM \cite{hollowllm}, and the \emph{source-integrity gap} for predicate
proofs is named by zk-attach \cite{zkattach}; both concern what the proof system
was asked to prove, whereas Proposition \ref{thm:gap} concerns where the proved
string is delivered. We state the naive version of our own claim because it
would be false against our own system.

\subsection{A quantitative binding--hiding tradeoff}
\label{sec:tradeoff}

The certificate must publish \emph{something} tying it to $T$; otherwise it
binds nothing at all. We therefore ask how much that public value leaks. The
answer depends on one measurable quantity --- the guessing entropy of the
response given everything the adversary already knows --- and on one structural
property of the public value itself.

\begin{definition}[Guessing entropy of a response]
\label{def:k}
For a response distribution $\mathcal{D}$ and adversary side information $Z$,
let
\[
\mathsf{Guess}(T \mid Z) \;\triangleq\; \max_{\mathcal{A}} \Pr[\mathcal{A}(Z) = T],
\]
which equals $\mathbb{E}_z[\max_t \Pr[T = t \mid Z = z]]$, and
\[
k \;\triangleq\; -\log_2 \mathsf{Guess}(T \mid Z),
\]
the maximum ranging over all time-unbounded adversaries.
\end{definition}

\begin{remark}[The definition is operational, and it is the average]
\label{rem:kmeas}
$k$ is \emph{measured}, not assumed: sampling one model on one prompt many
times estimates $\mathsf{Guess}(T \mid Z)$ directly
(\S\ref{sec:eval}), so Theorem \ref{thm:t1} is instantiated with a
number we obtain rather than a number we posit. The definition must average
over $z$ rather than take the pointwise worst case
$k_z = -\log \max_t \Pr[T = t \mid Z = z]$: the proof takes an expectation over
$z$, and $\mathbb{E}_z[2^{-k_z}] = 2^{-k}$ holds only for the average form.
\end{remark}

\begin{theorem}[Binding, and the hypothesis it needs]
\label{thm:t1}
Let $\beta : \mathcal{M} \to \{0,1\}^\lambda$ be public, deterministic and
efficiently computable, and let the certificate publish $\beta(T)$. Let
$\mathcal{A}$ be any adversary, unbounded in time, that is given $Z$ and
evaluates $\beta$ at most $q$ times.
\begin{enumerate}
\item[(H)] \emph{(Randomised public value.)} Suppose
  $\beta(T) = H(\mathsf{dom} \,\|\, \mathsf{u32be}(|n|) \,\|\, n \,\|\, T)$ for a
  random oracle $H$ and a nonce $n$ drawn uniformly and \emph{independently of
  $T$}. Then
  \[
  \Pr\!\left[\mathcal{A}^\beta\!\left(Z, \beta(T)\right) = T\right]
  \;\le\; q \cdot 2^{-k} \;+\; q \cdot 2^{-\lambda}.
  \]
  Consequently hiding with advantage $\varepsilon$ requires
  $k \ge \log q + \log(1/\varepsilon)$.
\item[(L)] \emph{(Lossy public value.)} For an arbitrary deterministic public
  $\beta$ the correct bound is over the \emph{preimage class}:
  \begin{multline*}
  \Pr\!\left[\mathcal{A}^\beta\!\left(Z, \beta(T)\right) = T\right]
  \;\le\; \mathsf{Guess}\!\left(T \mid Z, \beta(T)\right) \\
  \triangleq\; \mathbb{E}_y\!\left[\max_t \Pr[T = t \mid Z, \beta(T) = y]\right].
  \end{multline*}
\end{enumerate}
\end{theorem}

\begin{proof}[Proof sketch]
Part (L) is immediate after conditioning on $\beta(T)$; no hypothesis on
$\beta$ is used. For (H), fix $z$ and decompose the adversary's success by the
first query that hits $T$. The oracle answers it has already seen are uniform
and independent of $T$, so each query is itself independent of $T$ and succeeds
with probability at most $2^{-k_z}$. Taking the expectation over $z$ and
applying Definition \ref{def:k} yields $q \cdot 2^{-k}$; the term
$q \cdot 2^{-\lambda}$ charges a query that hits the uniform value $\beta(T)$
by chance. Appendix~\ref{app:proofs} gives the full argument.
\end{proof}

\fixme{$\lambda$ 比特随机与随机预言机只给 \emph{heuristic} 保证：实践中的 SHA-256
不是真的随机预言机（骨架 :712--715）。正文措辞必须写成 heuristic，
\textbf{不得}写成信息论界。}

\begin{remark}[The hypothesis in (H) is necessary, and our two instances straddle it]
\label{rem:betanecessary}
Without the hypothesis, part (H) is not merely unproven but \emph{false}: take
$\beta = \mathrm{id}$, or $\beta = \mathsf{encode}$ as in
Proposition \ref{prop:semantic}. The adversary outputs $\beta(T)$ and succeeds
with probability $1$, independently of $q$ and $k$. Our own two public values
therefore sit on opposite sides of the hypothesis---$b =
H(\mathsf{dom}\|n\|T)$ satisfies it, $\mathsf{encode}(T)$ does not---and this is
what makes the pair a theorem about \emph{our design} rather than a citation of
a standard bound. An earlier draft of this work stated (H) without its
hypothesis; the counterexample above is why the hypothesis is now explicit.
\end{remark}

\begin{corollary}[Re-execution, not enumeration]
\label{cor:reexec}
Under Assumption \ref{as:Z}, when the decoding is greedy or low-temperature,
$k \approx 0$ and the attack of Theorem \ref{thm:t1} is realised by
\emph{re-running the model}, at the cost of one inference---independent of
$|T|$ and requiring \emph{no} evaluation of $\beta$. Enumeration over the string
space is \emph{not} the relevant attack.
\end{corollary}

\fixme{原 docs/sp1-zk-audit.md §4 写的是"熵(~10^2--10^3 bit)远低于 256 bit，
因此枚举可行"——这是非因果推论（2^100 与 2^256 一样枚举不了；真正的攻击是\textbf{重执行}，
对手的已知信息是 (prompt, 模型, 解码参数)）。
2026-10-12 已全库改正，docs/tdsc-fixlist.md §G 的 G1--G6 逐条结案。}

\begin{remark}[A design observation, not a separation]
\label{rem:designated}
Our verification interface is \emph{designated}: binding is checked against a
value $b = H(\mathsf{dom}\|n\|T)$ in which a verifier-issued nonce $n$ appears.
It is tempting to promote this into a separation (``public verification is
impossible when $k$ is small''). \textbf{We claim no such theorem, and
Corollary \ref{cor:reexec} is why.} When $k \approx 0$, hiding fails for
\emph{every} scheme, designated or not, because the adversary re-runs the
model; when $k$ is large, a zero-knowledge proof of ``$\exists\, T : b =
H(\mathsf{dom}\|n\|T) \wedge J(\pi, T) = 1$'' is simultaneously publicly
verifiable and hiding. The dichotomy is thus governed by $k$ alone and carries
no independent content. What does survive---and what \S\ref{sec:discussion}
develops---is that \emph{third-party auditing of delivery} forces disclosure of
$T$, which is an informational remark rather than a theorem. We record the
demotion here because an earlier draft asserted the separation as a theorem.
\end{remark}

We stress that the designated-verifier shape is a \emph{retrospective
justification, not a novel construction}. The system publishes
$b = H(\mathsf{dom}\|n\|T)$ with a verifier-issued nonce $n$, and
Remark \ref{rem:designated} explains the shape of that choice after the fact.
The security model does not pin down whether $n$ is public: making it public
forfeits hiding, while holding it at the verifier is what makes verification
designated.

\begin{remark}[A hypothesis that can be settled by experiment]
\label{rem:kexp}
Theorem \ref{thm:t1} is only as useful as the number put into it. Because
Definition \ref{def:k} is operational, $k$ can be measured rather than assumed:
sampling one model on one prompt many times estimates
$\mathsf{Guess}(T \mid Z)$ directly, and the estimate substitutes into the
threshold $k \ge \log q + \log(1/\varepsilon)$. To our knowledge no prior work
measures this quantity for agent responses, which makes the measurement a
contribution in its own right; \S\ref{sec:eval} describes the protocol and
reports the result.
\end{remark}

\begin{proposition}[The lossy instance of Theorem \ref{thm:t1}(L)]
\label{prop:semantic}
For the semantic rule class, the public values contain $\beta(T) =
\mathsf{encode}(T)$, a content-preserving feature vector of the response. This
$\beta$ is deterministic and injective on the features but not on strings, so
it violates the hypothesis of Theorem \ref{thm:t1}(H) and instantiates part
(L): the binding force is exactly
$\mathsf{Guess}(T \mid Z, \mathsf{encode}(T))$, i.e.\ the guessing probability
\emph{within the preimage class} of the published vector. Here $k$ need not be
argued from perplexity at all: the adversary's search space is that preimage,
which is dramatically smaller than the string space. The relevant asymmetry is
not ``low entropy'' but a \emph{lossy public encoding}.
\end{proposition}
\begin{proof}
Immediate. The measurement is direct: publish $\mathsf{encode}(T)$ for real
responses and measure how much of $T$ is recoverable.
\end{proof}

Proposition \ref{prop:semantic} is the cheapest experiment in this paper and
the strongest anchor for Theorem \ref{thm:t1}: it is measurable in a day, and it
makes the bound concrete rather than cited. It also shows that the relevant
asymmetry is not entropy at all --- for the semantic rule class there is no
entropy argument available, because $\mathsf{encode}(T)$ simply \emph{is} a
content-preserving function of the response.

\subsection{Quantitative request binding}
\label{sec:t3}

\begin{theorem}[Selective-presentation advantage]
\label{thm:t3}
Fix a policy $\pi$ and a request distribution. A selective-presentation
adversary $\mathcal{B}$ may, \emph{before} seeing the receiver's challenge,
compute a set $S$ of accepted transcripts. The receiver then contributes $f$
fresh, uniformly random bits $r$ that enter the proven statement. Then
\[
\Pr\!\left[\mathcal{B}\ \text{is accepted on } r\right] \;\le\; |S| \cdot 2^{-f}.
\]
If instead the receiver contributes no fresh randomness ($f = 0$; e.g.\ an
empty nonce), the advantage is $1$, and \emph{does not decay in any parameter
of the system}.
\end{theorem}
\begin{proof}
The receiver's freshness $r$ enters the proven statement through
$b = H(\mathsf{dom} \,\|\, r \,\|\, T)$. For a certificate in $S$ to be accepted
on $r$, the value $r$ must fall inside the challenge set that certificate
already covers --- an event of probability at most $2^{-f}$ for $f$ fresh
uniform bits --- and a union bound over $S$ gives $|S| \cdot 2^{-f}$. When
$f = 0$ every element of $S$ is a permanently valid replay, so the advantage is
$1$ and does not decay in any parameter of the system.

The $f = 0$ branch is not hypothetical.
\texttt{circuits/\allowbreak types/\allowbreak src/\allowbreak lib.rs}
(lines 950--995) documents an empty nonce as a legal input that ``provides no
replay protection''; it is therefore a \emph{valid input to the current API},
not a straw man.
\end{proof}

Theorem \ref{thm:t3} is \emph{isomorphic} to Theorem \ref{thm:t1}, and that is
what earns it a place here. Both state the same fact about different inputs: a
proof binds only as much as the committed quantity the adversary cannot select
in advance has bits for. On the response side that quantity is the response and
the resource is entropy $k$; on the request side it is freshness $f$. Neither
can be bought with more computation. This is the spine announced in
\S\ref{sec:intro}, and the two theorems are its two sides.

\begin{remark}[A taxonomy, not a trichotomy]
\label{rem:taxonomy}
Every policy-proof deployment we surveyed closes the divergent-delivery gap of
Proposition \ref{thm:gap} in one of three ways:
(a) root trust in hardware (attested enclaves);
(b) have the \emph{end user} contribute a secret to the proven statement (the
user issues the nonce); or
(c) gate delivery on successful verification (\emph{prove-before-deliver}).
\textbf{We offer this as a classification, not as an exhaustiveness theorem}:
we neither prove that a fourth family cannot exist, nor that the three are
disjoint. zk-attach \cite{zkattach} walks (a); NiyamAI and ZKAP walk (c); PoP
walks (b).
\end{remark}

\paragraph{Posture, stated plainly.} Our system already carries
\texttt{response\_binding}, but $n$ is issued by the \emph{verifier}. Moving it
to the \emph{end user}, so that $f > 0$ becomes the receiver's contribution,
closes the divergent-delivery branch of Proposition \ref{thm:gap}. We are
explicit that this is a small protocol amendment rather than a new system: the
mechanism is already implemented, and only the issuer of the nonce changes. What
Theorem \ref{thm:t3} buys is the \emph{price}: with $f = 0$ there is no closure
at any cost.

% =====================================================================
\section{Policy-Private Compliance}
\label{sec:private-policy}

For a policy whose rules are deterministic and a response that is public
anyway, a signed hash chain over the response suffices: the signature already
binds the bytes, and anyone can re-evaluate the rules. The one scenario a
signature cannot cover is a policy that must itself stay secret---where the
operator has to certify compliance to a third party without handing over the
rules it is being judged against. That is the mode this section is about, and we
lead with what the mode \emph{cannot} buy, because that determines whether it
means anything at all.

\subsection{The observation that reframes this work package}
\label{sec:wp4-finding}

\texttt{spec\_canonical} is already a \emph{guest input}---a witness---rather
than a public value; only its hash,
\[
\mathsf{policy\_hash} \;=\; \mathsf{SHA256}(\texttt{spec\_canonical}),
\]
enters the public values (\texttt{circuits/\allowbreak types/\allowbreak
src/\allowbreak lib.rs}:1--16, 415--416, 918, 1003--1004). The instruction
``do not reveal the policy text'' is therefore
\emph{already true of the architecture}; what leaks is everything \emph{derived}
from the policy. Policy privacy is thus not a new proof system to be built but a
\emph{disclosure policy to be audited}, and we write it that way---as a
systematic account of what a certificate reveals about the policy it certifies.

\begin{table}[t]
\centering
\caption{What a policy-private certificate leaks today, and what closing each
leak costs.}
\label{tab:wp4}
\tabfont
\begin{tabular}{@{}L{0.04\linewidth}L{0.46\linewidth}L{0.36\linewidth}@{}}
\toprule
\# & Leak & Difficulty / dependency \\
\midrule
1 & The policy text itself & \textbf{Already closed} (it is a witness). \\
2 & Witness hiding in the proof system & Medium; \emph{inherited from the
    proving tier}: only \texttt{groth16}/\texttt{plonk} is witness-hiding, per
    \texttt{docs/\allowbreak sp1-zk-audit.md}. \\
3 & \texttt{PrivateViolation}\allowbreak\{\texttt{rule},\allowbreak\texttt{kind}\}
    names the rule that fired & Medium; needs a coarsened violation class plus a
    commitment to it. \\
4 & Evidence opening reveals the rule
    (\texttt{policydsl/\allowbreak privacy/\allowbreak commit.py}:319) & Medium;
    disclose the \emph{existence} of a violation, not its identity. \\
5 & The \texttt{delegated} flag exposes policy structure (which rules were
    routed to the \texttt{ezkl} path) & Medium; easiest to overlook, cheapest to
    fix. \\
6 & Semantic rules publish $\beta(T)=\mathsf{encode}(T)$
    (Proposition~\ref{prop:semantic}) & Medium--high; either bound the
    preimage, or disable this rule class in policy-private mode. \\
\bottomrule
\end{tabular}
\end{table}

The honest structure of Table~\ref{tab:wp4} is worth stating, because it sets
the reader's expectations for the rest of the section. Sub-problem~1 is a
one-paragraph re-reading of existing code, not a contribution. Sub-problem~2 is
inherited wholesale from the proving tier and is not specific to this work.
Sub-problems~3--5 are engineering with a single formal core---the granularity of
disclosure. Sub-problem~6 is where the work actually is, and it collides with
\S\ref{sec:tradeoff}: one cannot hide the policy while publishing the lossy
encoding of the response that the semantic rule class requires
(Proposition~\ref{prop:semantic}).

\subsection{The tension worth writing up: policy privacy vs.\ accountability}
\label{sec:wp4-tension}

\paragraph{The prerequisite, stated as a prerequisite.} Let
$H = \mathsf{SHA256}(\pi)$. On its own $H$ is \emph{inert}. A prover can always
emit a self-consistent certificate for any policy $\pi'$ it prefers, simply by
placing $H' = \mathsf{SHA256}(\pi')$ in that certificate. \emph{Secrecy of
$\pi$ is what creates the gap}, and no cryptographic mechanism can close it:
closing it would require showing $\pi$ to the verifier, which is precisely what
the mode exists to avoid. What closes it is an \emph{issuing act}, so we state
it as an explicit prerequisite of the mode.

\begin{description}
  \item[(P-PP)] An \emph{issuing authority} signs $\pi$ and registers $H$ in a
  record readable by third parties. A policy is \emph{accountable} iff its
  \texttt{policy\_hash} can be traced to such a record.
\end{description}

Without (P-PP) the mode is not merely weaker---the verifier's conclusion is
empty. This is the same assumption as the ``honest policy author'' of
\texttt{docs/security-model.md}~\S0.2; what is new is scoping it to the
verification interface, which had not been done, and which a referee would
otherwise use to dismiss the whole mode in one sentence.

\paragraph{Three parties.} ``Three'' means \emph{authority / operator / third
party}; an auditor is a named instance of the third party, not a fourth trust
anchor.

\begin{table}[t]
\centering
\caption{The three-party model of policy-private mode.}
\label{tab:wp4-parties}
\tabfont
\setlength{\tabcolsep}{3pt}
\begin{tabular}{@{}L{0.16\columnwidth}L{0.14\columnwidth}L{0.62\columnwidth}@{}}
\toprule
Party & Honest? & Holds; wants \\
\midrule
Authority $R$ & \textbf{Yes}, by (P-PP) & Holds $\pi$ and the issuing record;
wants $\pi$ actually enforced, and wants to be able to check that it was. \\
Operator $O$ & \textbf{No} (the adversary) & Holds $\pi$, $T$, the proving key;
wants $\pi$ hidden from third parties while still being able to certify. \\
Third party $V$ & Yes & Holds only public values---\texttt{policy\_\allowbreak hash},
\texttt{circuit\_\allowbreak hash}, \texttt{response\_\allowbreak commitment},
\texttt{passed}, the proof; wants to confirm that \emph{some issued} policy
passed, and must \emph{not} learn $\pi$. \\
\emph{(Auditor $V^{*}$)} & Yes & As $V$, plus $\pi$ once $O$ discloses it;
wants ex-post accountability. \\
\bottomrule
\end{tabular}
\end{table}

Two boundaries must be stated with the model. First, secrecy of $\pi$ comes from
a \emph{data-flow fact}---$\pi$ never enters the public values---and not from
zero-knowledge; any rule that writes a fragment of $\pi$ into the public values
(the leaks in Table~\ref{tab:wp4}, items~3 and~5) breaks secrecy on that
fragment. Second, accountability for $V^{*}$ is \emph{voluntary}: $V^{*}$ can
check only when $O$ chooses to disclose $\pi$. The mode offers no forced
disclosure; that would require the release-gating family discussed in
\texttt{docs/tdsc-plan.md}~\S1.6(b), which we do not claim.

\paragraph{What each verifier may conclude.} Table~\ref{tab:wp4-capability} is
the contribution of this subsection: the cryptography is downstream of it, and
the table survives even if the cryptography only partially lands.

\begin{table*}[t]
\centering
\caption{What each verifier may and may not conclude in policy-private mode.}
\label{tab:wp4-capability}
\tabfont
\setlength{\tabcolsep}{4pt}
\begin{tabular}{@{}L{0.13\textwidth}L{0.30\textwidth}L{0.30\textwidth}L{0.19\textwidth}@{}}
\toprule
Verifier & \textbf{May} conclude & \textbf{May not} conclude & Rests on \\
\midrule
Third party $V$ & (i) the commitment binds $T$
(Definition~\ref{def:system}); (ii) \emph{some} policy with hash $H$ passed;
(iii) $H$ appears in the issuing record & (iv) the \emph{content} of $\pi$;
(v) how many rules $\pi$ has, or of which classes; (vi) ``the operator enforces
the regulation \emph{I} care about''---only ``some issued regulation'' &
(P-PP); Assumption~\ref{as:Z}; the $k$ of Theorem~\ref{thm:t1} \\
End user (also receiver) & As $V$, and additionally: the proved $T$ is the
string it received, since its nonce is bound (Theorem~\ref{thm:t3}, requires
$f>0$) & Only the string it received; \emph{not} what was delivered to another
downstream agent & $f>0$ \\
Authority $R$ & Everything; it holds $\pi$ and can recompute $H$ & --- & --- \\
Auditor $V^{*}$ (after disclosure) & (vii) $\mathsf{SHA256}(\pi) =
\texttt{policy\_\allowbreak hash}$, i.e.\ what ran was the issued $\pi$;
(viii) the disclosed violation record is consistent & (ix) the undisclosed part
of $T$ & voluntary disclosure \\
On-chain verifier & (i)--(iii) only; it hashes and verifies signatures, and does
not interpret policy semantics & (iv)--(ix) & (P-PP) \\
\bottomrule
\end{tabular}
\end{table*}

The $V$ row is the reason this subsection exists. A third party does not obtain
``compliance''; it obtains ``some issued policy was satisfied''. Conflating the
two is the sentence a referee would use to reject the mode, and preventing that
sentence is what the table is for.

\paragraph{Prior art.} \emph{Policy hiding is not a new primitive}, and we claim
no primacy for it. Hidden credentials \cite{hiddencred}, secret handshakes
\cite{handshake}, privacy-preserving access control over XACML policies
\cite{xacmlpriv}, and attribute-hiding and fuzzy predicate encryption
\cite{sw05,ksw08} are one family, and policy-private mode is an instance of it.
What is new here is narrow: (a) the \emph{application} to LLM output compliance,
together with selective disclosure of violation type and location, and (b) the
accountability tension of this subsection.

\need{⚠️ 骨架 §4.3「Expected results, stated in advance」（子问题置信度百分比表）
\emph{未}迁入正文。那是 WP 规划的产物，不是论文内容——期刊稿里出现
「我们对子问题 4 有 70\% 把握」是硬伤。若作者要保留这份校准，应改写为
\S\ref{sec:discussion} 里一段关于适用范围与未完成项的讨论，而不是概率表。}

\need{(P-PP) 的落地形式尚未定：链上已签发策略哈希注册表，还是签名清单？
属作者决策。先例引文已按骨架 :623--630 的要求入 bibliography，
但 \cite{hiddencred} / \cite{handshake} / \cite{xacmlpriv} / \cite{sw05} /
\cite{ksw08} 五条\emph{均待逐条打开原文核实}（外部资源阻塞）。}

% =====================================================================
\section{Evaluation}
\label{sec:eval}

\subsection{Cycle count: length versus rule count}
\label{sec:eval-cycles}

Response length was pushed from 200 to 100{,}000 characters and the rule count
from 1 to 6, giving 40 measurement points (\texttt{bench/bench\_cycles.py}; the
full matrix is in \texttt{bench/results/cycles.md}). The response text comes from
an evaluation artefact the repository's own end-to-end demo produced (372
characters, three segments deduplicated by content); longer points are tiled from
it. What is measured is therefore the \emph{curve} of scanning cost, not ``a real
100{,}000-character document''. Every point was first confirmed to match no rule:
a matching rule short-circuits, and the cost of a short-circuit says nothing.

\begin{table}[t]
\centering
\caption{Cycle count against length, per rule set. Five rule sets, six length
points each; $R^2 \approx 1.0000$ throughout.}
\label{tab:cycles}
\tabfont
\begin{tabular}{@{}rrrrr@{}}
\toprule
Rules & \texttt{keyword} & \texttt{pattern} & Slope (cyc/char) & Intercept \\
\midrule
1 & 0 & 0 & 73.9      & 50{,}137 \\
2 & 1 & 0 & 105.3     & 65{,}942 \\
3 & 1 & 1 & 4{,}289.8 & 387{,}501 \\
4 & 2 & 1 & 4{,}048.8 & 405{,}785 \\
6 & 3 & 2 & 5{,}403.5 & 993{,}257 \\
\bottomrule
\end{tabular}
\end{table}

Three findings, in decreasing order of how much they change what may be claimed.

\emph{(1) With the rule set fixed, cycle count is exactly linear in length.}
Five rule sets, six length points each, $R^2 \approx 1.0000$ throughout
(Table~\ref{tab:cycles}).

\emph{(2) Scanning cost is dominated by the NFA, not by the rule count.} One
keyword rule costs about $+31.6$ cycles per character at the margin; one NFA
regular expression (an email pattern) costs about $+3{,}634$---roughly
$115\times$ more. Equivalently, a rule set containing a regular expression has a
per-character slope ($4{,}289.8$) $58\times$ that of a rule set with only string
rules ($73.9$).

\emph{(3) Unit prices are not additive: a unit price is a property of a rule
\emph{in context}, not of the rule.} Fitting the slopes as
$(\text{keyword},\text{pattern})$ unit prices gives $R^2 = 0.9169$ and a
\emph{negative} keyword unit price ($-477.2$ cycles per character). This is not
noise; two probes give the mechanism directly:

\begin{itemize}
  \item \textbf{Ordering probe} (same rules, same text, $L = 10^5$): declared
  order costs $541{,}342{,}069$ cycles against $570{,}736{,}989$ in reverse,
  a difference of $+5.43\%$, while \texttt{passed} is identical in both. \emph{The
  semantics of the decision do not depend on order; the cost does.}
  \item \textbf{Marginal probe} ($L = 10^5$, filling the one matrix cell the
  sweep could not reach---regular expressions with no keyword rule):
  $[\text{len},\text{pat}]$ over $[\text{len}]$ has a marginal unit price of
  $3{,}633.7$, whereas ``add one more keyword rule'' has two different answers:
  $+31.6$ alone, $+585.6$ in a policy that already contains a regular expression
  ($-153.7$ if inserted before it). The same action changes sign between
  $-154$ and $+586$.
\end{itemize}

The mechanism lies in zkVM memory allocation: every \texttt{ascii\_lower} and
every Pike VM match allocates, and allocator state depends on previous
allocations. We therefore quote \emph{per-rule-set slopes} (the table above,
directly citable) rather than any ``unit price $\times$ count'' expression. The
simple model $\text{cycles} \approx a\cdot|T|\cdot\text{rules} + b$ that this
work anticipated is refuted by the data ($R^2 = 0.872$, negative intercept), and
we report that too: \emph{a model refuted by its own data is itself an evaluation
result.}

\subsection{Matcher ablation: quadratic degradation on a pathological input}
\label{sec:eval-ablation}

\emph{Construction.} Pattern \texttt{a+b} against input
\texttt{a}$^{\times n}$ containing no \texttt{b}, so the match is necessarily
absent and both implementations must scan to the end. The naive implementation
reinvents and restarts at every start position, giving $\Theta(n^2)$; the Pike VM
decides in a single linear pass (\texttt{bench/bench\_ablation.py}).

\begin{table}[t]
\centering
\caption{Matcher ablation. Both paths decide identically; only the cost differs.}
\label{tab:ablation}
\tabfont
\begin{tabular}{@{}rrrrr@{}}
\toprule
$n$ & Pike (cycles) & Naive (cycles) & Ratio & Naive growth/char \\
\midrule
100 & 643{,}894   & 12{,}382{,}720  & 19.2$\times$ & --- \\
200 & 1{,}145{,}209 & 48{,}638{,}946 & 42.5$\times$ & $+96.4\%$ \\
400 & 2{,}143{,}617 & 193{,}207{,}130 & 90.1$\times$ & $+98.6\%$ \\
\bottomrule
\end{tabular}
\end{table}

Read per-character cost rather than absolute figures. Doubling $n$ leaves the
Pike VM's per-character cost essentially unchanged ($-11.1\%$, $-6.4\%$) while
the naive path's nearly doubles ($+96.4\%$, $+98.6\%$). At $n = 400$ the naive
path is already $90\times$ slower, and that ratio doubles again with each
doubling of $n$. \emph{This is an attack surface}: the input length for policy
compliance is chosen externally---an agent's reply, a concatenated session---so
an adversary who needs only to write a long reply drives a naive implementation's
cost quadratically. The comparison is meaningful only because the two paths
decide identically, which \texttt{tests/test\_ablation.py} pins down (Python and
Rust \texttt{match\_search} $\equiv$ \texttt{match\_search\_naive}). To our
knowledge this is the first ablation evidence specific to policy-compliance
matching.

\subsection{Real proofs: time, size, memory---and the ceiling}
\label{sec:eval-proofs}

Sampling points share the corpus of \S\ref{sec:eval-cycles}. Each point runs in
its own subprocess (SP1's memory accumulates within a process, so two proofs in
one process do not price one proof), and failures are reported in the table
rather than omitted. Rows marked OOM are a \emph{result}, not a malfunction: they
measure where this machine cannot go.

\begin{table}[t]
\centering
\caption{Real proofs on a 24-core / 12\,GB CPU machine. ``Config'' is
characters $\times$ rules; $k$ abbreviates $10^3$. Every succeeded
point is at the \texttt{core} (non-zero-knowledge) tier.}
\label{tab:proofs}
\tabfont
\begin{tabular}{@{}lrrrl@{}}
\toprule
Config & Time & Size & Peak RSS & Result \\
\midrule
$200 \times 1$  & 123.5\,s & 2716.4\,KiB & 10{,}389\,MB & \checkmark \\
$200 \times 2$  & 118.9\,s & 2716.9\,KiB & 10{,}438\,MB & \checkmark \\
$2k \times 1$   & 138.1\,s & 2718.6\,KiB & 10{,}449\,MB & \checkmark \\
$10k \times 1$  & 172.5\,s & 2726.6\,KiB & 10{,}506\,MB & \checkmark \\
$200 \times 3$  & ---      & ---          & ---          & OOM \\
$20k \times 1$  & ---      & ---          & ---          & OOM \\
\bottomrule
\end{tabular}
\end{table}

\begin{itemize}
  \item \textbf{A fixed floor of $\approx$10.15\,GiB.} The smallest
  configuration---200 characters, one rule---already peaks at 10{,}389\,MB: the
  prover's own overhead (the trace and permutation structures of a \texttt{core}
  proof) is almost independent of the trace. All four successes fall inside the
  narrow band 10{,}389--10{,}506\,MB. Length and rule count have little marginal
  effect on memory---\emph{but the headroom is too thin for that to help}.
  \item \textbf{A two-step staircase, not a slope.} One rule is provable up to
  10k characters; two rules at about 200 characters; \textbf{three or more do not
  come out at all}. The third rule in the sweep is a \texttt{pattern\_block}: a
  regular-expression match activates another family of AIR chips and lifts the
  trace area in one step, so the cliff is at $2 \to 3$ rules; in length it is at
  $10\text{k} \to 20\text{k}$.
  \item \textbf{Memory binds, not CPU.} The cycle sweep of
  \S\ref{sec:eval-cycles} reaches 100k characters with 6 rules because it only
  executes in the zkVM and produces no proof; at 100k characters the cycle count
  is still on the order of $10^7$, which any CPU handles. \emph{The proving-side
  ceiling and the execution-side ceiling are in different places.} ``How large a
  policy can be proved'' is, on this machine, a fixed floor plus two steps rather
  than a cost curve.
  \item \textbf{Artefact size is nearly independent of configuration}: from 200
  to 10{,}000 characters the proof grows only from 2716.4 to 2726.6\,KiB. Pure
  verification takes $\approx$90\,ms (the \texttt{vkey} setup is 1.5\,s, once).
  \item \textbf{Reproducibility, stated precisely.} The \emph{position} of the
  wall reproduces (each failing point failed twice); the \emph{timings} at the
  wall do not. Re-running one group in isolation shows $\pm 10\%$ jitter in time
  and $\pm 1\%$ in memory, so we quote time only to an order of magnitude.
  \item \textbf{Session and multiparty.} The session aggregation proof (three
  certificates) measured 152.5\,s; multiparty verification needs one proof
  \emph{per non-empty slice}, so its cost grows linearly in the number of slices
  (a two-proof example policy). Neither introduces a new proving system---a slice
  and a session are ordinary inputs to the guest---but both are subject to the
  same floor: every additional proof adds another $\approx$10\,GiB peak, so
  multiparty proving must run in separate processes, serially.
\end{itemize}

\need{⚠️ 本节的 groth16 档与 ZK/非 ZK 交叉点属 WP1，尚未打通。骨架 :665--668 的
矩阵规格为规则 $\{1,2,3,4,6,8\}$ $\times$ 长度 $\{200,2k,10k,50k\}$ $\times$
模式 $\{$transparent, private, private-policy$\}$ $\times$ 证明档 $\{$core, groth16$\}$
（\texttt{bench/bench\_scale.py}）。表 \ref{tab:proofs} 只覆盖 core 档。
groth16 列落地前不得填入任何 ZK 数字。}

\subsection{The non-ZK baseline: what a signature cannot do}
\label{sec:eval-baseline}

A signed Merkle log is the obvious alternative to a proof, and on cost it is not
close: everything is milliseconds. We measured it so that the comparison is not
rhetorical. Policy package \texttt{eu\_ai\_act\_v1.json}, 103-character response,
one challenge nonce, a 1024-leaf log; each timing is the median of $n = 500$
repetitions.

\begin{table}[t]
\centering
\caption{Signed-Merkle baseline (WP9). No proofs, no zkVM: this arm is pure
Python, which is why the figures are milliseconds.}
\label{tab:baseline}
\tabfont
\begin{tabular}{@{}lrl@{}}
\toprule
Metric & Median & Note \\
\midrule
Sign certificate (Ed25519 + DSSE) & 0.132\,ms & $n=500$ \\
Merkle append (leaf + root recompute) & 2.131\,ms & $n=500$ \\
Public verification (sig + inclusion) & 0.322\,ms & $n=500$ \\
Re-run verification ($+J(\pi,T)$) & 0.772\,ms & $n=500$ \\
Certificate size & 648\,B & \\
Inclusion proof (10 siblings) & 761\,B & \\
\bottomrule
\end{tabular}
\end{table}

\need{⚠️ 引用 2.131 ms 时必须写明它是 \texttt{proofs/session.py} 的\emph{静态树重算}
$\Theta(n)$，\emph{不是} RFC-6962 增量日志的摊销 $\Theta(\log n)$。
⚠️ 共享机器有 $\pm 20\%$ 漂移。数据源 \texttt{bench/results/baseline.md}。}

The asymmetry is the point, and it is what this arm contributes. The verification
chain is \emph{verify signature} $\to$ \emph{verify log inclusion} $\to$
\emph{re-run $J(\pi,T)$}. The first two steps need only public data; the third
needs the plaintext of both $\pi$ and $T$. That is not an implementation defect
but a structural consequence of the signature paradigm: a signature binds
\emph{who said it} to the bytes, never \emph{whether it is true}. We measured the
consequence rather than asserting it---\textbf{the same key signs a certificate
whose declared content is false, and the public path passes green}, because what
it sees is a certificate with a valid signature and a valid inclusion proof. Only
re-executing $J(\pi,T)$ recovers the truth, and taking that step away from the
verifier is exactly what the zero-knowledge design buys. The machine-checked
assertions are enumerated in \texttt{bench/results/baseline.md}; one row of that
file is worth repeating because it is the honest one: on
\emph{unforgeability} the baseline is \emph{not weaker} than the proof path, both
being Ed25519 and vkey binding, and we do not score a point for ZK there.

\subsection{End-to-end and auditing}
\label{sec:eval-e2e}

The end-to-end demo produces 14 certificates (streaming chain with early stop,
MCP argument and response checks, one zk proof) and anchors them; a third party
verifies with public artefacts alone, and the ledger chain, signature,
\texttt{policy\_hash}, anchoring, stream chains and zk proof all pass. In the
on-chain variant (a local Anvil node, \texttt{Anchor.sol} deployed, 14
certificate digests registered) the verifier reads back each certificate's
\texttt{anchoredAt} and cross-checks the local ledger, passing 14/14; the
negative control (an unregistered digest reads back as zero) confirms the check
is not vacuously true. Cross-validation submits 19 vectors to the Python
reference decision and to both the SP1 host and proving paths, comparing
\texttt{passed} and the violated rule set row by row; a full re-run with real
\texttt{core} proofs gave 19/19 on both paths.

\subsection{Measuring \texorpdfstring{$k$}{k}}
\label{sec:eval-k}

\need{对一组真实 prompt、同一模型多次采样，估 $\mathsf{Guess}(T\mid Z)$，
代入 Thm \ref{thm:t1}(H) 的阈值 $k \ge \log q + \log(1/\varepsilon)$（承接
Rem \ref{rem:kexp}）。据先例扫查无人测过这个量，本身是一块贡献；
同时它也是 Prop \ref{prop:semantic} 的测量（发布 $\mathsf{encode}(T)$ 后
$T$ 的可恢复比例）。\textbf{⚠️ 需真实模型多次采样，当前阻塞（外部资源）。}}

\need{⚠️ 「Why zero-knowledge?」决策图（骨架 :670--675）\emph{未}放入本节：
它讲的是设计结论（ZK 只在 $\{$私有响应, 私有策略$\}$ 两角有正收益），是结论不是测量，
按 docs/tdsc-writing-plan.md 的建议移至 §7 Discussion。与骨架现状不同，\textbf{待作者认可}。}

% =====================================================================
\section{Related Work}
\label{sec:related}

We position PoP against five lineages. For each, Table~\ref{tab:related} records
the slot it occupies---which is also, read the other way, what this paper does
not have to claim.

\begin{table}[t]
\centering
\caption{Positioning against the five prior-art lineages. The right-hand column
names what each lineage already occupies, and therefore what we do not claim.}
\label{tab:related}
\tabfont
\setlength{\tabcolsep}{3pt}
\begin{tabular}{@{}L{0.26\columnwidth}L{0.42\columnwidth}L{0.26\columnwidth}@{}}
\toprule
Lineage & Representative work & Slot already occupied \\
\midrule
Policy $\to$ circuit & zkGen \cite{zkgen}, Circuitree, Prezta & DSL-to-circuit
translation \\
Agent/LLM compliance & NiyamAI, Zaguk, PCO \cite{pco}, WITNESS & policy-level
private-mode kernel \\
Model computation & DeepProve, zkGPT, zkAgent \cite{zkagent} & inference
integrity \\
Attacks / negative results & Hollow-LLM \cite{hollowllm}, zk-attach
\cite{zkattach} & naming the gaps \\
Non-ZK industry baseline & \texttt{@brivora/verify}, VAOL & signed audit logs \\
\bottomrule
\end{tabular}
\end{table}

One of the five is the lineage this paper is actually in dialogue with. The
attacks-and-negative-results row names two gaps---Hollow-LLM's \emph{effort gap}
and zk-attach's \emph{source-integrity gap}---and both concern \emph{what the
proof system was asked to prove}. Proposition~\ref{thm:gap} concerns a third
gap on a different axis, \emph{where the proved string is delivered}, and we are
not aware of prior work that quantifies the binding/hiding tradeoff in the terms
of \S\ref{sec:binding}. That is why we state it as a bound with a stated
hypothesis rather than cite it.

\paragraph{zkAgent: orthogonal, not competing.} zkAgent \cite{zkagent} is the
closest contemporaneous system, and the comparison is easy to get wrong in both
directions. It proves that a provider executed the declared model and tools; PoP
proves that a response satisfies a public policy. PoP proves nothing about model
computation, and zkAgent proves nothing about policy compliance, so the two
compose: their conjunction covers ``inference was executed \emph{and} the policy
held''. Table~\ref{tab:zkagent} places the reported figures side by side.

\begin{table}[t]
\centering
\caption{PoP against zkAgent. The two prove different obligations, so the rows
are a statement of orthogonality, not a race.}
\label{tab:zkagent}
\tabfont
\setlength{\tabcolsep}{3pt}
\begin{tabular}{@{}L{0.22\columnwidth}L{0.38\columnwidth}L{0.30\columnwidth}@{}}
\toprule
& \textbf{zkAgent} \cite{zkagent} & \textbf{PoP} \\
\midrule
Obligation & provider executed the declared model and tools (inference integrity)
& response/trace satisfies policy $\pi$ (compliance) \\
Depends on the model & yes (weights, arithmetisation) & \textbf{no} \\
Cost grows with & model size, sequence length & response length $\times$ rule
count \\
Content privacy & none & committed, selective disclosure \\
Hardware reported & 32 cores / 512\,GB & 24 cores / 12\,GB \\
Proving & 194.4\,s (Shout, $T{=}512$) & 118.9--172.5\,s (200--10k chars) \\
Proof size & 96\,MiB (Shout); 3.1\,MiB (LogUp) & $\approx$2.7\,MiB \\
Verification & 0.42\,s (Shout); 0.038\,s (LogUp) & $\approx$0.090\,s \\
\bottomrule
\end{tabular}
\end{table}

We deliberately do not claim that PoP is faster, and the temptation to do so is
worth naming because the raw numbers invite it. The correct statement is
narrower, and we make only it: \emph{for the compliance obligation}, PoP
completes in a proving time of the same order as zkAgent's, with a smaller memory
footprint and a comparable proof size, without touching the model---and it adds
committed content privacy that zkAgent does not provide. Two caveats must travel
with those figures. The hardware is not matched (24 cores and 12\,GB against 32
cores and 512\,GB), and PoP's ceiling is memory-bound rather than cycle-bound, so
its feasible region is a two-step staircase rather than a curve.
\S\ref{sec:eval} reports the measurements and their limits.

\need{⚠️ tab:zkagent 里 PoP 那一列全部取自 SP1 \texttt{core}（非零知识）档，
机器为 24 核 / 12\,GB WSL2；zkAgent 一列为 ePrint 2026/199 的报告值（32 核 / 512\,GB）。
数据源：\texttt{bench/comparison\_zkagent.md}。该文件自身警告：
\emph{不得}宣称「PoP 比 zkAgent 快」——正确口径是「对『合规』这一义务…
同量级证明时间 / 更小内存」。groth16 档（WP1）未打通前，本表不得出现 ZK 档数字。}

The one influence we should name plainly, rather than leave to resemblance:
zk-regex's ``match offline, verify the path in circuit'' pattern is the shape of
our NFA contract, and verifiable document redaction \cite{vdr} inspired the bit
selector behind \texttt{mask\_covered}. The borrowing stops at the circuit
technique and does not carry the security argument across: VDR assumes a
\emph{signed original document} exists, whereas in PoP's private mode both the
original and the redacted version are produced by the operator, so there is no
trusted original to appeal to. We use the mechanism, not the proof.

\fixme{引文纪律：除 Hollow-LLM \cite{hollowllm} / zkGen \cite{zkgen} /
zkAgent \cite{zkagent} / PCO \cite{pco} 四条外，其余引文必须逐条打开原文
核实后再进 bibliography（docs/tdsc-fixlist.md \S2、WP8）。
\textbf{骨架原先把「ePrint 2024/813」列进已核实四条，那是错的}，已剔除：
\textbf{绝不能再引 IACR ePrint 2024/813 作为 VDR}——那是
Kiraz/Larraia/Vaughan, ``How to Redact the Bitcoin Backbone Protocol''，
与 Verifiable Document Redacting 无关。见 fixlist \S A。
⚠️ 口径：WITNESS 参考实现为 Apache-2.0（并非闭源），故「开源」不构成差异化。}

% =====================================================================
\section{Discussion and Limitations}
\label{sec:discussion}

Everything above was kept deliberately narrow, and this section says how narrow.
It covers the four places where a reader is most likely to round our claims up:
the caliber of the privacy statement (\S\ref{sec:disc-privacy}), where
zero-knowledge actually pays (\S\ref{sec:disc-zkmap}), availability against an
adversarial response (\S\ref{sec:disc-dos}), coverage
(\S\ref{sec:disc-coverage}), and composition with an inference proof
(\S\ref{sec:disc-composition}). It closes with the engineering limitations that
remain (\S\ref{sec:disc-rest}).

\subsection{The caliber of the privacy claims}
\label{sec:disc-privacy}

\textbf{``Content privacy'' means \emph{no plaintext of $T$ in the public
values}, not \emph{content is unrecoverable}.} Two independent reasons force
this reading---one at the proof-system level, one at the design level---and the
weaker reading is the one our own measurements support.

\emph{Proof-system level.} The default proving tier is SP1's \texttt{core}, and
SP1's published security model states explicitly that individual STARK proofs do
not currently satisfy the zero-knowledge property; the source agrees, with no
blinding anywhere in the proof stack and the shard proof carrying the trace
Merkle root and trace openings in the clear.%
\need{⚠️ 引文债：「SP1 公开安全模型承认其 STARK 证明当前不满足零知识」是
\emph{外部}断言，出处（SP1 官方安全模型文档的具体章节/URL）尚未核实到可引的
一级来源，故此处只陈述、不给 \cite{}（WP8）。若核不到原文，须撤回该句，
只保留我们自己的源码审计结论那半句。}
Soundness (G1) therefore holds
while the proof artefact itself does not hide the witness. The
\texttt{gr\-oth16} and \texttt{plonk} tiers wrap the inner STARK proof as a
private witness---the only mode in which hiding is even possible---but that is a
wrapper-level claim, it is not post-quantum, it needs $\ge 16$\,GB, and we did
not measure it (\S\ref{sec:eval-proofs}).

\emph{Design level, and the more fundamental one.} The public values
\texttt{response\_binding} and \texttt{response\_commitment} are both
\emph{public, offline-recomputable functions of $T$}. This is not an oversight:
\texttt{response\_binding} \emph{must} be publicly recomputable, because that is
how a verifier independently checks that the delivered $T'$ is the $T$ that was
proved---the swap-response defence. But the same property is a guess-and-verify
oracle for $T$.

It is worth being precise about the attack, because the obvious version of the
argument is wrong. ``A one-line response carries only $10^2$--$10^3$ bits of
entropy, far below 256, so enumeration is feasible'' does \emph{not} follow:
$2^{100}$ is as unenumerable as $2^{256}$. The real attack is
\emph{re-execution} (Corollary \ref{cor:reexec}). The adversary's side
information is not the response space but the tuple (prompt, model, decoding
parameters) (Assumption \ref{as:Z}); under greedy or low-temperature decoding
the response is nearly determined, so $k = -\log_2 \mathsf{Guess}(T \mid Z)
\approx 0$ (Definition \ref{def:k}) and the attack costs \emph{one inference},
independent of $|T|$. Under the two premises that a verifier independently
recomputes the binding and that $k < \log q + \log(1/\varepsilon)$, response
binding and response-content hiding are mutually exclusive.

We therefore restate (G2) as what is actually established: \emph{in private mode
the public values contain no plaintext of $T$}---which the leak experiment
confirms directly---and not as content irrecoverability. The consequence for the
comparison in \S\ref{sec:related} is that the ``content privacy'' cells on
\emph{both} sides must not be read as ``content is unrecoverable.'' This is a
bound on our own system first.

\subsection{Where zero-knowledge pays, and where it does not}
\label{sec:disc-zkmap}

Section \ref{sec:eval-baseline} measured the non-ZK alternative---a signed
Merkle log---and everything in it is milliseconds. A proof therefore has to earn
its cost against what a signature cannot do, and our measurements locate that
boundary instead of asserting it.

The map has two axes: is the response secret, and is the policy secret?
Zero-knowledge has positive value only in the corners where at least one of them
is unknown to the verifier. In the transparent corner---public response, public
policy---a signed log certifies exactly the same thing at three to four orders
of magnitude less cost, and adding a proof is redundancy rather than rigour. The
familiar criticism that ``a signature would have done'' is correct in that one
corner and wrong in the other three; the decision map is what turns the
criticism into a measured design conclusion.

One dependency belongs in the open. The zero-knowledge property of a certificate
is a property of the \emph{wrapper}, not of the system as a whole, and the
default tier does not have it (\S\ref{sec:disc-privacy}). If \texttt{gr\-oth16}
does not become a delivery tier, we do not claim the artefact is
zero-knowledge, and the word must leave the title.

\subsection{Availability: the unanchored-pattern cost attack}
\label{sec:disc-dos}

A regular expression with an unanchored pattern---one with no left anchor, such
as \texttt{.*foo}---gives the matcher no way to reject a prefix of the response
early: every position stays a live start position until the whole string has
been read. The prover's work therefore scales with the response, and a party who
\emph{chooses} the response can drive it quadratically. In the agent setting
that party is often the agent itself, or whoever controls its output.

Two things are worth separating. This is a denial of service against the
\emph{prover}, not an attack on soundness: no false certificate is produced, and
the cost is simply borne by the honest side. And it is specific to proving
rather than to matching, because the trace---not merely the verdict---is what
gets proved. Mitigations exist and are partly in place: a Pike VM (already
implemented) simulates the NFA in linear time, committing to the response length
up front removes the adversary's freedom to inflate it, and chunked matching
bounds the per-chunk work. Prior work on zero-knowledge regular expressions
imposes the DFA restriction but does not analyse denial of service in a zkVM
setting, which is where this differs.

\need{⚠️ 作者决策未定：骨架 :649--651 把这条写成一条定理（$\Omega(n^2)$ 证明
工作量），但证明没写。需定：写成定理（并补证明），还是降为观察（Observation）。
本稿先按观察陈述、不给定量断言，以免留下一条无证明的定理。}

\subsection{Coverage: a provider can simply withhold certificates}
\label{sec:disc-coverage}

Nothing in the cryptographic layer can compel a provider to issue a certificate
at all. A provider that never produces one is indistinguishable, from the
verifier's side, from a provider with nothing to certify. We treat this as an
open problem rather than a defect of the construction: without trusted hardware
there is no anchor that forces issuance, and the mechanisms that could---%
sampling audits, economic penalties, disclosure obligations---live above the
cryptographic layer and outside the scope of this paper. Naming it as an open
problem is more useful than proposing a mitigation we cannot evaluate.

\subsection{Why composition alone does not close the gap}
\label{sec:disc-composition}

The natural next step is to compose: prove that the inference produced $T$ from
the request, conjoin that with a proof that $\pi(T)$ holds, and conclude that
the agent \emph{sent} a compliant response. We describe the mechanism, and we
record why the conclusion does not follow for free.

\emph{Composition by itself does not close the gap.} A model proof bound to
\emph{some} input shows only that \emph{some} input yields a compliant output;
the divergent-delivery gap of Proposition \ref{thm:gap} closes only when the
model proof is bound to the \emph{input the user submitted}, $x$, so that
$M(x) = T \wedge \pi(T)$ becomes a statement about the bytes actually sent. That
requirement is not an extra condition bolted on: it is exactly $f > 0$ in
Theorem \ref{thm:t3}, with the user-submitted $x$ playing the role of the $f$
fresh bits. The two halves of this paper meet here---the request-side theorem
names the condition the composition would have to satisfy, and the response-side
theorem is what makes that condition non-vacuous.

The composition contract is correspondingly short to state: the inference proof
must expose $T$ in its public values, and it must be bound to $x$.

\need{⚠️ 阻塞：真模型规模（$\ge 10^6$ 参数）与组合成本数据属 WP4 子问题② /
WP3 的 DeepProve 对照臂，需外部资源。代理规模下「组合成本由推理证明主导」
\emph{未被验证}（\texttt{bench/results/compose.md}），不得写入。
骨架 :641--642 的组合契约定义可在此处补。}

\subsection{Remaining limitations}
\label{sec:disc-rest}

\begin{itemize}
  \item \textbf{Three-party collusion is undetectable at the signature layer.}
    If the policy author, the operator, and the prover all collude, the
    certificate is indistinguishable from an honest one. The multi-prover
    construction partitions responsibility; it does not prevent agreement
    (\texttt{docs/security-model.md}).
  \item \textbf{Signing and key custody are not production-grade.} The on-chain
    path currently takes plaintext private-key parameters, the demo uses a
    public Anvil test key, and deployment is a local Anvil instance. Production
    signing (Ed25519/HSM) and key custody remain future work.
  \item \textbf{Semantic-rule model quality.} The delegated semantic rule works
    end to end, but the bundled model is a small demonstration proxy rather than
    a usable learned guardrail: ``$P(\text{harmful})$ is low'' is a statement
    about \emph{this} model. The companion proof's public instance also contains
    $\mathsf{encode}(T)$, so semantic rules are \emph{public-mode only}.
  \item \textbf{The inference proof is faithful, not good.} The inference half
    is a deterministic fixed-point MLP proxy whose weights are generated from a
    compile-time seed and baked into the program, hence committed by the
    verification key. It proves that \emph{this} graph computed \emph{this}
    output on \emph{this} response---not that the graph is any good. Swapping the
    proxy for a real prover touches only the prover; the composition layer does
    not move.
  \item \textbf{Verification remains $O(N)$.} There is no recursive aggregation,
    so a verifier pays per certificate. The supported regex subset and its ASCII
    semantics are bounded in the same way, and proof-cost optimisation (lookup
    arguments, parallelism, precomputation) is unexplored.
  \item \textbf{Slices share one guest.} Multiple provers hold distinct keys,
    but that is a division of responsibility rather than a per-role circuit: the
    repository does not build one circuit per role.
\end{itemize}

% =====================================================================
\section{Conclusion}
\label{sec:conclusion}

This paper began from a small question---what does a compliance certificate
actually bind?---and the answer came out the same on both sides of the
interaction. A proof binds only the bits that were submitted into it \emph{and}
that the adversary could not select in advance. On the response side those bits
are the entropy $k$ of $T$ given everything the adversary already knows, and no
amount of computation buys more of them; on the request side they are the $f$
bits of freshness the receiver contributes, and a certificate without them is
replayable without limit. Theorem \ref{thm:t1} and Theorem \ref{thm:t3} are two
faces of that one statement, and Proposition \ref{thm:gap}---a gap we state
against our own design---is what makes it bite.

PoP is the constructive instance of that analysis, not a claim to have invented
a new proof system. It composes an existing zkVM, an existing policy compiler,
and a signature layer into a certificate that a verifier holding only public
values can check without trusting hardware, and it is explicit about what such a
certificate does not do: it does not make the response unrecoverable, it does not
hide the witness at the default proving tier, and it cannot compel a reluctant
provider to issue anything at all. Those are limits of the problem, stated as
limits, rather than of one implementation.

The design conclusions are accordingly modest, and we hope reusable: publish the
leak profile of a certificate and audit it like a disclosure policy; let a
privacy claim be only as strong as its arithmetic supports; and put the
freshness where the binding is. A certificate is worth exactly as much as the
bits that went into it.

% =====================================================================
% ⚠️ IEEEtran 的 \appendix 是「单个附录」命令：它把 \section 重定义成
% \@IEEEdestroythesectionargument，直接丢弃 \section{...} 的参数（只打 warning）。
% 多个附录必须用 \appendices（复数），它把 section 计数器置 -1，
% 之后每个 \section{X} 才依次成为 A、B、……  （2026-10-12 实测：
% 用 \appendix 时 A/B 两个附录的 \ref 都解析成 "A"，且 title 字段退化成 \appendixname，
% 锚点退化成 section*.1；「0 undefined ref」检查抓不到这种错。）
\appendices
\section{Proofs}
\label{app:proofs}

\begin{proof}[Proof of Theorem \ref{thm:t1}]
\emph{(H).} Fix $z$ and let the adversary's query set be
$Q = \{q_1, \dots, q_q\}$. Decompose by first hit,
\[
\Pr[\exists i : q_i = T] \;=\; \sum_i \Pr\!\left[q_i = T \wedge \neg E_{<i}\right],
\]
where $E_{<i} \triangleq \{\exists j < i : q_j = T\}$.
Conditioned on $\neg E_{<i}$, every oracle answer $a_1,\dots,a_{i-1}$ already
revealed is uniform and independent of $T$ (no queried point equals $T$, and $H$
is a fresh random function), and $\beta(T)$ itself is uniform and independent of
$T$. Hence $q_i$---a deterministic function of $(z, \beta(T), a_{<i})$---is
\emph{independent of $T$}, so
\begin{multline*}
\Pr[q_i = T \mid Z = z]
= \sum_t \Pr[q_i = t]\,\Pr[T = t \mid Z = z] \\
\le \max_t \Pr[T = t \mid Z = z] = 2^{-k_z},
\end{multline*}
giving $\Pr[\exists i : q_i = T \mid z] \le q\,2^{-k_z}$. Taking the expectation
over $z$ and using $\mathbb{E}_z[2^{-k_z}] = \mathsf{Guess}(T \mid Z) = 2^{-k}$
from Definition \ref{def:k} yields $q \cdot 2^{-k}$. An adversary that outputs
an unqueried value wins with probability at most $2^{-k_z} \le q\,2^{-k_z}$,
which the same bound absorbs. The term $q \cdot 2^{-\lambda}$ charges the
probability that a query hits the uniform value $\beta(T)$ by chance.

\emph{(L).} Immediate from Definition \ref{def:k} after conditioning on
$\beta(T)$; no hypothesis on $\beta$ is used.
\end{proof}

\need{Thm \ref{thm:t1}(L) 无闭式界，只给测量口径（\S\ref{sec:eval}）。
三分法已降级为分类 Rem \ref{rem:taxonomy}，无需证明。
原 T2「分离定理」已\emph{停用}（见 Rem \ref{rem:designated}）—— 不要为它写附录。}

\need{\textbf{待补}：把 Thm \ref{thm:t1}(H) 的随机预言机模型换成标准模型，
或明确声明本文只在 ROM 下主张。当前证明依赖 $H$ 是新鲜随机函数这一条；
实践中的 SHA-256 只给计算意义下的近似，故正文措辞必须写成
\emph{heuristic}，不得写成信息论界。}

\section{Reproduction}
\label{app:repro}

The full walkthrough---environment, a single compliance proof, cross-checking
against the reference layer, private mode, certificate issuance, the one-command
demo, and third-party verification---is maintained in \texttt{docs/reproduce.md}
together with the exact commands and their expected output. This appendix records
only what a reader needs in order to judge the claims of \S\ref{sec:eval}.

\paragraph{Environment.} Ubuntu or WSL2 on x86-64, no GPU. The Rust/SP1 side
needs Rust stable, \texttt{cargo-prove}~v6.7.0, Go $\ge$\,1.24, and the usual
native toolchain; the reference layer is pure Python~3.10+ and runs on its own.
Memory, not time, is the binding constraint: the CPU prover carries a fixed
floor of approximately 10.15\,GiB regardless of the trace
(\S\ref{sec:eval-proofs}), and 12\,GB suffices only for small policies.

\paragraph{What reproduces without the prover.} The reference layer, the policy
compiler, and the whole test suite run in seconds and need no Rust. This is
enough to reproduce the non-ZK baseline of \S\ref{sec:eval-baseline}, to inspect
the certificate and inclusion proofs, and to check the cross-layer contract---%
that the Python reference implementation and the zkVM guest consume the same
compiled policy bytes.

\paragraph{What needs the prover.} Every timing, size, and memory figure in
\S\ref{sec:eval-cycles}--\S\ref{sec:eval-proofs}, and the end-to-end result of
\S\ref{sec:eval-e2e}, requires a real proof. The measurements were taken on a
24-core CPU machine with 12\,GB of RAM, which is what places the ceiling reported
there: the successes and the two out-of-memory configurations in
Table~\ref{tab:proofs} are both properties of that machine.

\paragraph{Where each number comes from.} Every measurement in \S\ref{sec:eval}
is the content of a checked-in result file with a machine-readable companion:
\texttt{bench/results/} holds \texttt{cycles}, \texttt{ablation},
\texttt{proofs}, \texttt{baseline}, \texttt{verify}, \texttt{semantic},
\texttt{streaming}, \texttt{compose}, and \texttt{prover\_knobs}, each as a
\texttt{.json} alongside its \texttt{.md}. The reproduction guide records the
commit, the hardware, and the prover binary digest for each run, so that a
reported figure can be traced to one specific execution rather than to a state
of the tree.

\need{⚠️ 数字漂移（本轮实测发现，需作者定夺）：正文 §2.2 写「801 tests / 15 skipped」，
其依据是实测（\texttt{docs/dev-plan.md:506} 的落地记录）\emph{以及}本轮对
\texttt{tests/} 的静态清点（\texttt{def test\_} 计 801 个）——\textbf{801 是对的}。
但 \texttt{docs/reproduce.md:31,428}、\texttt{docs/security-model.md:829}、
\texttt{docs/development.md:93,290}、\texttt{docs/modules/08-tests-bench.md:11,58,504,526}、
\texttt{docs/modules/README.md:94,161,201} 仍写 \textbf{783}（陈旧值，
\texttt{docs/report/advisor-report.md:586} 已把这条「文档数字漂移」登记在案）。
本附录既然指向 \texttt{reproduce.md}，照抄就会让复现者跑出 801 却以为环境坏了——
\textbf{这是真实的复现缺陷}。按仓库自己的「过期就标注、不静默改」纪律，
本稿不擅自改那些文档，只在此记债。作者需定：把五处文档回填为 801
（\texttt{advisor-report.md:599} 的建议），还是在附录里并列两套口径。}

% =====================================================================
\begin{thebibliography}{99}

\bibitem{hollowllm}
C.~Gong, B.~Liu, and M.~Li, ``Hollow-LLM Attack: Computationally Trivial
Weights in Zero-Knowledge Verification of LLM Inference,'' in \emph{IEEE S\&P},
2026, pp.~4339--4355, DOI 10.1109/SP63933.2026.00258.

\bibitem{zkattach}
\need{zk-attach, arXiv:2608.30083 --- 打开原文核实作者后再引}

\bibitem{zkgen}
J.~Lauinger, J.~Ernstberger, and S.~Steinhorst, ``zkGen: Policy-to-Circuit
Transpiler,'' in \emph{IEEE ICBC}, 2024, pp.~620--624,
DOI 10.1109/ICBC59979.2024.10634440.

\bibitem{pco}
\need{PCO, IACR ePrint 2026/994 --- 打开原文核实作者（Dolev/Yung 等）后再引}

\bibitem{zkagent}
L.~Wang, H.~Lou, C.~Li, Y.~Yu, and Y.~Hu, ``zkAgent: Verifiable LLM Agent
Execution via One-Shot Transcript Proofs,'' IACR ePrint 2026/199.

\bibitem{hiddencred}
D.~Balfanz and J.~Schweitzer, ``Hidden credentials,''
\need{核实出处：2002 年 unpublished manuscript / 2003 年会议版本，二者择一并核准页码}.

\bibitem{handshake}
\need{核实作者列表与页码（首发 S&P 2003, pp.~139--150）}D.~Balfanz, G.~Durfee,
N.~Shankar, D.~K.~Smetters, J.~Staddon, and H.-C.~Wong, ``Secret handshakes from
pairing-based key agreements,'' in \emph{IEEE S\&P}, 2003.

\bibitem{xacmlpriv}
\need{核实作者与出处。骨架点名的方向是「XACML 策略的隐私保护访问控制」，
候选：C.~A.~Ardagna 等关于 XACML 隐私的系列工作。\textbf{未核实前不得写成
具体条目}。}

\bibitem{sw05}
\need{核实页码（LNCS 3494, pp.~457--473）}A.~Sahai and B.~Waters, ``Fuzzy
identity-based encryption,'' in \emph{EUROCRYPT}, 2005.

\bibitem{ksw08}
\need{核实页码（LNCS 4965, pp.~515--534）}J.~Katz, A.~Sahai, and B.~Waters,
``Predicate encryption supporting disjunctions, polynomial equations, and inner
products,'' in \emph{EUROCRYPT}, 2008.

\bibitem{vdr}
\need{Verifiable Document Redacting, Springer, DOI 10.1007/978-3-319-66402-6_20
--- 核实作者与年份。\textbf{绝不能再引 IACR ePrint 2024/813}（那是
Kiraz/Larraia/Vaughan, "How to Redact the Bitcoin Backbone Protocol"）。
见 docs/tdsc-fixlist.md \S A。}

\end{thebibliography}

\end{document}
"   % 两遍，第二遍收交叉引用
%  \SKELQUIET 让所有中文标记宏吞掉参数；不带它则需要能排版中文的引擎
%  （xelatex + CJK 字体）。正文本身全英文，标记里才有中文。
%
%  ⚠️ IEEEtran.cls 需要装在 TeX 发行版里。WSL 侧已装进
%     /home/dong/texlive/texmf-local/tex/latex/ieeetran/（2026-10-12）。
% =====================================================================
\documentclass[journal]{IEEEtran}

\usepackage{amsmath,amssymb,amsthm}
\usepackage{booktabs}
\usepackage{array}
\usepackage{graphicx}
\usepackage{xcolor}
\usepackage{hyperref}

% 正文用 Times（ptm，窄），\texttt 用 Courier（pcr，等宽且宽）。
% 一行里只要有一个长 Courier 标识符（\texttt{normalized\_keyword\_block}
% 这类）就几乎没有可拉伸的胶，TeX 只能拉出 badness 10000 的大空隙。
% emergencystretch 允许第三遍在这些行上额外拉伸，把最坏情况吸收掉；
% 代价是那些行略松，远好于一道 10000 的豁口。
% （2026-10-12：装上真正的 Times 后此类行从 3 涨到 12 才暴露。）
\setlength{\emergencystretch}{2em}

% microtype：字形微伸缩（pdftex 下 ±2%）+ 标点悬挂。
% emergencystretch 只能吸收「可拉伸胶不足」，但 badness 在 10000 处封顶，
% 段首被后面长片段（consequences---customer-facing、normalized\_keyword\_block）
% 挤空的行照样报 10000。microtype 从字形宽度上补，是这一类的标准解法。
\usepackage{microtype}

% 窄栏里用 p{} 会因两端对齐产生大片字距（编译时刷一堆 Underfull \hbox）。
% 表格正文列一律用 L{}：左对齐、允许自然参差。
\newcolumntype{L}[1]{>{\raggedright\arraybackslash}p{#1}}
\newcommand{\tabfont}{\footnotesize}

\newtheorem{theorem}{Theorem}
\newtheorem{lemma}[theorem]{Lemma}
\newtheorem{proposition}[theorem]{Proposition}
\newtheorem{corollary}[theorem]{Corollary}
\newtheorem{definition}[theorem]{Definition}
\newtheorem{assumption}[theorem]{Assumption}
\newtheorem{remark}[theorem]{Remark}

% --- 工作标记：提交前全部清掉 ---
% 注意：宏名不能叫 \wp —— \wp 已是 LaTeX/amssymb 的 Weierstrass ℘，
% \newcommand{\wp} 会直接报 "Command \wp already defined" 并使后文级联报错。
\ifdefined\SKELQUIET
  \newcommand{\wpmark}[1]{}
  \newcommand{\need}[1]{}
  \newcommand{\fixme}[1]{}
\else
  \newcommand{\wpmark}[1]{\textcolor{red}{\textsuperscript{[WP#1]}}}
  \newcommand{\need}[1]{\textcolor{blue}{[[ NEED: #1 ]]}}
  \newcommand{\fixme}[1]{\textcolor{orange}{[[ FIX: #1 ]]}}
\fi

\begin{document}

% ---------------------------------------------------------------------
% 标题：决策 D5 —— 保留 "Zero-Knowledge" 的前提是 groth16 成为交付档 (WP1)。
%       若 WP1 未打通 groth16，改用下面被注释掉的那个标题，
%       并同步删掉摘要与 §8 里的 zero-knowledge 措辞。
% ---------------------------------------------------------------------
\title{Proof-of-Policy: Zero-Knowledge Compliance Certificates\\ for LLM Agent Outputs}
% 备选标题（若 groth16 未打通）：
% \title{Proof-of-Policy: Succinct Compliance Certificates for LLM Agent Outputs}

\author{\need{authors}}
\maketitle

% ---------------------------------------------------------------------
\begin{abstract}
% 决策 E2：内容隐私上界必须出现在摘要，不能埋进 §Limitations。
% 决策 D1：贡献陈述按"分析 x 系统各半"写 —— 分析在前，系统是分析所需的那一实例。
Proof-of-Policy (PoP) lets an auditor check that an LLM agent's response satisfies
a published policy, without trusted hardware and without the operator disclosing
the response in the clear. Our contribution is analytical before it is
constructive. We give a quantitative bound on what such a certificate binds---a
guessing-entropy term that no additional computation can improve---together with
the structural assumption the bound requires, and a proof that the assumption is
necessary rather than convenient. The same accounting governs the request side,
where a certificate without receiver-supplied freshness is replayable without
limit. PoP is then presented as the instance these bounds require---a policy
compiler, a
% 决策 D5 条件措辞：若 groth16 未成为交付档，把下行的 "zero-knowledge" 换成 "succinct"，
% 并同步改标题与 §8。
zero-knowledge circuit, and a signed certificate layer---measured end to end
rather than asserted, with the leak profile of the certificate reported as a
result in its own right.

\smallskip
\noindent\textbf{Scope of privacy.} PoP provides \emph{policy}-zero-knowledge:
the certificate reveals neither the policy evaluation trace nor, in
policy-private mode, the policy itself. It does \emph{not} guarantee that $T$
is unrecoverable: as we show in \S\ref{sec:binding}, any publicly recomputable
binding to $T$ admits a re-execution attack whenever the adversary knows the
prompt and the model. The same bit-budget governs the request side: a
certificate is replayable without limit unless the receiver contributes fresh
randomness (Theorem \ref{thm:t3}).
\end{abstract}

\begin{IEEEkeywords}
zero-knowledge proofs, LLM agents, policy compliance, verifiable computation,
accountability
\end{IEEEkeywords}

% =====================================================================
\section{Introduction}
\label{sec:intro}

LLM agents increasingly take actions that carry con\-se\-quences---\allowbreak{}
customer-facing replies, tool invocations, financial and medical content. Regulation such as the
EU AI Act (Articles~12 and~13) and enterprise governance both require that
compliance be auditable. In practice it can only be taken on trust: logs can be
forged, guardrail claims cannot be checked, and content privacy forbids retaining
the plaintext that would settle the question.

Four lines of work each occupy one corner of this problem without spanning it.
\emph{TEE execution proofs}---Proof-of-Guardrail and its relatives---show that
guardrail code was executed and that inputs and outputs were signed; as that work
itself observes, \emph{execution is not compliance}, and the approach carries
roughly $34\%$ latency overhead together with a hardware dependency.
\emph{TEE governance stacks} (Agentrust, cMCP, cA2A) inherit both the hardware
and the vendor facility. \emph{Formal verification} (Lean-Agent and similar)
requires the entire policy to be formalised by hand, and does not apply to
black-box responses. \emph{ZK regex and PII primitives} (zk-regex~V2, verifiable
document redaction \cite{vdr}, ZK Email) are building blocks; none of them is a
policy-level system.

The nearest work, WITNESS, supports two policy primitives, has no peer-reviewed
paper, and---by its own project page---ships an Apache-2.0 reference
implementation, so ``open source'' is not a differentiator; it also offers no
content privacy. \texttt{@brivora/verify} is not zero-knowledge, and DeepProve
proves the \emph{model computation} layer rather than the \emph{policy
compliance} layer. The combination of a general policy DSL compiled to ZK
constraints, \emph{without} a TEE, with selective disclosure of response content
and a formal security model, is an explicit vacancy. This paper occupies it with
PoP---but the system is not the contribution we lead with. Leading with a system
invites the objection that a signature would do; leading with the analysis is
what makes the system necessary, and it is the analysis that this paper is
organised around.

\need{⚠️ 上面点名的 Proof-of-Guardrail / Agentrust,cMCP,cA2A / Lean-Agent /
zk-regex V2 / ZK Email / WITNESS / \texttt{@brivora/verify} / DeepProve，
以及 EU AI Act Art.~12/13，\emph{均未给 \cite{}}：条目尚未逐条核实
（docs/tdsc-fixlist.md \S2、WP8）。核实前不得凭空写 bibitem，故此处只留名。}

\smallskip
\noindent\textbf{What we do \emph{not} claim.} We do not claim that policy proofs
bind nothing: our own system already binds the response through a verifier-issued
nonce (Definition~\ref{def:system}). We do not claim to be the first to observe
that a proof can fail to bind the artifact it purports to describe: the
\emph{effort gap} for ZK inference proofs is formalised by Hollow-LLM
\cite{hollowllm}, and the \emph{source-integrity gap} for predicate proofs is
named by zk-attach \cite{zkattach}. And we do not claim the separation theorem
that our designated-verifier interface makes tempting (``public verification is
impossible when $k$ is small''); Remark~\ref{rem:designated} gives the reason,
and Corollary~\ref{cor:reexec} is the short version of it.

\smallskip
\noindent\textbf{Contributions.}
\begin{enumerate}
  \item We isolate the gap that survives in every design we surveyed---\emph{divergent
  delivery}---and show that neither hardware nor any proof closes it, while having
  the receiver contribute to the proven statement does
  (Proposition~\ref{thm:gap}).
  \item We give a quantitative bound on what a public binding reveals, together
  with the \emph{structural hypothesis} that bound needs, and a demonstration
  that the hypothesis is necessary: our own two public values sit on opposite
  sides of it (Theorem~\ref{thm:t1}, Remark~\ref{rem:betanecessary}). The
  governing quantity $k$ is measurable rather than assumed.
  \item We price the fix. A certificate is replayable without limit unless the
  receiver contributes fresh randomness, and with an empty nonce the adversary's
  advantage is $1$ regardless of every other parameter of the system
  (Theorem~\ref{thm:t3}).
  \item We describe PoP, a system that realises this design without a TEE
  (\S\ref{sec:system}), and we report what its certificates leak about the policy
  they certify (\S\ref{sec:private-policy}).
\end{enumerate}

\smallskip
\noindent\textbf{One spine, two sides.} Theorems~\ref{thm:t1} and~\ref{thm:t3}
state the same fact about two different inputs: \emph{a proof binds only as much
as the committed, adversary-non-selectable quantity has bits for}. The response
side is governed by entropy, $k = -\log_2 \mathsf{Guess}(T \mid Z)$
(Definition~\ref{def:k}); the request side by freshness, $f$. Neither can be
bought with more computation.

\fixme{全文禁止出现"比 Proof-of-Guardrail 更强"或任何强弱对比。
Proof-of-Guardrail 与本文正交：它证"计算发生过"，本文证"性质成立"。}

% =====================================================================
\section{System and Threat Model}
\label{sec:system}

The object of study is a \emph{policy-proof system}: a pipeline that compiles a
policy written in a small DSL into a circuit, evaluates an agent's response
against that circuit inside a zkVM, and emits a certificate binding the verdict
to the response. This section fixes the object and its adversary before any
theorem is stated.

\subsection{The system}
\label{sec:system-overview}

\paragraph{Pipeline and the cross-layer contract.} A policy package is compiled
to a \texttt{ConstraintSpec}---a JSON artefact carrying the compiled NFA---which
is the \emph{only} object crossing the layer boundary: the Python reference
implementation and the Rust/zkVM guest consume the same bytes, so cross-layer
drift is structurally impossible rather than merely tested against. The guest
runs a deterministic evaluator and commits its outcome. Of eight rule classes,
seven are decided \emph{inside} the circuit
(\texttt{keyword\_\allowbreak block},
\texttt{normalized\_\allowbreak keyword\_\allowbreak block},
\texttt{length\_\allowbreak bound}, \texttt{pattern\_\allowbreak block},
\texttt{format\_\allowbreak check}, \texttt{tool\_\allowbreak arg\_\allowbreak guard},
\texttt{budget\_\allowbreak bound}). The eighth,
\texttt{semantic\_bound}---a rule of the form ``harm probability $\le$
threshold''---is precisely the kind for which no formal proof can be learned, and
the circuit does not decide it. The circuit registers only \emph{that the rule
was delegated}, in the public value \texttt{delegated}, and a verifier must
\emph{conjunctively} check a companion \texttt{ezkl}/halo2 proof. Verifying the
zkVM side alone does not establish compliance, and we say so here because it is
the most likely misreading of the certificate.

\paragraph{What the verdict is bound to.} Rules that reason about the agent's
\emph{tool calls} cannot be judged against a list the prover fills in itself.
Receipts are therefore issued by a tool gateway after each call and chained
(\texttt{seq}, \texttt{prev}, signature); the circuit checks the chain's
\emph{structure}---entries consecutive, mutually consistent, each digest
recomputed from its own content---and fails closed to \texttt{trace\_unbound}
otherwise. Signature verification happens \emph{off circuit} (Ed25519 inside the
zkVM is prohibitively expensive), a boundary we state rather than gloss: the
circuit is blind to a forgery in the chain's final element, and two independent
off-circuit gates---digest comparison and signature verification---carry that
obligation, neither substituting for the other. Truncating the chain's tail is a
separate and weaker opening, closed by a session-end \texttt{ToolSeal}
commitment from the gateway, and it remains conditional on the verifier holding
the gateway's public key. \S\ref{sec:discussion} develops these residual
boundaries; they are not incidental to the design.

\paragraph{Dual privacy modes.} In \emph{transparent} mode the response and the
verdict are public. In \emph{private} mode the certificate publishes a response
commitment, a commitment per violation, and a redaction proof
(\texttt{redaction\_ok} together with \texttt{mask\_covered}, the latter
requiring the masked positions to fall inside a span the circuit verified to be a
genuine complete match). The guarantee is that \emph{public values contain no
plaintext}---not that $T$ cannot be recovered. A commitment is a public function
of $T$ against which anyone can guess and check, while the response binding must
be publicly recomputable for a verifier to check the string it was delivered.
Those two requirements are in direct tension, and \S\ref{sec:binding} quantifies
exactly when.

\paragraph{Certificates and anchoring.} A certificate is signed with a DSSE
envelope under Ed25519 with the key held by the issuing party, and its digest is
appended to a hash-chain anchor ledger. Two backends share one interface: a file
ledger, verifiable offline, and an on-chain registration storing only a 32-byte
digest. The transaction hash and on-chain timestamp are written back into the
local entry, so the offline chain and the public timestamp cross-check each
other and tampering with either is detectable.

\paragraph{Framework integration.} One gateway serves LangChain callbacks,
LangGraph (\texttt{guard\_node} plus full event streams), and MCP tool argument
and response checks, so trace certificates and proofs share a single chain of
evidence. The implementation passes 801 tests with 15 skipped; the skipped set is
enumerated in Appendix~\ref{app:repro}.

\begin{definition}[Policy-proof system]
\label{def:system}
A policy-proof system is a tuple $(\mathsf{Prove}, \mathsf{Ver}, \pi, \mathcal{M})$.
On challenge $n$ from a verifier, $\mathsf{Prove}(\pi, T, n)$ outputs a
certificate $c$ whose public values include
\begin{multline*}
b \;=\; H(\texttt{pop-bind-v1} \,\|\, \mathsf{u32be}(|n|) \,\|\, n \,\|\, T),\\
C = H(T).
\end{multline*}
$\mathsf{Ver}(c, n, T')$ accepts iff the policy relation holds \emph{and}
$b$ recomputed from the received $T'$ matches.
\end{definition}

\fixme{\textbf{H3/H4}：早先写的 ``the relation contains no term referring to the
delivered bytes'' 与仓库事实冲突——公开值里\emph{已有} $b$，验证时由送达的
$T'$ 现场重算（\texttt{circuits/types/src/lib.rs:444}、
\texttt{docs/security-model.md:232}）。命题已按上面的定义改写。见
\texttt{docs/tdsc-fixlist.md} \S H。}

\subsection{Threat model}
\label{sec:threat}

The adversary is the \emph{agent operator}: the party that holds the policy, runs
the model, and produces the certificate. The verifier is a regulator, auditor, or
third party holding only public values. The policy author issues $\pi$; in this
section we treat it as honest, and \S\ref{sec:wp4-tension} makes that assumption
explicit and scopes it, because in policy-private mode its honesty is doing real
work and cannot be waved through. We state the goals, and more importantly the
non-goals, before the theorems.

\begin{assumption}[Adversary side information]
\label{as:Z}
The adversary knows the prompt, the model (possibly its weights), and the
decoding parameters. \emph{This is the assumption that makes the theorem
non-trivial}: under greedy decoding, $k \approx 0$.
\end{assumption}

\paragraph{Goals.}
\begin{description}
  \item[(G1)] \emph{Soundness}---a response $T$ violating $\pi$ cannot yield an
  accepted certificate.
  \item[(G2)] \emph{Committed privacy}---in private mode the public values
  contain no plaintext of $T$. This is explicitly \emph{not} the claim that $T$
  is unrecoverable: once the guessing entropy $k$ of Definition~\ref{def:k} falls
  below $\log q + \log(1/\varepsilon)$, a publicly recomputable binding and
  content hiding are mutually exclusive (Theorem~\ref{thm:t1}).
  \item[(G3)] \emph{Traceable binding}---policy version, circuit, and proof
  artefact are bound to one another and recomputed at verification.
  \item[(G4)] \emph{Selective disclosure}---redaction and evidence opening are
  themselves provable.
  \item[(G5)] \emph{Record integrity}---tampering with, or reordering, the audit
  ledger is detectable.
\end{description}

\paragraph{Non-goals.} We do not prove the model's inference, and we do not cover
training data or model cards. (G1) is modulo the soundness of the zkVM and of the
hash function; we use those primitives and do not re-derive them.

\paragraph{Where the policy-private case is developed.} The three-party model,
the issuing prerequisite (P-PP), and the table of what each verifier may and may
not conclude are specific to policy-private mode, where the policy itself is
secret. They are stated once, in \S\ref{sec:wp4-tension}, and are deliberately
\emph{not} duplicated here---at a twelve-page limit they cannot be. The general
adversary of this section is the same operator, viewed without the
policy-secrecy refinement.

% =====================================================================
\section{Binding of Policy Proofs}
\label{sec:binding}

A policy proof is meant to bind a response: a verifier who accepts a
certificate should conclude that \emph{this} string was checked against the
policy. This section asks what that binding is actually worth. We first isolate
the gap that survives in every design we surveyed (\S\ref{sec:gap}), then give
a quantitative bound on what a public value reveals, together with the
structural hypothesis that bound needs (\S\ref{sec:tradeoff}), and finally
state the same fact on the request side, where the governing resource is
freshness rather than entropy (\S\ref{sec:t3}).

\subsection{The gap}
\label{sec:gap}

\begin{proposition}[Divergent delivery]
\label{thm:gap}
Definition \ref{def:system} binds $T'$ \emph{as received by the verifying party}. If the
provider delivers a different string $T'' \ne T'$ to the end user, verification
still succeeds on $T'$. Neither hardware nor a proof can prevent this; only the
end user holding $n$ and recomputing $b$ can.
\end{proposition}
\begin{proof}
$b$ is a function of $(n,T)$ only; the delivery channel to the end user carries
no authenticated term. One line.
\end{proof}

Two remarks scope this proposition. First, it describes a \emph{residual} gap,
not the claim that policy proofs bind nothing. Our own system already binds $T$
through the verifier-issued nonce $n$ of Definition \ref{def:system}, and a
verifier who receives $T'$ and recomputes $b$ does learn that $T'$ was the
proved string; the gap is that the binding is anchored to the
\emph{verifying party's} copy. Second, we do not claim to be the first to
observe that a proof can fail to bind the artifact it purports to describe. The
\emph{effort gap} for zero-knowledge inference proofs is formalized by
Hollow-LLM \cite{hollowllm}, and the \emph{source-integrity gap} for predicate
proofs is named by zk-attach \cite{zkattach}; both concern what the proof system
was asked to prove, whereas Proposition \ref{thm:gap} concerns where the proved
string is delivered. We state the naive version of our own claim because it
would be false against our own system.

\subsection{A quantitative binding--hiding tradeoff}
\label{sec:tradeoff}

The certificate must publish \emph{something} tying it to $T$; otherwise it
binds nothing at all. We therefore ask how much that public value leaks. The
answer depends on one measurable quantity --- the guessing entropy of the
response given everything the adversary already knows --- and on one structural
property of the public value itself.

\begin{definition}[Guessing entropy of a response]
\label{def:k}
For a response distribution $\mathcal{D}$ and adversary side information $Z$,
let
\[
\mathsf{Guess}(T \mid Z) \;\triangleq\; \max_{\mathcal{A}} \Pr[\mathcal{A}(Z) = T],
\]
which equals $\mathbb{E}_z[\max_t \Pr[T = t \mid Z = z]]$, and
\[
k \;\triangleq\; -\log_2 \mathsf{Guess}(T \mid Z),
\]
the maximum ranging over all time-unbounded adversaries.
\end{definition}

\begin{remark}[The definition is operational, and it is the average]
\label{rem:kmeas}
$k$ is \emph{measured}, not assumed: sampling one model on one prompt many
times estimates $\mathsf{Guess}(T \mid Z)$ directly
(\S\ref{sec:eval}), so Theorem \ref{thm:t1} is instantiated with a
number we obtain rather than a number we posit. The definition must average
over $z$ rather than take the pointwise worst case
$k_z = -\log \max_t \Pr[T = t \mid Z = z]$: the proof takes an expectation over
$z$, and $\mathbb{E}_z[2^{-k_z}] = 2^{-k}$ holds only for the average form.
\end{remark}

\begin{theorem}[Binding, and the hypothesis it needs]
\label{thm:t1}
Let $\beta : \mathcal{M} \to \{0,1\}^\lambda$ be public, deterministic and
efficiently computable, and let the certificate publish $\beta(T)$. Let
$\mathcal{A}$ be any adversary, unbounded in time, that is given $Z$ and
evaluates $\beta$ at most $q$ times.
\begin{enumerate}
\item[(H)] \emph{(Randomised public value.)} Suppose
  $\beta(T) = H(\mathsf{dom} \,\|\, \mathsf{u32be}(|n|) \,\|\, n \,\|\, T)$ for a
  random oracle $H$ and a nonce $n$ drawn uniformly and \emph{independently of
  $T$}. Then
  \[
  \Pr\!\left[\mathcal{A}^\beta\!\left(Z, \beta(T)\right) = T\right]
  \;\le\; q \cdot 2^{-k} \;+\; q \cdot 2^{-\lambda}.
  \]
  Consequently hiding with advantage $\varepsilon$ requires
  $k \ge \log q + \log(1/\varepsilon)$.
\item[(L)] \emph{(Lossy public value.)} For an arbitrary deterministic public
  $\beta$ the correct bound is over the \emph{preimage class}:
  \begin{multline*}
  \Pr\!\left[\mathcal{A}^\beta\!\left(Z, \beta(T)\right) = T\right]
  \;\le\; \mathsf{Guess}\!\left(T \mid Z, \beta(T)\right) \\
  \triangleq\; \mathbb{E}_y\!\left[\max_t \Pr[T = t \mid Z, \beta(T) = y]\right].
  \end{multline*}
\end{enumerate}
\end{theorem}

\begin{proof}[Proof sketch]
Part (L) is immediate after conditioning on $\beta(T)$; no hypothesis on
$\beta$ is used. For (H), fix $z$ and decompose the adversary's success by the
first query that hits $T$. The oracle answers it has already seen are uniform
and independent of $T$, so each query is itself independent of $T$ and succeeds
with probability at most $2^{-k_z}$. Taking the expectation over $z$ and
applying Definition \ref{def:k} yields $q \cdot 2^{-k}$; the term
$q \cdot 2^{-\lambda}$ charges a query that hits the uniform value $\beta(T)$
by chance. Appendix~\ref{app:proofs} gives the full argument.
\end{proof}

\fixme{$\lambda$ 比特随机与随机预言机只给 \emph{heuristic} 保证：实践中的 SHA-256
不是真的随机预言机（骨架 :712--715）。正文措辞必须写成 heuristic，
\textbf{不得}写成信息论界。}

\begin{remark}[The hypothesis in (H) is necessary, and our two instances straddle it]
\label{rem:betanecessary}
Without the hypothesis, part (H) is not merely unproven but \emph{false}: take
$\beta = \mathrm{id}$, or $\beta = \mathsf{encode}$ as in
Proposition \ref{prop:semantic}. The adversary outputs $\beta(T)$ and succeeds
with probability $1$, independently of $q$ and $k$. Our own two public values
therefore sit on opposite sides of the hypothesis---$b =
H(\mathsf{dom}\|n\|T)$ satisfies it, $\mathsf{encode}(T)$ does not---and this is
what makes the pair a theorem about \emph{our design} rather than a citation of
a standard bound. An earlier draft of this work stated (H) without its
hypothesis; the counterexample above is why the hypothesis is now explicit.
\end{remark}

\begin{corollary}[Re-execution, not enumeration]
\label{cor:reexec}
Under Assumption \ref{as:Z}, when the decoding is greedy or low-temperature,
$k \approx 0$ and the attack of Theorem \ref{thm:t1} is realised by
\emph{re-running the model}, at the cost of one inference---independent of
$|T|$ and requiring \emph{no} evaluation of $\beta$. Enumeration over the string
space is \emph{not} the relevant attack.
\end{corollary}

\fixme{原 docs/sp1-zk-audit.md §4 写的是"熵(~10^2--10^3 bit)远低于 256 bit，
因此枚举可行"——这是非因果推论（2^100 与 2^256 一样枚举不了；真正的攻击是\textbf{重执行}，
对手的已知信息是 (prompt, 模型, 解码参数)）。
2026-10-12 已全库改正，docs/tdsc-fixlist.md §G 的 G1--G6 逐条结案。}

\begin{remark}[A design observation, not a separation]
\label{rem:designated}
Our verification interface is \emph{designated}: binding is checked against a
value $b = H(\mathsf{dom}\|n\|T)$ in which a verifier-issued nonce $n$ appears.
It is tempting to promote this into a separation (``public verification is
impossible when $k$ is small''). \textbf{We claim no such theorem, and
Corollary \ref{cor:reexec} is why.} When $k \approx 0$, hiding fails for
\emph{every} scheme, designated or not, because the adversary re-runs the
model; when $k$ is large, a zero-knowledge proof of ``$\exists\, T : b =
H(\mathsf{dom}\|n\|T) \wedge J(\pi, T) = 1$'' is simultaneously publicly
verifiable and hiding. The dichotomy is thus governed by $k$ alone and carries
no independent content. What does survive---and what \S\ref{sec:discussion}
develops---is that \emph{third-party auditing of delivery} forces disclosure of
$T$, which is an informational remark rather than a theorem. We record the
demotion here because an earlier draft asserted the separation as a theorem.
\end{remark}

We stress that the designated-verifier shape is a \emph{retrospective
justification, not a novel construction}. The system publishes
$b = H(\mathsf{dom}\|n\|T)$ with a verifier-issued nonce $n$, and
Remark \ref{rem:designated} explains the shape of that choice after the fact.
The security model does not pin down whether $n$ is public: making it public
forfeits hiding, while holding it at the verifier is what makes verification
designated.

\begin{remark}[A hypothesis that can be settled by experiment]
\label{rem:kexp}
Theorem \ref{thm:t1} is only as useful as the number put into it. Because
Definition \ref{def:k} is operational, $k$ can be measured rather than assumed:
sampling one model on one prompt many times estimates
$\mathsf{Guess}(T \mid Z)$ directly, and the estimate substitutes into the
threshold $k \ge \log q + \log(1/\varepsilon)$. To our knowledge no prior work
measures this quantity for agent responses, which makes the measurement a
contribution in its own right; \S\ref{sec:eval} describes the protocol and
reports the result.
\end{remark}

\begin{proposition}[The lossy instance of Theorem \ref{thm:t1}(L)]
\label{prop:semantic}
For the semantic rule class, the public values contain $\beta(T) =
\mathsf{encode}(T)$, a content-preserving feature vector of the response. This
$\beta$ is deterministic and injective on the features but not on strings, so
it violates the hypothesis of Theorem \ref{thm:t1}(H) and instantiates part
(L): the binding force is exactly
$\mathsf{Guess}(T \mid Z, \mathsf{encode}(T))$, i.e.\ the guessing probability
\emph{within the preimage class} of the published vector. Here $k$ need not be
argued from perplexity at all: the adversary's search space is that preimage,
which is dramatically smaller than the string space. The relevant asymmetry is
not ``low entropy'' but a \emph{lossy public encoding}.
\end{proposition}
\begin{proof}
Immediate. The measurement is direct: publish $\mathsf{encode}(T)$ for real
responses and measure how much of $T$ is recoverable.
\end{proof}

Proposition \ref{prop:semantic} is the cheapest experiment in this paper and
the strongest anchor for Theorem \ref{thm:t1}: it is measurable in a day, and it
makes the bound concrete rather than cited. It also shows that the relevant
asymmetry is not entropy at all --- for the semantic rule class there is no
entropy argument available, because $\mathsf{encode}(T)$ simply \emph{is} a
content-preserving function of the response.

\subsection{Quantitative request binding}
\label{sec:t3}

\begin{theorem}[Selective-presentation advantage]
\label{thm:t3}
Fix a policy $\pi$ and a request distribution. A selective-presentation
adversary $\mathcal{B}$ may, \emph{before} seeing the receiver's challenge,
compute a set $S$ of accepted transcripts. The receiver then contributes $f$
fresh, uniformly random bits $r$ that enter the proven statement. Then
\[
\Pr\!\left[\mathcal{B}\ \text{is accepted on } r\right] \;\le\; |S| \cdot 2^{-f}.
\]
If instead the receiver contributes no fresh randomness ($f = 0$; e.g.\ an
empty nonce), the advantage is $1$, and \emph{does not decay in any parameter
of the system}.
\end{theorem}
\begin{proof}
The receiver's freshness $r$ enters the proven statement through
$b = H(\mathsf{dom} \,\|\, r \,\|\, T)$. For a certificate in $S$ to be accepted
on $r$, the value $r$ must fall inside the challenge set that certificate
already covers --- an event of probability at most $2^{-f}$ for $f$ fresh
uniform bits --- and a union bound over $S$ gives $|S| \cdot 2^{-f}$. When
$f = 0$ every element of $S$ is a permanently valid replay, so the advantage is
$1$ and does not decay in any parameter of the system.

The $f = 0$ branch is not hypothetical.
\texttt{circuits/\allowbreak types/\allowbreak src/\allowbreak lib.rs}
(lines 950--995) documents an empty nonce as a legal input that ``provides no
replay protection''; it is therefore a \emph{valid input to the current API},
not a straw man.
\end{proof}

Theorem \ref{thm:t3} is \emph{isomorphic} to Theorem \ref{thm:t1}, and that is
what earns it a place here. Both state the same fact about different inputs: a
proof binds only as much as the committed quantity the adversary cannot select
in advance has bits for. On the response side that quantity is the response and
the resource is entropy $k$; on the request side it is freshness $f$. Neither
can be bought with more computation. This is the spine announced in
\S\ref{sec:intro}, and the two theorems are its two sides.

\begin{remark}[A taxonomy, not a trichotomy]
\label{rem:taxonomy}
Every policy-proof deployment we surveyed closes the divergent-delivery gap of
Proposition \ref{thm:gap} in one of three ways:
(a) root trust in hardware (attested enclaves);
(b) have the \emph{end user} contribute a secret to the proven statement (the
user issues the nonce); or
(c) gate delivery on successful verification (\emph{prove-before-deliver}).
\textbf{We offer this as a classification, not as an exhaustiveness theorem}:
we neither prove that a fourth family cannot exist, nor that the three are
disjoint. zk-attach \cite{zkattach} walks (a); NiyamAI and ZKAP walk (c); PoP
walks (b).
\end{remark}

\paragraph{Posture, stated plainly.} Our system already carries
\texttt{response\_binding}, but $n$ is issued by the \emph{verifier}. Moving it
to the \emph{end user}, so that $f > 0$ becomes the receiver's contribution,
closes the divergent-delivery branch of Proposition \ref{thm:gap}. We are
explicit that this is a small protocol amendment rather than a new system: the
mechanism is already implemented, and only the issuer of the nonce changes. What
Theorem \ref{thm:t3} buys is the \emph{price}: with $f = 0$ there is no closure
at any cost.

% =====================================================================
\section{Policy-Private Compliance}
\label{sec:private-policy}

For a policy whose rules are deterministic and a response that is public
anyway, a signed hash chain over the response suffices: the signature already
binds the bytes, and anyone can re-evaluate the rules. The one scenario a
signature cannot cover is a policy that must itself stay secret---where the
operator has to certify compliance to a third party without handing over the
rules it is being judged against. That is the mode this section is about, and we
lead with what the mode \emph{cannot} buy, because that determines whether it
means anything at all.

\subsection{The observation that reframes this work package}
\label{sec:wp4-finding}

\texttt{spec\_canonical} is already a \emph{guest input}---a witness---rather
than a public value; only its hash,
\[
\mathsf{policy\_hash} \;=\; \mathsf{SHA256}(\texttt{spec\_canonical}),
\]
enters the public values (\texttt{circuits/\allowbreak types/\allowbreak
src/\allowbreak lib.rs}:1--16, 415--416, 918, 1003--1004). The instruction
``do not reveal the policy text'' is therefore
\emph{already true of the architecture}; what leaks is everything \emph{derived}
from the policy. Policy privacy is thus not a new proof system to be built but a
\emph{disclosure policy to be audited}, and we write it that way---as a
systematic account of what a certificate reveals about the policy it certifies.

\begin{table}[t]
\centering
\caption{What a policy-private certificate leaks today, and what closing each
leak costs.}
\label{tab:wp4}
\tabfont
\begin{tabular}{@{}L{0.04\linewidth}L{0.46\linewidth}L{0.36\linewidth}@{}}
\toprule
\# & Leak & Difficulty / dependency \\
\midrule
1 & The policy text itself & \textbf{Already closed} (it is a witness). \\
2 & Witness hiding in the proof system & Medium; \emph{inherited from the
    proving tier}: only \texttt{groth16}/\texttt{plonk} is witness-hiding, per
    \texttt{docs/\allowbreak sp1-zk-audit.md}. \\
3 & \texttt{PrivateViolation}\allowbreak\{\texttt{rule},\allowbreak\texttt{kind}\}
    names the rule that fired & Medium; needs a coarsened violation class plus a
    commitment to it. \\
4 & Evidence opening reveals the rule
    (\texttt{policydsl/\allowbreak privacy/\allowbreak commit.py}:319) & Medium;
    disclose the \emph{existence} of a violation, not its identity. \\
5 & The \texttt{delegated} flag exposes policy structure (which rules were
    routed to the \texttt{ezkl} path) & Medium; easiest to overlook, cheapest to
    fix. \\
6 & Semantic rules publish $\beta(T)=\mathsf{encode}(T)$
    (Proposition~\ref{prop:semantic}) & Medium--high; either bound the
    preimage, or disable this rule class in policy-private mode. \\
\bottomrule
\end{tabular}
\end{table}

The honest structure of Table~\ref{tab:wp4} is worth stating, because it sets
the reader's expectations for the rest of the section. Sub-problem~1 is a
one-paragraph re-reading of existing code, not a contribution. Sub-problem~2 is
inherited wholesale from the proving tier and is not specific to this work.
Sub-problems~3--5 are engineering with a single formal core---the granularity of
disclosure. Sub-problem~6 is where the work actually is, and it collides with
\S\ref{sec:tradeoff}: one cannot hide the policy while publishing the lossy
encoding of the response that the semantic rule class requires
(Proposition~\ref{prop:semantic}).

\subsection{The tension worth writing up: policy privacy vs.\ accountability}
\label{sec:wp4-tension}

\paragraph{The prerequisite, stated as a prerequisite.} Let
$H = \mathsf{SHA256}(\pi)$. On its own $H$ is \emph{inert}. A prover can always
emit a self-consistent certificate for any policy $\pi'$ it prefers, simply by
placing $H' = \mathsf{SHA256}(\pi')$ in that certificate. \emph{Secrecy of
$\pi$ is what creates the gap}, and no cryptographic mechanism can close it:
closing it would require showing $\pi$ to the verifier, which is precisely what
the mode exists to avoid. What closes it is an \emph{issuing act}, so we state
it as an explicit prerequisite of the mode.

\begin{description}
  \item[(P-PP)] An \emph{issuing authority} signs $\pi$ and registers $H$ in a
  record readable by third parties. A policy is \emph{accountable} iff its
  \texttt{policy\_hash} can be traced to such a record.
\end{description}

Without (P-PP) the mode is not merely weaker---the verifier's conclusion is
empty. This is the same assumption as the ``honest policy author'' of
\texttt{docs/security-model.md}~\S0.2; what is new is scoping it to the
verification interface, which had not been done, and which a referee would
otherwise use to dismiss the whole mode in one sentence.

\paragraph{Three parties.} ``Three'' means \emph{authority / operator / third
party}; an auditor is a named instance of the third party, not a fourth trust
anchor.

\begin{table}[t]
\centering
\caption{The three-party model of policy-private mode.}
\label{tab:wp4-parties}
\tabfont
\setlength{\tabcolsep}{3pt}
\begin{tabular}{@{}L{0.16\columnwidth}L{0.14\columnwidth}L{0.62\columnwidth}@{}}
\toprule
Party & Honest? & Holds; wants \\
\midrule
Authority $R$ & \textbf{Yes}, by (P-PP) & Holds $\pi$ and the issuing record;
wants $\pi$ actually enforced, and wants to be able to check that it was. \\
Operator $O$ & \textbf{No} (the adversary) & Holds $\pi$, $T$, the proving key;
wants $\pi$ hidden from third parties while still being able to certify. \\
Third party $V$ & Yes & Holds only public values---\texttt{policy\_\allowbreak hash},
\texttt{circuit\_\allowbreak hash}, \texttt{response\_\allowbreak commitment},
\texttt{passed}, the proof; wants to confirm that \emph{some issued} policy
passed, and must \emph{not} learn $\pi$. \\
\emph{(Auditor $V^{*}$)} & Yes & As $V$, plus $\pi$ once $O$ discloses it;
wants ex-post accountability. \\
\bottomrule
\end{tabular}
\end{table}

Two boundaries must be stated with the model. First, secrecy of $\pi$ comes from
a \emph{data-flow fact}---$\pi$ never enters the public values---and not from
zero-knowledge; any rule that writes a fragment of $\pi$ into the public values
(the leaks in Table~\ref{tab:wp4}, items~3 and~5) breaks secrecy on that
fragment. Second, accountability for $V^{*}$ is \emph{voluntary}: $V^{*}$ can
check only when $O$ chooses to disclose $\pi$. The mode offers no forced
disclosure; that would require the release-gating family discussed in
\texttt{docs/tdsc-plan.md}~\S1.6(b), which we do not claim.

\paragraph{What each verifier may conclude.} Table~\ref{tab:wp4-capability} is
the contribution of this subsection: the cryptography is downstream of it, and
the table survives even if the cryptography only partially lands.

\begin{table*}[t]
\centering
\caption{What each verifier may and may not conclude in policy-private mode.}
\label{tab:wp4-capability}
\tabfont
\setlength{\tabcolsep}{4pt}
\begin{tabular}{@{}L{0.13\textwidth}L{0.30\textwidth}L{0.30\textwidth}L{0.19\textwidth}@{}}
\toprule
Verifier & \textbf{May} conclude & \textbf{May not} conclude & Rests on \\
\midrule
Third party $V$ & (i) the commitment binds $T$
(Definition~\ref{def:system}); (ii) \emph{some} policy with hash $H$ passed;
(iii) $H$ appears in the issuing record & (iv) the \emph{content} of $\pi$;
(v) how many rules $\pi$ has, or of which classes; (vi) ``the operator enforces
the regulation \emph{I} care about''---only ``some issued regulation'' &
(P-PP); Assumption~\ref{as:Z}; the $k$ of Theorem~\ref{thm:t1} \\
End user (also receiver) & As $V$, and additionally: the proved $T$ is the
string it received, since its nonce is bound (Theorem~\ref{thm:t3}, requires
$f>0$) & Only the string it received; \emph{not} what was delivered to another
downstream agent & $f>0$ \\
Authority $R$ & Everything; it holds $\pi$ and can recompute $H$ & --- & --- \\
Auditor $V^{*}$ (after disclosure) & (vii) $\mathsf{SHA256}(\pi) =
\texttt{policy\_\allowbreak hash}$, i.e.\ what ran was the issued $\pi$;
(viii) the disclosed violation record is consistent & (ix) the undisclosed part
of $T$ & voluntary disclosure \\
On-chain verifier & (i)--(iii) only; it hashes and verifies signatures, and does
not interpret policy semantics & (iv)--(ix) & (P-PP) \\
\bottomrule
\end{tabular}
\end{table*}

The $V$ row is the reason this subsection exists. A third party does not obtain
``compliance''; it obtains ``some issued policy was satisfied''. Conflating the
two is the sentence a referee would use to reject the mode, and preventing that
sentence is what the table is for.

\paragraph{Prior art.} \emph{Policy hiding is not a new primitive}, and we claim
no primacy for it. Hidden credentials \cite{hiddencred}, secret handshakes
\cite{handshake}, privacy-preserving access control over XACML policies
\cite{xacmlpriv}, and attribute-hiding and fuzzy predicate encryption
\cite{sw05,ksw08} are one family, and policy-private mode is an instance of it.
What is new here is narrow: (a) the \emph{application} to LLM output compliance,
together with selective disclosure of violation type and location, and (b) the
accountability tension of this subsection.

\need{⚠️ 骨架 §4.3「Expected results, stated in advance」（子问题置信度百分比表）
\emph{未}迁入正文。那是 WP 规划的产物，不是论文内容——期刊稿里出现
「我们对子问题 4 有 70\% 把握」是硬伤。若作者要保留这份校准，应改写为
\S\ref{sec:discussion} 里一段关于适用范围与未完成项的讨论，而不是概率表。}

\need{(P-PP) 的落地形式尚未定：链上已签发策略哈希注册表，还是签名清单？
属作者决策。先例引文已按骨架 :623--630 的要求入 bibliography，
但 \cite{hiddencred} / \cite{handshake} / \cite{xacmlpriv} / \cite{sw05} /
\cite{ksw08} 五条\emph{均待逐条打开原文核实}（外部资源阻塞）。}

% =====================================================================
\section{Evaluation}
\label{sec:eval}

\subsection{Cycle count: length versus rule count}
\label{sec:eval-cycles}

Response length was pushed from 200 to 100{,}000 characters and the rule count
from 1 to 6, giving 40 measurement points (\texttt{bench/bench\_cycles.py}; the
full matrix is in \texttt{bench/results/cycles.md}). The response text comes from
an evaluation artefact the repository's own end-to-end demo produced (372
characters, three segments deduplicated by content); longer points are tiled from
it. What is measured is therefore the \emph{curve} of scanning cost, not ``a real
100{,}000-character document''. Every point was first confirmed to match no rule:
a matching rule short-circuits, and the cost of a short-circuit says nothing.

\begin{table}[t]
\centering
\caption{Cycle count against length, per rule set. Five rule sets, six length
points each; $R^2 \approx 1.0000$ throughout.}
\label{tab:cycles}
\tabfont
\begin{tabular}{@{}rrrrr@{}}
\toprule
Rules & \texttt{keyword} & \texttt{pattern} & Slope (cyc/char) & Intercept \\
\midrule
1 & 0 & 0 & 73.9      & 50{,}137 \\
2 & 1 & 0 & 105.3     & 65{,}942 \\
3 & 1 & 1 & 4{,}289.8 & 387{,}501 \\
4 & 2 & 1 & 4{,}048.8 & 405{,}785 \\
6 & 3 & 2 & 5{,}403.5 & 993{,}257 \\
\bottomrule
\end{tabular}
\end{table}

Three findings, in decreasing order of how much they change what may be claimed.

\emph{(1) With the rule set fixed, cycle count is exactly linear in length.}
Five rule sets, six length points each, $R^2 \approx 1.0000$ throughout
(Table~\ref{tab:cycles}).

\emph{(2) Scanning cost is dominated by the NFA, not by the rule count.} One
keyword rule costs about $+31.6$ cycles per character at the margin; one NFA
regular expression (an email pattern) costs about $+3{,}634$---roughly
$115\times$ more. Equivalently, a rule set containing a regular expression has a
per-character slope ($4{,}289.8$) $58\times$ that of a rule set with only string
rules ($73.9$).

\emph{(3) Unit prices are not additive: a unit price is a property of a rule
\emph{in context}, not of the rule.} Fitting the slopes as
$(\text{keyword},\text{pattern})$ unit prices gives $R^2 = 0.9169$ and a
\emph{negative} keyword unit price ($-477.2$ cycles per character). This is not
noise; two probes give the mechanism directly:

\begin{itemize}
  \item \textbf{Ordering probe} (same rules, same text, $L = 10^5$): declared
  order costs $541{,}342{,}069$ cycles against $570{,}736{,}989$ in reverse,
  a difference of $+5.43\%$, while \texttt{passed} is identical in both. \emph{The
  semantics of the decision do not depend on order; the cost does.}
  \item \textbf{Marginal probe} ($L = 10^5$, filling the one matrix cell the
  sweep could not reach---regular expressions with no keyword rule):
  $[\text{len},\text{pat}]$ over $[\text{len}]$ has a marginal unit price of
  $3{,}633.7$, whereas ``add one more keyword rule'' has two different answers:
  $+31.6$ alone, $+585.6$ in a policy that already contains a regular expression
  ($-153.7$ if inserted before it). The same action changes sign between
  $-154$ and $+586$.
\end{itemize}

The mechanism lies in zkVM memory allocation: every \texttt{ascii\_lower} and
every Pike VM match allocates, and allocator state depends on previous
allocations. We therefore quote \emph{per-rule-set slopes} (the table above,
directly citable) rather than any ``unit price $\times$ count'' expression. The
simple model $\text{cycles} \approx a\cdot|T|\cdot\text{rules} + b$ that this
work anticipated is refuted by the data ($R^2 = 0.872$, negative intercept), and
we report that too: \emph{a model refuted by its own data is itself an evaluation
result.}

\subsection{Matcher ablation: quadratic degradation on a pathological input}
\label{sec:eval-ablation}

\emph{Construction.} Pattern \texttt{a+b} against input
\texttt{a}$^{\times n}$ containing no \texttt{b}, so the match is necessarily
absent and both implementations must scan to the end. The naive implementation
reinvents and restarts at every start position, giving $\Theta(n^2)$; the Pike VM
decides in a single linear pass (\texttt{bench/bench\_ablation.py}).

\begin{table}[t]
\centering
\caption{Matcher ablation. Both paths decide identically; only the cost differs.}
\label{tab:ablation}
\tabfont
\begin{tabular}{@{}rrrrr@{}}
\toprule
$n$ & Pike (cycles) & Naive (cycles) & Ratio & Naive growth/char \\
\midrule
100 & 643{,}894   & 12{,}382{,}720  & 19.2$\times$ & --- \\
200 & 1{,}145{,}209 & 48{,}638{,}946 & 42.5$\times$ & $+96.4\%$ \\
400 & 2{,}143{,}617 & 193{,}207{,}130 & 90.1$\times$ & $+98.6\%$ \\
\bottomrule
\end{tabular}
\end{table}

Read per-character cost rather than absolute figures. Doubling $n$ leaves the
Pike VM's per-character cost essentially unchanged ($-11.1\%$, $-6.4\%$) while
the naive path's nearly doubles ($+96.4\%$, $+98.6\%$). At $n = 400$ the naive
path is already $90\times$ slower, and that ratio doubles again with each
doubling of $n$. \emph{This is an attack surface}: the input length for policy
compliance is chosen externally---an agent's reply, a concatenated session---so
an adversary who needs only to write a long reply drives a naive implementation's
cost quadratically. The comparison is meaningful only because the two paths
decide identically, which \texttt{tests/test\_ablation.py} pins down (Python and
Rust \texttt{match\_search} $\equiv$ \texttt{match\_search\_naive}). To our
knowledge this is the first ablation evidence specific to policy-compliance
matching.

\subsection{Real proofs: time, size, memory---and the ceiling}
\label{sec:eval-proofs}

Sampling points share the corpus of \S\ref{sec:eval-cycles}. Each point runs in
its own subprocess (SP1's memory accumulates within a process, so two proofs in
one process do not price one proof), and failures are reported in the table
rather than omitted. Rows marked OOM are a \emph{result}, not a malfunction: they
measure where this machine cannot go.

\begin{table}[t]
\centering
\caption{Real proofs on a 24-core / 12\,GB CPU machine. ``Config'' is
characters $\times$ rules; $k$ abbreviates $10^3$. Every succeeded
point is at the \texttt{core} (non-zero-knowledge) tier.}
\label{tab:proofs}
\tabfont
\begin{tabular}{@{}lrrrl@{}}
\toprule
Config & Time & Size & Peak RSS & Result \\
\midrule
$200 \times 1$  & 123.5\,s & 2716.4\,KiB & 10{,}389\,MB & \checkmark \\
$200 \times 2$  & 118.9\,s & 2716.9\,KiB & 10{,}438\,MB & \checkmark \\
$2k \times 1$   & 138.1\,s & 2718.6\,KiB & 10{,}449\,MB & \checkmark \\
$10k \times 1$  & 172.5\,s & 2726.6\,KiB & 10{,}506\,MB & \checkmark \\
$200 \times 3$  & ---      & ---          & ---          & OOM \\
$20k \times 1$  & ---      & ---          & ---          & OOM \\
\bottomrule
\end{tabular}
\end{table}

\begin{itemize}
  \item \textbf{A fixed floor of $\approx$10.15\,GiB.} The smallest
  configuration---200 characters, one rule---already peaks at 10{,}389\,MB: the
  prover's own overhead (the trace and permutation structures of a \texttt{core}
  proof) is almost independent of the trace. All four successes fall inside the
  narrow band 10{,}389--10{,}506\,MB. Length and rule count have little marginal
  effect on memory---\emph{but the headroom is too thin for that to help}.
  \item \textbf{A two-step staircase, not a slope.} One rule is provable up to
  10k characters; two rules at about 200 characters; \textbf{three or more do not
  come out at all}. The third rule in the sweep is a \texttt{pattern\_block}: a
  regular-expression match activates another family of AIR chips and lifts the
  trace area in one step, so the cliff is at $2 \to 3$ rules; in length it is at
  $10\text{k} \to 20\text{k}$.
  \item \textbf{Memory binds, not CPU.} The cycle sweep of
  \S\ref{sec:eval-cycles} reaches 100k characters with 6 rules because it only
  executes in the zkVM and produces no proof; at 100k characters the cycle count
  is still on the order of $10^7$, which any CPU handles. \emph{The proving-side
  ceiling and the execution-side ceiling are in different places.} ``How large a
  policy can be proved'' is, on this machine, a fixed floor plus two steps rather
  than a cost curve.
  \item \textbf{Artefact size is nearly independent of configuration}: from 200
  to 10{,}000 characters the proof grows only from 2716.4 to 2726.6\,KiB. Pure
  verification takes $\approx$90\,ms (the \texttt{vkey} setup is 1.5\,s, once).
  \item \textbf{Reproducibility, stated precisely.} The \emph{position} of the
  wall reproduces (each failing point failed twice); the \emph{timings} at the
  wall do not. Re-running one group in isolation shows $\pm 10\%$ jitter in time
  and $\pm 1\%$ in memory, so we quote time only to an order of magnitude.
  \item \textbf{Session and multiparty.} The session aggregation proof (three
  certificates) measured 152.5\,s; multiparty verification needs one proof
  \emph{per non-empty slice}, so its cost grows linearly in the number of slices
  (a two-proof example policy). Neither introduces a new proving system---a slice
  and a session are ordinary inputs to the guest---but both are subject to the
  same floor: every additional proof adds another $\approx$10\,GiB peak, so
  multiparty proving must run in separate processes, serially.
\end{itemize}

\need{⚠️ 本节的 groth16 档与 ZK/非 ZK 交叉点属 WP1，尚未打通。骨架 :665--668 的
矩阵规格为规则 $\{1,2,3,4,6,8\}$ $\times$ 长度 $\{200,2k,10k,50k\}$ $\times$
模式 $\{$transparent, private, private-policy$\}$ $\times$ 证明档 $\{$core, groth16$\}$
（\texttt{bench/bench\_scale.py}）。表 \ref{tab:proofs} 只覆盖 core 档。
groth16 列落地前不得填入任何 ZK 数字。}

\subsection{The non-ZK baseline: what a signature cannot do}
\label{sec:eval-baseline}

A signed Merkle log is the obvious alternative to a proof, and on cost it is not
close: everything is milliseconds. We measured it so that the comparison is not
rhetorical. Policy package \texttt{eu\_ai\_act\_v1.json}, 103-character response,
one challenge nonce, a 1024-leaf log; each timing is the median of $n = 500$
repetitions.

\begin{table}[t]
\centering
\caption{Signed-Merkle baseline (WP9). No proofs, no zkVM: this arm is pure
Python, which is why the figures are milliseconds.}
\label{tab:baseline}
\tabfont
\begin{tabular}{@{}lrl@{}}
\toprule
Metric & Median & Note \\
\midrule
Sign certificate (Ed25519 + DSSE) & 0.132\,ms & $n=500$ \\
Merkle append (leaf + root recompute) & 2.131\,ms & $n=500$ \\
Public verification (sig + inclusion) & 0.322\,ms & $n=500$ \\
Re-run verification ($+J(\pi,T)$) & 0.772\,ms & $n=500$ \\
Certificate size & 648\,B & \\
Inclusion proof (10 siblings) & 761\,B & \\
\bottomrule
\end{tabular}
\end{table}

\need{⚠️ 引用 2.131 ms 时必须写明它是 \texttt{proofs/session.py} 的\emph{静态树重算}
$\Theta(n)$，\emph{不是} RFC-6962 增量日志的摊销 $\Theta(\log n)$。
⚠️ 共享机器有 $\pm 20\%$ 漂移。数据源 \texttt{bench/results/baseline.md}。}

The asymmetry is the point, and it is what this arm contributes. The verification
chain is \emph{verify signature} $\to$ \emph{verify log inclusion} $\to$
\emph{re-run $J(\pi,T)$}. The first two steps need only public data; the third
needs the plaintext of both $\pi$ and $T$. That is not an implementation defect
but a structural consequence of the signature paradigm: a signature binds
\emph{who said it} to the bytes, never \emph{whether it is true}. We measured the
consequence rather than asserting it---\textbf{the same key signs a certificate
whose declared content is false, and the public path passes green}, because what
it sees is a certificate with a valid signature and a valid inclusion proof. Only
re-executing $J(\pi,T)$ recovers the truth, and taking that step away from the
verifier is exactly what the zero-knowledge design buys. The machine-checked
assertions are enumerated in \texttt{bench/results/baseline.md}; one row of that
file is worth repeating because it is the honest one: on
\emph{unforgeability} the baseline is \emph{not weaker} than the proof path, both
being Ed25519 and vkey binding, and we do not score a point for ZK there.

\subsection{End-to-end and auditing}
\label{sec:eval-e2e}

The end-to-end demo produces 14 certificates (streaming chain with early stop,
MCP argument and response checks, one zk proof) and anchors them; a third party
verifies with public artefacts alone, and the ledger chain, signature,
\texttt{policy\_hash}, anchoring, stream chains and zk proof all pass. In the
on-chain variant (a local Anvil node, \texttt{Anchor.sol} deployed, 14
certificate digests registered) the verifier reads back each certificate's
\texttt{anchoredAt} and cross-checks the local ledger, passing 14/14; the
negative control (an unregistered digest reads back as zero) confirms the check
is not vacuously true. Cross-validation submits 19 vectors to the Python
reference decision and to both the SP1 host and proving paths, comparing
\texttt{passed} and the violated rule set row by row; a full re-run with real
\texttt{core} proofs gave 19/19 on both paths.

\subsection{Measuring \texorpdfstring{$k$}{k}}
\label{sec:eval-k}

\need{对一组真实 prompt、同一模型多次采样，估 $\mathsf{Guess}(T\mid Z)$，
代入 Thm \ref{thm:t1}(H) 的阈值 $k \ge \log q + \log(1/\varepsilon)$（承接
Rem \ref{rem:kexp}）。据先例扫查无人测过这个量，本身是一块贡献；
同时它也是 Prop \ref{prop:semantic} 的测量（发布 $\mathsf{encode}(T)$ 后
$T$ 的可恢复比例）。\textbf{⚠️ 需真实模型多次采样，当前阻塞（外部资源）。}}

\need{⚠️ 「Why zero-knowledge?」决策图（骨架 :670--675）\emph{未}放入本节：
它讲的是设计结论（ZK 只在 $\{$私有响应, 私有策略$\}$ 两角有正收益），是结论不是测量，
按 docs/tdsc-writing-plan.md 的建议移至 §7 Discussion。与骨架现状不同，\textbf{待作者认可}。}

% =====================================================================
\section{Related Work}
\label{sec:related}

We position PoP against five lineages. For each, Table~\ref{tab:related} records
the slot it occupies---which is also, read the other way, what this paper does
not have to claim.

\begin{table}[t]
\centering
\caption{Positioning against the five prior-art lineages. The right-hand column
names what each lineage already occupies, and therefore what we do not claim.}
\label{tab:related}
\tabfont
\setlength{\tabcolsep}{3pt}
\begin{tabular}{@{}L{0.26\columnwidth}L{0.42\columnwidth}L{0.26\columnwidth}@{}}
\toprule
Lineage & Representative work & Slot already occupied \\
\midrule
Policy $\to$ circuit & zkGen \cite{zkgen}, Circuitree, Prezta & DSL-to-circuit
translation \\
Agent/LLM compliance & NiyamAI, Zaguk, PCO \cite{pco}, WITNESS & policy-level
private-mode kernel \\
Model computation & DeepProve, zkGPT, zkAgent \cite{zkagent} & inference
integrity \\
Attacks / negative results & Hollow-LLM \cite{hollowllm}, zk-attach
\cite{zkattach} & naming the gaps \\
Non-ZK industry baseline & \texttt{@brivora/verify}, VAOL & signed audit logs \\
\bottomrule
\end{tabular}
\end{table}

One of the five is the lineage this paper is actually in dialogue with. The
attacks-and-negative-results row names two gaps---Hollow-LLM's \emph{effort gap}
and zk-attach's \emph{source-integrity gap}---and both concern \emph{what the
proof system was asked to prove}. Proposition~\ref{thm:gap} concerns a third
gap on a different axis, \emph{where the proved string is delivered}, and we are
not aware of prior work that quantifies the binding/hiding tradeoff in the terms
of \S\ref{sec:binding}. That is why we state it as a bound with a stated
hypothesis rather than cite it.

\paragraph{zkAgent: orthogonal, not competing.} zkAgent \cite{zkagent} is the
closest contemporaneous system, and the comparison is easy to get wrong in both
directions. It proves that a provider executed the declared model and tools; PoP
proves that a response satisfies a public policy. PoP proves nothing about model
computation, and zkAgent proves nothing about policy compliance, so the two
compose: their conjunction covers ``inference was executed \emph{and} the policy
held''. Table~\ref{tab:zkagent} places the reported figures side by side.

\begin{table}[t]
\centering
\caption{PoP against zkAgent. The two prove different obligations, so the rows
are a statement of orthogonality, not a race.}
\label{tab:zkagent}
\tabfont
\setlength{\tabcolsep}{3pt}
\begin{tabular}{@{}L{0.22\columnwidth}L{0.38\columnwidth}L{0.30\columnwidth}@{}}
\toprule
& \textbf{zkAgent} \cite{zkagent} & \textbf{PoP} \\
\midrule
Obligation & provider executed the declared model and tools (inference integrity)
& response/trace satisfies policy $\pi$ (compliance) \\
Depends on the model & yes (weights, arithmetisation) & \textbf{no} \\
Cost grows with & model size, sequence length & response length $\times$ rule
count \\
Content privacy & none & committed, selective disclosure \\
Hardware reported & 32 cores / 512\,GB & 24 cores / 12\,GB \\
Proving & 194.4\,s (Shout, $T{=}512$) & 118.9--172.5\,s (200--10k chars) \\
Proof size & 96\,MiB (Shout); 3.1\,MiB (LogUp) & $\approx$2.7\,MiB \\
Verification & 0.42\,s (Shout); 0.038\,s (LogUp) & $\approx$0.090\,s \\
\bottomrule
\end{tabular}
\end{table}

We deliberately do not claim that PoP is faster, and the temptation to do so is
worth naming because the raw numbers invite it. The correct statement is
narrower, and we make only it: \emph{for the compliance obligation}, PoP
completes in a proving time of the same order as zkAgent's, with a smaller memory
footprint and a comparable proof size, without touching the model---and it adds
committed content privacy that zkAgent does not provide. Two caveats must travel
with those figures. The hardware is not matched (24 cores and 12\,GB against 32
cores and 512\,GB), and PoP's ceiling is memory-bound rather than cycle-bound, so
its feasible region is a two-step staircase rather than a curve.
\S\ref{sec:eval} reports the measurements and their limits.

\need{⚠️ tab:zkagent 里 PoP 那一列全部取自 SP1 \texttt{core}（非零知识）档，
机器为 24 核 / 12\,GB WSL2；zkAgent 一列为 ePrint 2026/199 的报告值（32 核 / 512\,GB）。
数据源：\texttt{bench/comparison\_zkagent.md}。该文件自身警告：
\emph{不得}宣称「PoP 比 zkAgent 快」——正确口径是「对『合规』这一义务…
同量级证明时间 / 更小内存」。groth16 档（WP1）未打通前，本表不得出现 ZK 档数字。}

The one influence we should name plainly, rather than leave to resemblance:
zk-regex's ``match offline, verify the path in circuit'' pattern is the shape of
our NFA contract, and verifiable document redaction \cite{vdr} inspired the bit
selector behind \texttt{mask\_covered}. The borrowing stops at the circuit
technique and does not carry the security argument across: VDR assumes a
\emph{signed original document} exists, whereas in PoP's private mode both the
original and the redacted version are produced by the operator, so there is no
trusted original to appeal to. We use the mechanism, not the proof.

\fixme{引文纪律：除 Hollow-LLM \cite{hollowllm} / zkGen \cite{zkgen} /
zkAgent \cite{zkagent} / PCO \cite{pco} 四条外，其余引文必须逐条打开原文
核实后再进 bibliography（docs/tdsc-fixlist.md \S2、WP8）。
\textbf{骨架原先把「ePrint 2024/813」列进已核实四条，那是错的}，已剔除：
\textbf{绝不能再引 IACR ePrint 2024/813 作为 VDR}——那是
Kiraz/Larraia/Vaughan, ``How to Redact the Bitcoin Backbone Protocol''，
与 Verifiable Document Redacting 无关。见 fixlist \S A。
⚠️ 口径：WITNESS 参考实现为 Apache-2.0（并非闭源），故「开源」不构成差异化。}

% =====================================================================
\section{Discussion and Limitations}
\label{sec:discussion}

Everything above was kept deliberately narrow, and this section says how narrow.
It covers the four places where a reader is most likely to round our claims up:
the caliber of the privacy statement (\S\ref{sec:disc-privacy}), where
zero-knowledge actually pays (\S\ref{sec:disc-zkmap}), availability against an
adversarial response (\S\ref{sec:disc-dos}), coverage
(\S\ref{sec:disc-coverage}), and composition with an inference proof
(\S\ref{sec:disc-composition}). It closes with the engineering limitations that
remain (\S\ref{sec:disc-rest}).

\subsection{The caliber of the privacy claims}
\label{sec:disc-privacy}

\textbf{``Content privacy'' means \emph{no plaintext of $T$ in the public
values}, not \emph{content is unrecoverable}.} Two independent reasons force
this reading---one at the proof-system level, one at the design level---and the
weaker reading is the one our own measurements support.

\emph{Proof-system level.} The default proving tier is SP1's \texttt{core}, and
SP1's published security model states explicitly that individual STARK proofs do
not currently satisfy the zero-knowledge property; the source agrees, with no
blinding anywhere in the proof stack and the shard proof carrying the trace
Merkle root and trace openings in the clear.%
\need{⚠️ 引文债：「SP1 公开安全模型承认其 STARK 证明当前不满足零知识」是
\emph{外部}断言，出处（SP1 官方安全模型文档的具体章节/URL）尚未核实到可引的
一级来源，故此处只陈述、不给 \cite{}（WP8）。若核不到原文，须撤回该句，
只保留我们自己的源码审计结论那半句。}
Soundness (G1) therefore holds
while the proof artefact itself does not hide the witness. The
\texttt{gr\-oth16} and \texttt{plonk} tiers wrap the inner STARK proof as a
private witness---the only mode in which hiding is even possible---but that is a
wrapper-level claim, it is not post-quantum, it needs $\ge 16$\,GB, and we did
not measure it (\S\ref{sec:eval-proofs}).

\emph{Design level, and the more fundamental one.} The public values
\texttt{response\_binding} and \texttt{response\_commitment} are both
\emph{public, offline-recomputable functions of $T$}. This is not an oversight:
\texttt{response\_binding} \emph{must} be publicly recomputable, because that is
how a verifier independently checks that the delivered $T'$ is the $T$ that was
proved---the swap-response defence. But the same property is a guess-and-verify
oracle for $T$.

It is worth being precise about the attack, because the obvious version of the
argument is wrong. ``A one-line response carries only $10^2$--$10^3$ bits of
entropy, far below 256, so enumeration is feasible'' does \emph{not} follow:
$2^{100}$ is as unenumerable as $2^{256}$. The real attack is
\emph{re-execution} (Corollary \ref{cor:reexec}). The adversary's side
information is not the response space but the tuple (prompt, model, decoding
parameters) (Assumption \ref{as:Z}); under greedy or low-temperature decoding
the response is nearly determined, so $k = -\log_2 \mathsf{Guess}(T \mid Z)
\approx 0$ (Definition \ref{def:k}) and the attack costs \emph{one inference},
independent of $|T|$. Under the two premises that a verifier independently
recomputes the binding and that $k < \log q + \log(1/\varepsilon)$, response
binding and response-content hiding are mutually exclusive.

We therefore restate (G2) as what is actually established: \emph{in private mode
the public values contain no plaintext of $T$}---which the leak experiment
confirms directly---and not as content irrecoverability. The consequence for the
comparison in \S\ref{sec:related} is that the ``content privacy'' cells on
\emph{both} sides must not be read as ``content is unrecoverable.'' This is a
bound on our own system first.

\subsection{Where zero-knowledge pays, and where it does not}
\label{sec:disc-zkmap}

Section \ref{sec:eval-baseline} measured the non-ZK alternative---a signed
Merkle log---and everything in it is milliseconds. A proof therefore has to earn
its cost against what a signature cannot do, and our measurements locate that
boundary instead of asserting it.

The map has two axes: is the response secret, and is the policy secret?
Zero-knowledge has positive value only in the corners where at least one of them
is unknown to the verifier. In the transparent corner---public response, public
policy---a signed log certifies exactly the same thing at three to four orders
of magnitude less cost, and adding a proof is redundancy rather than rigour. The
familiar criticism that ``a signature would have done'' is correct in that one
corner and wrong in the other three; the decision map is what turns the
criticism into a measured design conclusion.

One dependency belongs in the open. The zero-knowledge property of a certificate
is a property of the \emph{wrapper}, not of the system as a whole, and the
default tier does not have it (\S\ref{sec:disc-privacy}). If \texttt{gr\-oth16}
does not become a delivery tier, we do not claim the artefact is
zero-knowledge, and the word must leave the title.

\subsection{Availability: the unanchored-pattern cost attack}
\label{sec:disc-dos}

A regular expression with an unanchored pattern---one with no left anchor, such
as \texttt{.*foo}---gives the matcher no way to reject a prefix of the response
early: every position stays a live start position until the whole string has
been read. The prover's work therefore scales with the response, and a party who
\emph{chooses} the response can drive it quadratically. In the agent setting
that party is often the agent itself, or whoever controls its output.

Two things are worth separating. This is a denial of service against the
\emph{prover}, not an attack on soundness: no false certificate is produced, and
the cost is simply borne by the honest side. And it is specific to proving
rather than to matching, because the trace---not merely the verdict---is what
gets proved. Mitigations exist and are partly in place: a Pike VM (already
implemented) simulates the NFA in linear time, committing to the response length
up front removes the adversary's freedom to inflate it, and chunked matching
bounds the per-chunk work. Prior work on zero-knowledge regular expressions
imposes the DFA restriction but does not analyse denial of service in a zkVM
setting, which is where this differs.

\need{⚠️ 作者决策未定：骨架 :649--651 把这条写成一条定理（$\Omega(n^2)$ 证明
工作量），但证明没写。需定：写成定理（并补证明），还是降为观察（Observation）。
本稿先按观察陈述、不给定量断言，以免留下一条无证明的定理。}

\subsection{Coverage: a provider can simply withhold certificates}
\label{sec:disc-coverage}

Nothing in the cryptographic layer can compel a provider to issue a certificate
at all. A provider that never produces one is indistinguishable, from the
verifier's side, from a provider with nothing to certify. We treat this as an
open problem rather than a defect of the construction: without trusted hardware
there is no anchor that forces issuance, and the mechanisms that could---%
sampling audits, economic penalties, disclosure obligations---live above the
cryptographic layer and outside the scope of this paper. Naming it as an open
problem is more useful than proposing a mitigation we cannot evaluate.

\subsection{Why composition alone does not close the gap}
\label{sec:disc-composition}

The natural next step is to compose: prove that the inference produced $T$ from
the request, conjoin that with a proof that $\pi(T)$ holds, and conclude that
the agent \emph{sent} a compliant response. We describe the mechanism, and we
record why the conclusion does not follow for free.

\emph{Composition by itself does not close the gap.} A model proof bound to
\emph{some} input shows only that \emph{some} input yields a compliant output;
the divergent-delivery gap of Proposition \ref{thm:gap} closes only when the
model proof is bound to the \emph{input the user submitted}, $x$, so that
$M(x) = T \wedge \pi(T)$ becomes a statement about the bytes actually sent. That
requirement is not an extra condition bolted on: it is exactly $f > 0$ in
Theorem \ref{thm:t3}, with the user-submitted $x$ playing the role of the $f$
fresh bits. The two halves of this paper meet here---the request-side theorem
names the condition the composition would have to satisfy, and the response-side
theorem is what makes that condition non-vacuous.

The composition contract is correspondingly short to state: the inference proof
must expose $T$ in its public values, and it must be bound to $x$.

\need{⚠️ 阻塞：真模型规模（$\ge 10^6$ 参数）与组合成本数据属 WP4 子问题② /
WP3 的 DeepProve 对照臂，需外部资源。代理规模下「组合成本由推理证明主导」
\emph{未被验证}（\texttt{bench/results/compose.md}），不得写入。
骨架 :641--642 的组合契约定义可在此处补。}

\subsection{Remaining limitations}
\label{sec:disc-rest}

\begin{itemize}
  \item \textbf{Three-party collusion is undetectable at the signature layer.}
    If the policy author, the operator, and the prover all collude, the
    certificate is indistinguishable from an honest one. The multi-prover
    construction partitions responsibility; it does not prevent agreement
    (\texttt{docs/security-model.md}).
  \item \textbf{Signing and key custody are not production-grade.} The on-chain
    path currently takes plaintext private-key parameters, the demo uses a
    public Anvil test key, and deployment is a local Anvil instance. Production
    signing (Ed25519/HSM) and key custody remain future work.
  \item \textbf{Semantic-rule model quality.} The delegated semantic rule works
    end to end, but the bundled model is a small demonstration proxy rather than
    a usable learned guardrail: ``$P(\text{harmful})$ is low'' is a statement
    about \emph{this} model. The companion proof's public instance also contains
    $\mathsf{encode}(T)$, so semantic rules are \emph{public-mode only}.
  \item \textbf{The inference proof is faithful, not good.} The inference half
    is a deterministic fixed-point MLP proxy whose weights are generated from a
    compile-time seed and baked into the program, hence committed by the
    verification key. It proves that \emph{this} graph computed \emph{this}
    output on \emph{this} response---not that the graph is any good. Swapping the
    proxy for a real prover touches only the prover; the composition layer does
    not move.
  \item \textbf{Verification remains $O(N)$.} There is no recursive aggregation,
    so a verifier pays per certificate. The supported regex subset and its ASCII
    semantics are bounded in the same way, and proof-cost optimisation (lookup
    arguments, parallelism, precomputation) is unexplored.
  \item \textbf{Slices share one guest.} Multiple provers hold distinct keys,
    but that is a division of responsibility rather than a per-role circuit: the
    repository does not build one circuit per role.
\end{itemize}

% =====================================================================
\section{Conclusion}
\label{sec:conclusion}

This paper began from a small question---what does a compliance certificate
actually bind?---and the answer came out the same on both sides of the
interaction. A proof binds only the bits that were submitted into it \emph{and}
that the adversary could not select in advance. On the response side those bits
are the entropy $k$ of $T$ given everything the adversary already knows, and no
amount of computation buys more of them; on the request side they are the $f$
bits of freshness the receiver contributes, and a certificate without them is
replayable without limit. Theorem \ref{thm:t1} and Theorem \ref{thm:t3} are two
faces of that one statement, and Proposition \ref{thm:gap}---a gap we state
against our own design---is what makes it bite.

PoP is the constructive instance of that analysis, not a claim to have invented
a new proof system. It composes an existing zkVM, an existing policy compiler,
and a signature layer into a certificate that a verifier holding only public
values can check without trusting hardware, and it is explicit about what such a
certificate does not do: it does not make the response unrecoverable, it does not
hide the witness at the default proving tier, and it cannot compel a reluctant
provider to issue anything at all. Those are limits of the problem, stated as
limits, rather than of one implementation.

The design conclusions are accordingly modest, and we hope reusable: publish the
leak profile of a certificate and audit it like a disclosure policy; let a
privacy claim be only as strong as its arithmetic supports; and put the
freshness where the binding is. A certificate is worth exactly as much as the
bits that went into it.

% =====================================================================
% ⚠️ IEEEtran 的 \appendix 是「单个附录」命令：它把 \section 重定义成
% \@IEEEdestroythesectionargument，直接丢弃 \section{...} 的参数（只打 warning）。
% 多个附录必须用 \appendices（复数），它把 section 计数器置 -1，
% 之后每个 \section{X} 才依次成为 A、B、……  （2026-10-12 实测：
% 用 \appendix 时 A/B 两个附录的 \ref 都解析成 "A"，且 title 字段退化成 \appendixname，
% 锚点退化成 section*.1；「0 undefined ref」检查抓不到这种错。）
\appendices
\section{Proofs}
\label{app:proofs}

\begin{proof}[Proof of Theorem \ref{thm:t1}]
\emph{(H).} Fix $z$ and let the adversary's query set be
$Q = \{q_1, \dots, q_q\}$. Decompose by first hit,
\[
\Pr[\exists i : q_i = T] \;=\; \sum_i \Pr\!\left[q_i = T \wedge \neg E_{<i}\right],
\]
where $E_{<i} \triangleq \{\exists j < i : q_j = T\}$.
Conditioned on $\neg E_{<i}$, every oracle answer $a_1,\dots,a_{i-1}$ already
revealed is uniform and independent of $T$ (no queried point equals $T$, and $H$
is a fresh random function), and $\beta(T)$ itself is uniform and independent of
$T$. Hence $q_i$---a deterministic function of $(z, \beta(T), a_{<i})$---is
\emph{independent of $T$}, so
\begin{multline*}
\Pr[q_i = T \mid Z = z]
= \sum_t \Pr[q_i = t]\,\Pr[T = t \mid Z = z] \\
\le \max_t \Pr[T = t \mid Z = z] = 2^{-k_z},
\end{multline*}
giving $\Pr[\exists i : q_i = T \mid z] \le q\,2^{-k_z}$. Taking the expectation
over $z$ and using $\mathbb{E}_z[2^{-k_z}] = \mathsf{Guess}(T \mid Z) = 2^{-k}$
from Definition \ref{def:k} yields $q \cdot 2^{-k}$. An adversary that outputs
an unqueried value wins with probability at most $2^{-k_z} \le q\,2^{-k_z}$,
which the same bound absorbs. The term $q \cdot 2^{-\lambda}$ charges the
probability that a query hits the uniform value $\beta(T)$ by chance.

\emph{(L).} Immediate from Definition \ref{def:k} after conditioning on
$\beta(T)$; no hypothesis on $\beta$ is used.
\end{proof}

\need{Thm \ref{thm:t1}(L) 无闭式界，只给测量口径（\S\ref{sec:eval}）。
三分法已降级为分类 Rem \ref{rem:taxonomy}，无需证明。
原 T2「分离定理」已\emph{停用}（见 Rem \ref{rem:designated}）—— 不要为它写附录。}

\need{\textbf{待补}：把 Thm \ref{thm:t1}(H) 的随机预言机模型换成标准模型，
或明确声明本文只在 ROM 下主张。当前证明依赖 $H$ 是新鲜随机函数这一条；
实践中的 SHA-256 只给计算意义下的近似，故正文措辞必须写成
\emph{heuristic}，不得写成信息论界。}

\section{Reproduction}
\label{app:repro}

The full walkthrough---environment, a single compliance proof, cross-checking
against the reference layer, private mode, certificate issuance, the one-command
demo, and third-party verification---is maintained in \texttt{docs/reproduce.md}
together with the exact commands and their expected output. This appendix records
only what a reader needs in order to judge the claims of \S\ref{sec:eval}.

\paragraph{Environment.} Ubuntu or WSL2 on x86-64, no GPU. The Rust/SP1 side
needs Rust stable, \texttt{cargo-prove}~v6.7.0, Go $\ge$\,1.24, and the usual
native toolchain; the reference layer is pure Python~3.10+ and runs on its own.
Memory, not time, is the binding constraint: the CPU prover carries a fixed
floor of approximately 10.15\,GiB regardless of the trace
(\S\ref{sec:eval-proofs}), and 12\,GB suffices only for small policies.

\paragraph{What reproduces without the prover.} The reference layer, the policy
compiler, and the whole test suite run in seconds and need no Rust. This is
enough to reproduce the non-ZK baseline of \S\ref{sec:eval-baseline}, to inspect
the certificate and inclusion proofs, and to check the cross-layer contract---%
that the Python reference implementation and the zkVM guest consume the same
compiled policy bytes.

\paragraph{What needs the prover.} Every timing, size, and memory figure in
\S\ref{sec:eval-cycles}--\S\ref{sec:eval-proofs}, and the end-to-end result of
\S\ref{sec:eval-e2e}, requires a real proof. The measurements were taken on a
24-core CPU machine with 12\,GB of RAM, which is what places the ceiling reported
there: the successes and the two out-of-memory configurations in
Table~\ref{tab:proofs} are both properties of that machine.

\paragraph{Where each number comes from.} Every measurement in \S\ref{sec:eval}
is the content of a checked-in result file with a machine-readable companion:
\texttt{bench/results/} holds \texttt{cycles}, \texttt{ablation},
\texttt{proofs}, \texttt{baseline}, \texttt{verify}, \texttt{semantic},
\texttt{streaming}, \texttt{compose}, and \texttt{prover\_knobs}, each as a
\texttt{.json} alongside its \texttt{.md}. The reproduction guide records the
commit, the hardware, and the prover binary digest for each run, so that a
reported figure can be traced to one specific execution rather than to a state
of the tree.

\need{⚠️ 数字漂移（本轮实测发现，需作者定夺）：正文 §2.2 写「801 tests / 15 skipped」，
其依据是实测（\texttt{docs/dev-plan.md:506} 的落地记录）\emph{以及}本轮对
\texttt{tests/} 的静态清点（\texttt{def test\_} 计 801 个）——\textbf{801 是对的}。
但 \texttt{docs/reproduce.md:31,428}、\texttt{docs/security-model.md:829}、
\texttt{docs/development.md:93,290}、\texttt{docs/modules/08-tests-bench.md:11,58,504,526}、
\texttt{docs/modules/README.md:94,161,201} 仍写 \textbf{783}（陈旧值，
\texttt{docs/report/advisor-report.md:586} 已把这条「文档数字漂移」登记在案）。
本附录既然指向 \texttt{reproduce.md}，照抄就会让复现者跑出 801 却以为环境坏了——
\textbf{这是真实的复现缺陷}。按仓库自己的「过期就标注、不静默改」纪律，
本稿不擅自改那些文档，只在此记债。作者需定：把五处文档回填为 801
（\texttt{advisor-report.md:599} 的建议），还是在附录里并列两套口径。}

% =====================================================================
\begin{thebibliography}{99}

\bibitem{hollowllm}
C.~Gong, B.~Liu, and M.~Li, ``Hollow-LLM Attack: Computationally Trivial
Weights in Zero-Knowledge Verification of LLM Inference,'' in \emph{IEEE S\&P},
2026, pp.~4339--4355, DOI 10.1109/SP63933.2026.00258.

\bibitem{zkattach}
\need{zk-attach, arXiv:2608.30083 --- 打开原文核实作者后再引}

\bibitem{zkgen}
J.~Lauinger, J.~Ernstberger, and S.~Steinhorst, ``zkGen: Policy-to-Circuit
Transpiler,'' in \emph{IEEE ICBC}, 2024, pp.~620--624,
DOI 10.1109/ICBC59979.2024.10634440.

\bibitem{pco}
\need{PCO, IACR ePrint 2026/994 --- 打开原文核实作者（Dolev/Yung 等）后再引}

\bibitem{zkagent}
L.~Wang, H.~Lou, C.~Li, Y.~Yu, and Y.~Hu, ``zkAgent: Verifiable LLM Agent
Execution via One-Shot Transcript Proofs,'' IACR ePrint 2026/199.

\bibitem{hiddencred}
D.~Balfanz and J.~Schweitzer, ``Hidden credentials,''
\need{核实出处：2002 年 unpublished manuscript / 2003 年会议版本，二者择一并核准页码}.

\bibitem{handshake}
\need{核实作者列表与页码（首发 S&P 2003, pp.~139--150）}D.~Balfanz, G.~Durfee,
N.~Shankar, D.~K.~Smetters, J.~Staddon, and H.-C.~Wong, ``Secret handshakes from
pairing-based key agreements,'' in \emph{IEEE S\&P}, 2003.

\bibitem{xacmlpriv}
\need{核实作者与出处。骨架点名的方向是「XACML 策略的隐私保护访问控制」，
候选：C.~A.~Ardagna 等关于 XACML 隐私的系列工作。\textbf{未核实前不得写成
具体条目}。}

\bibitem{sw05}
\need{核实页码（LNCS 3494, pp.~457--473）}A.~Sahai and B.~Waters, ``Fuzzy
identity-based encryption,'' in \emph{EUROCRYPT}, 2005.

\bibitem{ksw08}
\need{核实页码（LNCS 4965, pp.~515--534）}J.~Katz, A.~Sahai, and B.~Waters,
``Predicate encryption supporting disjunctions, polynomial equations, and inner
products,'' in \emph{EUROCRYPT}, 2008.

\bibitem{vdr}
\need{Verifiable Document Redacting, Springer, DOI 10.1007/978-3-319-66402-6_20
--- 核实作者与年份。\textbf{绝不能再引 IACR ePrint 2024/813}（那是
Kiraz/Larraia/Vaughan, "How to Redact the Bitcoin Backbone Protocol"）。
见 docs/tdsc-fixlist.md \S A。}

\end{thebibliography}

\end{document}
"   % 两遍，第二遍收交叉引用
%  \SKELQUIET 让所有中文标记宏吞掉参数；不带它则需要能排版中文的引擎
%  （xelatex + CJK 字体）。正文本身全英文，标记里才有中文。
%
%  ⚠️ IEEEtran.cls 需要装在 TeX 发行版里。WSL 侧已装进
%     /home/dong/texlive/texmf-local/tex/latex/ieeetran/（2026-10-12）。
% =====================================================================
\documentclass[journal]{IEEEtran}

\usepackage{amsmath,amssymb,amsthm}
\usepackage{booktabs}
\usepackage{array}
\usepackage{graphicx}
\usepackage{xcolor}
\usepackage{hyperref}

% 正文用 Times（ptm，窄），\texttt 用 Courier（pcr，等宽且宽）。
% 一行里只要有一个长 Courier 标识符（\texttt{normalized\_keyword\_block}
% 这类）就几乎没有可拉伸的胶，TeX 只能拉出 badness 10000 的大空隙。
% emergencystretch 允许第三遍在这些行上额外拉伸，把最坏情况吸收掉；
% 代价是那些行略松，远好于一道 10000 的豁口。
% （2026-10-12：装上真正的 Times 后此类行从 3 涨到 12 才暴露。）
\setlength{\emergencystretch}{2em}

% microtype：字形微伸缩（pdftex 下 ±2%）+ 标点悬挂。
% emergencystretch 只能吸收「可拉伸胶不足」，但 badness 在 10000 处封顶，
% 段首被后面长片段（consequences---customer-facing、normalized\_keyword\_block）
% 挤空的行照样报 10000。microtype 从字形宽度上补，是这一类的标准解法。
\usepackage{microtype}

% 窄栏里用 p{} 会因两端对齐产生大片字距（编译时刷一堆 Underfull \hbox）。
% 表格正文列一律用 L{}：左对齐、允许自然参差。
\newcolumntype{L}[1]{>{\raggedright\arraybackslash}p{#1}}
\newcommand{\tabfont}{\footnotesize}

\newtheorem{theorem}{Theorem}
\newtheorem{lemma}[theorem]{Lemma}
\newtheorem{proposition}[theorem]{Proposition}
\newtheorem{corollary}[theorem]{Corollary}
\newtheorem{definition}[theorem]{Definition}
\newtheorem{assumption}[theorem]{Assumption}
\newtheorem{remark}[theorem]{Remark}

% --- 工作标记：提交前全部清掉 ---
% 注意：宏名不能叫 \wp —— \wp 已是 LaTeX/amssymb 的 Weierstrass ℘，
% \newcommand{\wp} 会直接报 "Command \wp already defined" 并使后文级联报错。
\ifdefined\SKELQUIET
  \newcommand{\wpmark}[1]{}
  \newcommand{\need}[1]{}
  \newcommand{\fixme}[1]{}
\else
  \newcommand{\wpmark}[1]{\textcolor{red}{\textsuperscript{[WP#1]}}}
  \newcommand{\need}[1]{\textcolor{blue}{[[ NEED: #1 ]]}}
  \newcommand{\fixme}[1]{\textcolor{orange}{[[ FIX: #1 ]]}}
\fi

\begin{document}

% ---------------------------------------------------------------------
% 标题：决策 D5 —— 保留 "Zero-Knowledge" 的前提是 groth16 成为交付档 (WP1)。
%       若 WP1 未打通 groth16，改用下面被注释掉的那个标题，
%       并同步删掉摘要与 §8 里的 zero-knowledge 措辞。
% ---------------------------------------------------------------------
\title{Proof-of-Policy: Zero-Knowledge Compliance Certificates\\ for LLM Agent Outputs}
% 备选标题（若 groth16 未打通）：
% \title{Proof-of-Policy: Succinct Compliance Certificates for LLM Agent Outputs}

\author{\need{authors}}
\maketitle

% ---------------------------------------------------------------------
\begin{abstract}
% 决策 E2：内容隐私上界必须出现在摘要，不能埋进 §Limitations。
% 决策 D1：贡献陈述按"分析 x 系统各半"写 —— 分析在前，系统是分析所需的那一实例。
Proof-of-Policy (PoP) lets an auditor check that an LLM agent's response satisfies
a published policy, without trusted hardware and without the operator disclosing
the response in the clear. Our contribution is analytical before it is
constructive. We give a quantitative bound on what such a certificate binds---a
guessing-entropy term that no additional computation can improve---together with
the structural assumption the bound requires, and a proof that the assumption is
necessary rather than convenient. The same accounting governs the request side,
where a certificate without receiver-supplied freshness is replayable without
limit. PoP is then presented as the instance these bounds require---a policy
compiler, a
% 决策 D5 条件措辞：若 groth16 未成为交付档，把下行的 "zero-knowledge" 换成 "succinct"，
% 并同步改标题与 §8。
zero-knowledge circuit, and a signed certificate layer---measured end to end
rather than asserted, with the leak profile of the certificate reported as a
result in its own right.

\smallskip
\noindent\textbf{Scope of privacy.} PoP provides \emph{policy}-zero-knowledge:
the certificate reveals neither the policy evaluation trace nor, in
policy-private mode, the policy itself. It does \emph{not} guarantee that $T$
is unrecoverable: as we show in \S\ref{sec:binding}, any publicly recomputable
binding to $T$ admits a re-execution attack whenever the adversary knows the
prompt and the model. The same bit-budget governs the request side: a
certificate is replayable without limit unless the receiver contributes fresh
randomness (Theorem \ref{thm:t3}).
\end{abstract}

\begin{IEEEkeywords}
zero-knowledge proofs, LLM agents, policy compliance, verifiable computation,
accountability
\end{IEEEkeywords}

% =====================================================================
\section{Introduction}
\label{sec:intro}

LLM agents increasingly take actions that carry con\-se\-quences---\allowbreak{}
customer-facing replies, tool invocations, financial and medical content. Regulation such as the
EU AI Act (Articles~12 and~13) and enterprise governance both require that
compliance be auditable. In practice it can only be taken on trust: logs can be
forged, guardrail claims cannot be checked, and content privacy forbids retaining
the plaintext that would settle the question.

Four lines of work each occupy one corner of this problem without spanning it.
\emph{TEE execution proofs}---Proof-of-Guardrail and its relatives---show that
guardrail code was executed and that inputs and outputs were signed; as that work
itself observes, \emph{execution is not compliance}, and the approach carries
roughly $34\%$ latency overhead together with a hardware dependency.
\emph{TEE governance stacks} (Agentrust, cMCP, cA2A) inherit both the hardware
and the vendor facility. \emph{Formal verification} (Lean-Agent and similar)
requires the entire policy to be formalised by hand, and does not apply to
black-box responses. \emph{ZK regex and PII primitives} (zk-regex~V2, verifiable
document redaction \cite{vdr}, ZK Email) are building blocks; none of them is a
policy-level system.

The nearest work, WITNESS, supports two policy primitives, has no peer-reviewed
paper, and---by its own project page---ships an Apache-2.0 reference
implementation, so ``open source'' is not a differentiator; it also offers no
content privacy. \texttt{@brivora/verify} is not zero-knowledge, and DeepProve
proves the \emph{model computation} layer rather than the \emph{policy
compliance} layer. The combination of a general policy DSL compiled to ZK
constraints, \emph{without} a TEE, with selective disclosure of response content
and a formal security model, is an explicit vacancy. This paper occupies it with
PoP---but the system is not the contribution we lead with. Leading with a system
invites the objection that a signature would do; leading with the analysis is
what makes the system necessary, and it is the analysis that this paper is
organised around.

\need{⚠️ 上面点名的 Proof-of-Guardrail / Agentrust,cMCP,cA2A / Lean-Agent /
zk-regex V2 / ZK Email / WITNESS / \texttt{@brivora/verify} / DeepProve，
以及 EU AI Act Art.~12/13，\emph{均未给 \cite{}}：条目尚未逐条核实
（docs/tdsc-fixlist.md \S2、WP8）。核实前不得凭空写 bibitem，故此处只留名。}

\smallskip
\noindent\textbf{What we do \emph{not} claim.} We do not claim that policy proofs
bind nothing: our own system already binds the response through a verifier-issued
nonce (Definition~\ref{def:system}). We do not claim to be the first to observe
that a proof can fail to bind the artifact it purports to describe: the
\emph{effort gap} for ZK inference proofs is formalised by Hollow-LLM
\cite{hollowllm}, and the \emph{source-integrity gap} for predicate proofs is
named by zk-attach \cite{zkattach}. And we do not claim the separation theorem
that our designated-verifier interface makes tempting (``public verification is
impossible when $k$ is small''); Remark~\ref{rem:designated} gives the reason,
and Corollary~\ref{cor:reexec} is the short version of it.

\smallskip
\noindent\textbf{Contributions.}
\begin{enumerate}
  \item We isolate the gap that survives in every design we surveyed---\emph{divergent
  delivery}---and show that neither hardware nor any proof closes it, while having
  the receiver contribute to the proven statement does
  (Proposition~\ref{thm:gap}).
  \item We give a quantitative bound on what a public binding reveals, together
  with the \emph{structural hypothesis} that bound needs, and a demonstration
  that the hypothesis is necessary: our own two public values sit on opposite
  sides of it (Theorem~\ref{thm:t1}, Remark~\ref{rem:betanecessary}). The
  governing quantity $k$ is measurable rather than assumed.
  \item We price the fix. A certificate is replayable without limit unless the
  receiver contributes fresh randomness, and with an empty nonce the adversary's
  advantage is $1$ regardless of every other parameter of the system
  (Theorem~\ref{thm:t3}).
  \item We describe PoP, a system that realises this design without a TEE
  (\S\ref{sec:system}), and we report what its certificates leak about the policy
  they certify (\S\ref{sec:private-policy}).
\end{enumerate}

\smallskip
\noindent\textbf{One spine, two sides.} Theorems~\ref{thm:t1} and~\ref{thm:t3}
state the same fact about two different inputs: \emph{a proof binds only as much
as the committed, adversary-non-selectable quantity has bits for}. The response
side is governed by entropy, $k = -\log_2 \mathsf{Guess}(T \mid Z)$
(Definition~\ref{def:k}); the request side by freshness, $f$. Neither can be
bought with more computation.

\fixme{全文禁止出现"比 Proof-of-Guardrail 更强"或任何强弱对比。
Proof-of-Guardrail 与本文正交：它证"计算发生过"，本文证"性质成立"。}

% =====================================================================
\section{System and Threat Model}
\label{sec:system}

The object of study is a \emph{policy-proof system}: a pipeline that compiles a
policy written in a small DSL into a circuit, evaluates an agent's response
against that circuit inside a zkVM, and emits a certificate binding the verdict
to the response. This section fixes the object and its adversary before any
theorem is stated.

\subsection{The system}
\label{sec:system-overview}

\paragraph{Pipeline and the cross-layer contract.} A policy package is compiled
to a \texttt{ConstraintSpec}---a JSON artefact carrying the compiled NFA---which
is the \emph{only} object crossing the layer boundary: the Python reference
implementation and the Rust/zkVM guest consume the same bytes, so cross-layer
drift is structurally impossible rather than merely tested against. The guest
runs a deterministic evaluator and commits its outcome. Of eight rule classes,
seven are decided \emph{inside} the circuit
(\texttt{keyword\_\allowbreak block},
\texttt{normalized\_\allowbreak keyword\_\allowbreak block},
\texttt{length\_\allowbreak bound}, \texttt{pattern\_\allowbreak block},
\texttt{format\_\allowbreak check}, \texttt{tool\_\allowbreak arg\_\allowbreak guard},
\texttt{budget\_\allowbreak bound}). The eighth,
\texttt{semantic\_bound}---a rule of the form ``harm probability $\le$
threshold''---is precisely the kind for which no formal proof can be learned, and
the circuit does not decide it. The circuit registers only \emph{that the rule
was delegated}, in the public value \texttt{delegated}, and a verifier must
\emph{conjunctively} check a companion \texttt{ezkl}/halo2 proof. Verifying the
zkVM side alone does not establish compliance, and we say so here because it is
the most likely misreading of the certificate.

\paragraph{What the verdict is bound to.} Rules that reason about the agent's
\emph{tool calls} cannot be judged against a list the prover fills in itself.
Receipts are therefore issued by a tool gateway after each call and chained
(\texttt{seq}, \texttt{prev}, signature); the circuit checks the chain's
\emph{structure}---entries consecutive, mutually consistent, each digest
recomputed from its own content---and fails closed to \texttt{trace\_unbound}
otherwise. Signature verification happens \emph{off circuit} (Ed25519 inside the
zkVM is prohibitively expensive), a boundary we state rather than gloss: the
circuit is blind to a forgery in the chain's final element, and two independent
off-circuit gates---digest comparison and signature verification---carry that
obligation, neither substituting for the other. Truncating the chain's tail is a
separate and weaker opening, closed by a session-end \texttt{ToolSeal}
commitment from the gateway, and it remains conditional on the verifier holding
the gateway's public key. \S\ref{sec:discussion} develops these residual
boundaries; they are not incidental to the design.

\paragraph{Dual privacy modes.} In \emph{transparent} mode the response and the
verdict are public. In \emph{private} mode the certificate publishes a response
commitment, a commitment per violation, and a redaction proof
(\texttt{redaction\_ok} together with \texttt{mask\_covered}, the latter
requiring the masked positions to fall inside a span the circuit verified to be a
genuine complete match). The guarantee is that \emph{public values contain no
plaintext}---not that $T$ cannot be recovered. A commitment is a public function
of $T$ against which anyone can guess and check, while the response binding must
be publicly recomputable for a verifier to check the string it was delivered.
Those two requirements are in direct tension, and \S\ref{sec:binding} quantifies
exactly when.

\paragraph{Certificates and anchoring.} A certificate is signed with a DSSE
envelope under Ed25519 with the key held by the issuing party, and its digest is
appended to a hash-chain anchor ledger. Two backends share one interface: a file
ledger, verifiable offline, and an on-chain registration storing only a 32-byte
digest. The transaction hash and on-chain timestamp are written back into the
local entry, so the offline chain and the public timestamp cross-check each
other and tampering with either is detectable.

\paragraph{Framework integration.} One gateway serves LangChain callbacks,
LangGraph (\texttt{guard\_node} plus full event streams), and MCP tool argument
and response checks, so trace certificates and proofs share a single chain of
evidence. The implementation passes 801 tests with 15 skipped; the skipped set is
enumerated in Appendix~\ref{app:repro}.

\begin{definition}[Policy-proof system]
\label{def:system}
A policy-proof system is a tuple $(\mathsf{Prove}, \mathsf{Ver}, \pi, \mathcal{M})$.
On challenge $n$ from a verifier, $\mathsf{Prove}(\pi, T, n)$ outputs a
certificate $c$ whose public values include
\begin{multline*}
b \;=\; H(\texttt{pop-bind-v1} \,\|\, \mathsf{u32be}(|n|) \,\|\, n \,\|\, T),\\
C = H(T).
\end{multline*}
$\mathsf{Ver}(c, n, T')$ accepts iff the policy relation holds \emph{and}
$b$ recomputed from the received $T'$ matches.
\end{definition}

\fixme{\textbf{H3/H4}：早先写的 ``the relation contains no term referring to the
delivered bytes'' 与仓库事实冲突——公开值里\emph{已有} $b$，验证时由送达的
$T'$ 现场重算（\texttt{circuits/types/src/lib.rs:444}、
\texttt{docs/security-model.md:232}）。命题已按上面的定义改写。见
\texttt{docs/tdsc-fixlist.md} \S H。}

\subsection{Threat model}
\label{sec:threat}

The adversary is the \emph{agent operator}: the party that holds the policy, runs
the model, and produces the certificate. The verifier is a regulator, auditor, or
third party holding only public values. The policy author issues $\pi$; in this
section we treat it as honest, and \S\ref{sec:wp4-tension} makes that assumption
explicit and scopes it, because in policy-private mode its honesty is doing real
work and cannot be waved through. We state the goals, and more importantly the
non-goals, before the theorems.

\begin{assumption}[Adversary side information]
\label{as:Z}
The adversary knows the prompt, the model (possibly its weights), and the
decoding parameters. \emph{This is the assumption that makes the theorem
non-trivial}: under greedy decoding, $k \approx 0$.
\end{assumption}

\paragraph{Goals.}
\begin{description}
  \item[(G1)] \emph{Soundness}---a response $T$ violating $\pi$ cannot yield an
  accepted certificate.
  \item[(G2)] \emph{Committed privacy}---in private mode the public values
  contain no plaintext of $T$. This is explicitly \emph{not} the claim that $T$
  is unrecoverable: once the guessing entropy $k$ of Definition~\ref{def:k} falls
  below $\log q + \log(1/\varepsilon)$, a publicly recomputable binding and
  content hiding are mutually exclusive (Theorem~\ref{thm:t1}).
  \item[(G3)] \emph{Traceable binding}---policy version, circuit, and proof
  artefact are bound to one another and recomputed at verification.
  \item[(G4)] \emph{Selective disclosure}---redaction and evidence opening are
  themselves provable.
  \item[(G5)] \emph{Record integrity}---tampering with, or reordering, the audit
  ledger is detectable.
\end{description}

\paragraph{Non-goals.} We do not prove the model's inference, and we do not cover
training data or model cards. (G1) is modulo the soundness of the zkVM and of the
hash function; we use those primitives and do not re-derive them.

\paragraph{Where the policy-private case is developed.} The three-party model,
the issuing prerequisite (P-PP), and the table of what each verifier may and may
not conclude are specific to policy-private mode, where the policy itself is
secret. They are stated once, in \S\ref{sec:wp4-tension}, and are deliberately
\emph{not} duplicated here---at a twelve-page limit they cannot be. The general
adversary of this section is the same operator, viewed without the
policy-secrecy refinement.

% =====================================================================
\section{Binding of Policy Proofs}
\label{sec:binding}

A policy proof is meant to bind a response: a verifier who accepts a
certificate should conclude that \emph{this} string was checked against the
policy. This section asks what that binding is actually worth. We first isolate
the gap that survives in every design we surveyed (\S\ref{sec:gap}), then give
a quantitative bound on what a public value reveals, together with the
structural hypothesis that bound needs (\S\ref{sec:tradeoff}), and finally
state the same fact on the request side, where the governing resource is
freshness rather than entropy (\S\ref{sec:t3}).

\subsection{The gap}
\label{sec:gap}

\begin{proposition}[Divergent delivery]
\label{thm:gap}
Definition \ref{def:system} binds $T'$ \emph{as received by the verifying party}. If the
provider delivers a different string $T'' \ne T'$ to the end user, verification
still succeeds on $T'$. Neither hardware nor a proof can prevent this; only the
end user holding $n$ and recomputing $b$ can.
\end{proposition}
\begin{proof}
$b$ is a function of $(n,T)$ only; the delivery channel to the end user carries
no authenticated term. One line.
\end{proof}

Two remarks scope this proposition. First, it describes a \emph{residual} gap,
not the claim that policy proofs bind nothing. Our own system already binds $T$
through the verifier-issued nonce $n$ of Definition \ref{def:system}, and a
verifier who receives $T'$ and recomputes $b$ does learn that $T'$ was the
proved string; the gap is that the binding is anchored to the
\emph{verifying party's} copy. Second, we do not claim to be the first to
observe that a proof can fail to bind the artifact it purports to describe. The
\emph{effort gap} for zero-knowledge inference proofs is formalized by
Hollow-LLM \cite{hollowllm}, and the \emph{source-integrity gap} for predicate
proofs is named by zk-attach \cite{zkattach}; both concern what the proof system
was asked to prove, whereas Proposition \ref{thm:gap} concerns where the proved
string is delivered. We state the naive version of our own claim because it
would be false against our own system.

\subsection{A quantitative binding--hiding tradeoff}
\label{sec:tradeoff}

The certificate must publish \emph{something} tying it to $T$; otherwise it
binds nothing at all. We therefore ask how much that public value leaks. The
answer depends on one measurable quantity --- the guessing entropy of the
response given everything the adversary already knows --- and on one structural
property of the public value itself.

\begin{definition}[Guessing entropy of a response]
\label{def:k}
For a response distribution $\mathcal{D}$ and adversary side information $Z$,
let
\[
\mathsf{Guess}(T \mid Z) \;\triangleq\; \max_{\mathcal{A}} \Pr[\mathcal{A}(Z) = T],
\]
which equals $\mathbb{E}_z[\max_t \Pr[T = t \mid Z = z]]$, and
\[
k \;\triangleq\; -\log_2 \mathsf{Guess}(T \mid Z),
\]
the maximum ranging over all time-unbounded adversaries.
\end{definition}

\begin{remark}[The definition is operational, and it is the average]
\label{rem:kmeas}
$k$ is \emph{measured}, not assumed: sampling one model on one prompt many
times estimates $\mathsf{Guess}(T \mid Z)$ directly
(\S\ref{sec:eval}), so Theorem \ref{thm:t1} is instantiated with a
number we obtain rather than a number we posit. The definition must average
over $z$ rather than take the pointwise worst case
$k_z = -\log \max_t \Pr[T = t \mid Z = z]$: the proof takes an expectation over
$z$, and $\mathbb{E}_z[2^{-k_z}] = 2^{-k}$ holds only for the average form.
\end{remark}

\begin{theorem}[Binding, and the hypothesis it needs]
\label{thm:t1}
Let $\beta : \mathcal{M} \to \{0,1\}^\lambda$ be public, deterministic and
efficiently computable, and let the certificate publish $\beta(T)$. Let
$\mathcal{A}$ be any adversary, unbounded in time, that is given $Z$ and
evaluates $\beta$ at most $q$ times.
\begin{enumerate}
\item[(H)] \emph{(Randomised public value.)} Suppose
  $\beta(T) = H(\mathsf{dom} \,\|\, \mathsf{u32be}(|n|) \,\|\, n \,\|\, T)$ for a
  random oracle $H$ and a nonce $n$ drawn uniformly and \emph{independently of
  $T$}. Then
  \[
  \Pr\!\left[\mathcal{A}^\beta\!\left(Z, \beta(T)\right) = T\right]
  \;\le\; q \cdot 2^{-k} \;+\; q \cdot 2^{-\lambda}.
  \]
  Consequently hiding with advantage $\varepsilon$ requires
  $k \ge \log q + \log(1/\varepsilon)$.
\item[(L)] \emph{(Lossy public value.)} For an arbitrary deterministic public
  $\beta$ the correct bound is over the \emph{preimage class}:
  \begin{multline*}
  \Pr\!\left[\mathcal{A}^\beta\!\left(Z, \beta(T)\right) = T\right]
  \;\le\; \mathsf{Guess}\!\left(T \mid Z, \beta(T)\right) \\
  \triangleq\; \mathbb{E}_y\!\left[\max_t \Pr[T = t \mid Z, \beta(T) = y]\right].
  \end{multline*}
\end{enumerate}
\end{theorem}

\begin{proof}[Proof sketch]
Part (L) is immediate after conditioning on $\beta(T)$; no hypothesis on
$\beta$ is used. For (H), fix $z$ and decompose the adversary's success by the
first query that hits $T$. The oracle answers it has already seen are uniform
and independent of $T$, so each query is itself independent of $T$ and succeeds
with probability at most $2^{-k_z}$. Taking the expectation over $z$ and
applying Definition \ref{def:k} yields $q \cdot 2^{-k}$; the term
$q \cdot 2^{-\lambda}$ charges a query that hits the uniform value $\beta(T)$
by chance. Appendix~\ref{app:proofs} gives the full argument.
\end{proof}

\fixme{$\lambda$ 比特随机与随机预言机只给 \emph{heuristic} 保证：实践中的 SHA-256
不是真的随机预言机（骨架 :712--715）。正文措辞必须写成 heuristic，
\textbf{不得}写成信息论界。}

\begin{remark}[The hypothesis in (H) is necessary, and our two instances straddle it]
\label{rem:betanecessary}
Without the hypothesis, part (H) is not merely unproven but \emph{false}: take
$\beta = \mathrm{id}$, or $\beta = \mathsf{encode}$ as in
Proposition \ref{prop:semantic}. The adversary outputs $\beta(T)$ and succeeds
with probability $1$, independently of $q$ and $k$. Our own two public values
therefore sit on opposite sides of the hypothesis---$b =
H(\mathsf{dom}\|n\|T)$ satisfies it, $\mathsf{encode}(T)$ does not---and this is
what makes the pair a theorem about \emph{our design} rather than a citation of
a standard bound. An earlier draft of this work stated (H) without its
hypothesis; the counterexample above is why the hypothesis is now explicit.
\end{remark}

\begin{corollary}[Re-execution, not enumeration]
\label{cor:reexec}
Under Assumption \ref{as:Z}, when the decoding is greedy or low-temperature,
$k \approx 0$ and the attack of Theorem \ref{thm:t1} is realised by
\emph{re-running the model}, at the cost of one inference---independent of
$|T|$ and requiring \emph{no} evaluation of $\beta$. Enumeration over the string
space is \emph{not} the relevant attack.
\end{corollary}

\fixme{原 docs/sp1-zk-audit.md §4 写的是"熵(~10^2--10^3 bit)远低于 256 bit，
因此枚举可行"——这是非因果推论（2^100 与 2^256 一样枚举不了；真正的攻击是\textbf{重执行}，
对手的已知信息是 (prompt, 模型, 解码参数)）。
2026-10-12 已全库改正，docs/tdsc-fixlist.md §G 的 G1--G6 逐条结案。}

\begin{remark}[A design observation, not a separation]
\label{rem:designated}
Our verification interface is \emph{designated}: binding is checked against a
value $b = H(\mathsf{dom}\|n\|T)$ in which a verifier-issued nonce $n$ appears.
It is tempting to promote this into a separation (``public verification is
impossible when $k$ is small''). \textbf{We claim no such theorem, and
Corollary \ref{cor:reexec} is why.} When $k \approx 0$, hiding fails for
\emph{every} scheme, designated or not, because the adversary re-runs the
model; when $k$ is large, a zero-knowledge proof of ``$\exists\, T : b =
H(\mathsf{dom}\|n\|T) \wedge J(\pi, T) = 1$'' is simultaneously publicly
verifiable and hiding. The dichotomy is thus governed by $k$ alone and carries
no independent content. What does survive---and what \S\ref{sec:discussion}
develops---is that \emph{third-party auditing of delivery} forces disclosure of
$T$, which is an informational remark rather than a theorem. We record the
demotion here because an earlier draft asserted the separation as a theorem.
\end{remark}

We stress that the designated-verifier shape is a \emph{retrospective
justification, not a novel construction}. The system publishes
$b = H(\mathsf{dom}\|n\|T)$ with a verifier-issued nonce $n$, and
Remark \ref{rem:designated} explains the shape of that choice after the fact.
The security model does not pin down whether $n$ is public: making it public
forfeits hiding, while holding it at the verifier is what makes verification
designated.

\begin{remark}[A hypothesis that can be settled by experiment]
\label{rem:kexp}
Theorem \ref{thm:t1} is only as useful as the number put into it. Because
Definition \ref{def:k} is operational, $k$ can be measured rather than assumed:
sampling one model on one prompt many times estimates
$\mathsf{Guess}(T \mid Z)$ directly, and the estimate substitutes into the
threshold $k \ge \log q + \log(1/\varepsilon)$. To our knowledge no prior work
measures this quantity for agent responses, which makes the measurement a
contribution in its own right; \S\ref{sec:eval} describes the protocol and
reports the result.
\end{remark}

\begin{proposition}[The lossy instance of Theorem \ref{thm:t1}(L)]
\label{prop:semantic}
For the semantic rule class, the public values contain $\beta(T) =
\mathsf{encode}(T)$, a content-preserving feature vector of the response. This
$\beta$ is deterministic and injective on the features but not on strings, so
it violates the hypothesis of Theorem \ref{thm:t1}(H) and instantiates part
(L): the binding force is exactly
$\mathsf{Guess}(T \mid Z, \mathsf{encode}(T))$, i.e.\ the guessing probability
\emph{within the preimage class} of the published vector. Here $k$ need not be
argued from perplexity at all: the adversary's search space is that preimage,
which is dramatically smaller than the string space. The relevant asymmetry is
not ``low entropy'' but a \emph{lossy public encoding}.
\end{proposition}
\begin{proof}
Immediate. The measurement is direct: publish $\mathsf{encode}(T)$ for real
responses and measure how much of $T$ is recoverable.
\end{proof}

Proposition \ref{prop:semantic} is the cheapest experiment in this paper and
the strongest anchor for Theorem \ref{thm:t1}: it is measurable in a day, and it
makes the bound concrete rather than cited. It also shows that the relevant
asymmetry is not entropy at all --- for the semantic rule class there is no
entropy argument available, because $\mathsf{encode}(T)$ simply \emph{is} a
content-preserving function of the response.

\subsection{Quantitative request binding}
\label{sec:t3}

\begin{theorem}[Selective-presentation advantage]
\label{thm:t3}
Fix a policy $\pi$ and a request distribution. A selective-presentation
adversary $\mathcal{B}$ may, \emph{before} seeing the receiver's challenge,
compute a set $S$ of accepted transcripts. The receiver then contributes $f$
fresh, uniformly random bits $r$ that enter the proven statement. Then
\[
\Pr\!\left[\mathcal{B}\ \text{is accepted on } r\right] \;\le\; |S| \cdot 2^{-f}.
\]
If instead the receiver contributes no fresh randomness ($f = 0$; e.g.\ an
empty nonce), the advantage is $1$, and \emph{does not decay in any parameter
of the system}.
\end{theorem}
\begin{proof}
The receiver's freshness $r$ enters the proven statement through
$b = H(\mathsf{dom} \,\|\, r \,\|\, T)$. For a certificate in $S$ to be accepted
on $r$, the value $r$ must fall inside the challenge set that certificate
already covers --- an event of probability at most $2^{-f}$ for $f$ fresh
uniform bits --- and a union bound over $S$ gives $|S| \cdot 2^{-f}$. When
$f = 0$ every element of $S$ is a permanently valid replay, so the advantage is
$1$ and does not decay in any parameter of the system.

The $f = 0$ branch is not hypothetical.
\texttt{circuits/\allowbreak types/\allowbreak src/\allowbreak lib.rs}
(lines 950--995) documents an empty nonce as a legal input that ``provides no
replay protection''; it is therefore a \emph{valid input to the current API},
not a straw man.
\end{proof}

Theorem \ref{thm:t3} is \emph{isomorphic} to Theorem \ref{thm:t1}, and that is
what earns it a place here. Both state the same fact about different inputs: a
proof binds only as much as the committed quantity the adversary cannot select
in advance has bits for. On the response side that quantity is the response and
the resource is entropy $k$; on the request side it is freshness $f$. Neither
can be bought with more computation. This is the spine announced in
\S\ref{sec:intro}, and the two theorems are its two sides.

\begin{remark}[A taxonomy, not a trichotomy]
\label{rem:taxonomy}
Every policy-proof deployment we surveyed closes the divergent-delivery gap of
Proposition \ref{thm:gap} in one of three ways:
(a) root trust in hardware (attested enclaves);
(b) have the \emph{end user} contribute a secret to the proven statement (the
user issues the nonce); or
(c) gate delivery on successful verification (\emph{prove-before-deliver}).
\textbf{We offer this as a classification, not as an exhaustiveness theorem}:
we neither prove that a fourth family cannot exist, nor that the three are
disjoint. zk-attach \cite{zkattach} walks (a); NiyamAI and ZKAP walk (c); PoP
walks (b).
\end{remark}

\paragraph{Posture, stated plainly.} Our system already carries
\texttt{response\_binding}, but $n$ is issued by the \emph{verifier}. Moving it
to the \emph{end user}, so that $f > 0$ becomes the receiver's contribution,
closes the divergent-delivery branch of Proposition \ref{thm:gap}. We are
explicit that this is a small protocol amendment rather than a new system: the
mechanism is already implemented, and only the issuer of the nonce changes. What
Theorem \ref{thm:t3} buys is the \emph{price}: with $f = 0$ there is no closure
at any cost.

% =====================================================================
\section{Policy-Private Compliance}
\label{sec:private-policy}

For a policy whose rules are deterministic and a response that is public
anyway, a signed hash chain over the response suffices: the signature already
binds the bytes, and anyone can re-evaluate the rules. The one scenario a
signature cannot cover is a policy that must itself stay secret---where the
operator has to certify compliance to a third party without handing over the
rules it is being judged against. That is the mode this section is about, and we
lead with what the mode \emph{cannot} buy, because that determines whether it
means anything at all.

\subsection{The observation that reframes this work package}
\label{sec:wp4-finding}

\texttt{spec\_canonical} is already a \emph{guest input}---a witness---rather
than a public value; only its hash,
\[
\mathsf{policy\_hash} \;=\; \mathsf{SHA256}(\texttt{spec\_canonical}),
\]
enters the public values (\texttt{circuits/\allowbreak types/\allowbreak
src/\allowbreak lib.rs}:1--16, 415--416, 918, 1003--1004). The instruction
``do not reveal the policy text'' is therefore
\emph{already true of the architecture}; what leaks is everything \emph{derived}
from the policy. Policy privacy is thus not a new proof system to be built but a
\emph{disclosure policy to be audited}, and we write it that way---as a
systematic account of what a certificate reveals about the policy it certifies.

\begin{table}[t]
\centering
\caption{What a policy-private certificate leaks today, and what closing each
leak costs.}
\label{tab:wp4}
\tabfont
\begin{tabular}{@{}L{0.04\linewidth}L{0.46\linewidth}L{0.36\linewidth}@{}}
\toprule
\# & Leak & Difficulty / dependency \\
\midrule
1 & The policy text itself & \textbf{Already closed} (it is a witness). \\
2 & Witness hiding in the proof system & Medium; \emph{inherited from the
    proving tier}: only \texttt{groth16}/\texttt{plonk} is witness-hiding, per
    \texttt{docs/\allowbreak sp1-zk-audit.md}. \\
3 & \texttt{PrivateViolation}\allowbreak\{\texttt{rule},\allowbreak\texttt{kind}\}
    names the rule that fired & Medium; needs a coarsened violation class plus a
    commitment to it. \\
4 & Evidence opening reveals the rule
    (\texttt{policydsl/\allowbreak privacy/\allowbreak commit.py}:319) & Medium;
    disclose the \emph{existence} of a violation, not its identity. \\
5 & The \texttt{delegated} flag exposes policy structure (which rules were
    routed to the \texttt{ezkl} path) & Medium; easiest to overlook, cheapest to
    fix. \\
6 & Semantic rules publish $\beta(T)=\mathsf{encode}(T)$
    (Proposition~\ref{prop:semantic}) & Medium--high; either bound the
    preimage, or disable this rule class in policy-private mode. \\
\bottomrule
\end{tabular}
\end{table}

The honest structure of Table~\ref{tab:wp4} is worth stating, because it sets
the reader's expectations for the rest of the section. Sub-problem~1 is a
one-paragraph re-reading of existing code, not a contribution. Sub-problem~2 is
inherited wholesale from the proving tier and is not specific to this work.
Sub-problems~3--5 are engineering with a single formal core---the granularity of
disclosure. Sub-problem~6 is where the work actually is, and it collides with
\S\ref{sec:tradeoff}: one cannot hide the policy while publishing the lossy
encoding of the response that the semantic rule class requires
(Proposition~\ref{prop:semantic}).

\subsection{The tension worth writing up: policy privacy vs.\ accountability}
\label{sec:wp4-tension}

\paragraph{The prerequisite, stated as a prerequisite.} Let
$H = \mathsf{SHA256}(\pi)$. On its own $H$ is \emph{inert}. A prover can always
emit a self-consistent certificate for any policy $\pi'$ it prefers, simply by
placing $H' = \mathsf{SHA256}(\pi')$ in that certificate. \emph{Secrecy of
$\pi$ is what creates the gap}, and no cryptographic mechanism can close it:
closing it would require showing $\pi$ to the verifier, which is precisely what
the mode exists to avoid. What closes it is an \emph{issuing act}, so we state
it as an explicit prerequisite of the mode.

\begin{description}
  \item[(P-PP)] An \emph{issuing authority} signs $\pi$ and registers $H$ in a
  record readable by third parties. A policy is \emph{accountable} iff its
  \texttt{policy\_hash} can be traced to such a record.
\end{description}

Without (P-PP) the mode is not merely weaker---the verifier's conclusion is
empty. This is the same assumption as the ``honest policy author'' of
\texttt{docs/security-model.md}~\S0.2; what is new is scoping it to the
verification interface, which had not been done, and which a referee would
otherwise use to dismiss the whole mode in one sentence.

\paragraph{Three parties.} ``Three'' means \emph{authority / operator / third
party}; an auditor is a named instance of the third party, not a fourth trust
anchor.

\begin{table}[t]
\centering
\caption{The three-party model of policy-private mode.}
\label{tab:wp4-parties}
\tabfont
\setlength{\tabcolsep}{3pt}
\begin{tabular}{@{}L{0.16\columnwidth}L{0.14\columnwidth}L{0.62\columnwidth}@{}}
\toprule
Party & Honest? & Holds; wants \\
\midrule
Authority $R$ & \textbf{Yes}, by (P-PP) & Holds $\pi$ and the issuing record;
wants $\pi$ actually enforced, and wants to be able to check that it was. \\
Operator $O$ & \textbf{No} (the adversary) & Holds $\pi$, $T$, the proving key;
wants $\pi$ hidden from third parties while still being able to certify. \\
Third party $V$ & Yes & Holds only public values---\texttt{policy\_\allowbreak hash},
\texttt{circuit\_\allowbreak hash}, \texttt{response\_\allowbreak commitment},
\texttt{passed}, the proof; wants to confirm that \emph{some issued} policy
passed, and must \emph{not} learn $\pi$. \\
\emph{(Auditor $V^{*}$)} & Yes & As $V$, plus $\pi$ once $O$ discloses it;
wants ex-post accountability. \\
\bottomrule
\end{tabular}
\end{table}

Two boundaries must be stated with the model. First, secrecy of $\pi$ comes from
a \emph{data-flow fact}---$\pi$ never enters the public values---and not from
zero-knowledge; any rule that writes a fragment of $\pi$ into the public values
(the leaks in Table~\ref{tab:wp4}, items~3 and~5) breaks secrecy on that
fragment. Second, accountability for $V^{*}$ is \emph{voluntary}: $V^{*}$ can
check only when $O$ chooses to disclose $\pi$. The mode offers no forced
disclosure; that would require the release-gating family discussed in
\texttt{docs/tdsc-plan.md}~\S1.6(b), which we do not claim.

\paragraph{What each verifier may conclude.} Table~\ref{tab:wp4-capability} is
the contribution of this subsection: the cryptography is downstream of it, and
the table survives even if the cryptography only partially lands.

\begin{table*}[t]
\centering
\caption{What each verifier may and may not conclude in policy-private mode.}
\label{tab:wp4-capability}
\tabfont
\setlength{\tabcolsep}{4pt}
\begin{tabular}{@{}L{0.13\textwidth}L{0.30\textwidth}L{0.30\textwidth}L{0.19\textwidth}@{}}
\toprule
Verifier & \textbf{May} conclude & \textbf{May not} conclude & Rests on \\
\midrule
Third party $V$ & (i) the commitment binds $T$
(Definition~\ref{def:system}); (ii) \emph{some} policy with hash $H$ passed;
(iii) $H$ appears in the issuing record & (iv) the \emph{content} of $\pi$;
(v) how many rules $\pi$ has, or of which classes; (vi) ``the operator enforces
the regulation \emph{I} care about''---only ``some issued regulation'' &
(P-PP); Assumption~\ref{as:Z}; the $k$ of Theorem~\ref{thm:t1} \\
End user (also receiver) & As $V$, and additionally: the proved $T$ is the
string it received, since its nonce is bound (Theorem~\ref{thm:t3}, requires
$f>0$) & Only the string it received; \emph{not} what was delivered to another
downstream agent & $f>0$ \\
Authority $R$ & Everything; it holds $\pi$ and can recompute $H$ & --- & --- \\
Auditor $V^{*}$ (after disclosure) & (vii) $\mathsf{SHA256}(\pi) =
\texttt{policy\_\allowbreak hash}$, i.e.\ what ran was the issued $\pi$;
(viii) the disclosed violation record is consistent & (ix) the undisclosed part
of $T$ & voluntary disclosure \\
On-chain verifier & (i)--(iii) only; it hashes and verifies signatures, and does
not interpret policy semantics & (iv)--(ix) & (P-PP) \\
\bottomrule
\end{tabular}
\end{table*}

The $V$ row is the reason this subsection exists. A third party does not obtain
``compliance''; it obtains ``some issued policy was satisfied''. Conflating the
two is the sentence a referee would use to reject the mode, and preventing that
sentence is what the table is for.

\paragraph{Prior art.} \emph{Policy hiding is not a new primitive}, and we claim
no primacy for it. Hidden credentials \cite{hiddencred}, secret handshakes
\cite{handshake}, privacy-preserving access control over XACML policies
\cite{xacmlpriv}, and attribute-hiding and fuzzy predicate encryption
\cite{sw05,ksw08} are one family, and policy-private mode is an instance of it.
What is new here is narrow: (a) the \emph{application} to LLM output compliance,
together with selective disclosure of violation type and location, and (b) the
accountability tension of this subsection.

\need{⚠️ 骨架 §4.3「Expected results, stated in advance」（子问题置信度百分比表）
\emph{未}迁入正文。那是 WP 规划的产物，不是论文内容——期刊稿里出现
「我们对子问题 4 有 70\% 把握」是硬伤。若作者要保留这份校准，应改写为
\S\ref{sec:discussion} 里一段关于适用范围与未完成项的讨论，而不是概率表。}

\need{(P-PP) 的落地形式尚未定：链上已签发策略哈希注册表，还是签名清单？
属作者决策。先例引文已按骨架 :623--630 的要求入 bibliography，
但 \cite{hiddencred} / \cite{handshake} / \cite{xacmlpriv} / \cite{sw05} /
\cite{ksw08} 五条\emph{均待逐条打开原文核实}（外部资源阻塞）。}

% =====================================================================
\section{Evaluation}
\label{sec:eval}

\subsection{Cycle count: length versus rule count}
\label{sec:eval-cycles}

Response length was pushed from 200 to 100{,}000 characters and the rule count
from 1 to 6, giving 40 measurement points (\texttt{bench/bench\_cycles.py}; the
full matrix is in \texttt{bench/results/cycles.md}). The response text comes from
an evaluation artefact the repository's own end-to-end demo produced (372
characters, three segments deduplicated by content); longer points are tiled from
it. What is measured is therefore the \emph{curve} of scanning cost, not ``a real
100{,}000-character document''. Every point was first confirmed to match no rule:
a matching rule short-circuits, and the cost of a short-circuit says nothing.

\begin{table}[t]
\centering
\caption{Cycle count against length, per rule set. Five rule sets, six length
points each; $R^2 \approx 1.0000$ throughout.}
\label{tab:cycles}
\tabfont
\begin{tabular}{@{}rrrrr@{}}
\toprule
Rules & \texttt{keyword} & \texttt{pattern} & Slope (cyc/char) & Intercept \\
\midrule
1 & 0 & 0 & 73.9      & 50{,}137 \\
2 & 1 & 0 & 105.3     & 65{,}942 \\
3 & 1 & 1 & 4{,}289.8 & 387{,}501 \\
4 & 2 & 1 & 4{,}048.8 & 405{,}785 \\
6 & 3 & 2 & 5{,}403.5 & 993{,}257 \\
\bottomrule
\end{tabular}
\end{table}

Three findings, in decreasing order of how much they change what may be claimed.

\emph{(1) With the rule set fixed, cycle count is exactly linear in length.}
Five rule sets, six length points each, $R^2 \approx 1.0000$ throughout
(Table~\ref{tab:cycles}).

\emph{(2) Scanning cost is dominated by the NFA, not by the rule count.} One
keyword rule costs about $+31.6$ cycles per character at the margin; one NFA
regular expression (an email pattern) costs about $+3{,}634$---roughly
$115\times$ more. Equivalently, a rule set containing a regular expression has a
per-character slope ($4{,}289.8$) $58\times$ that of a rule set with only string
rules ($73.9$).

\emph{(3) Unit prices are not additive: a unit price is a property of a rule
\emph{in context}, not of the rule.} Fitting the slopes as
$(\text{keyword},\text{pattern})$ unit prices gives $R^2 = 0.9169$ and a
\emph{negative} keyword unit price ($-477.2$ cycles per character). This is not
noise; two probes give the mechanism directly:

\begin{itemize}
  \item \textbf{Ordering probe} (same rules, same text, $L = 10^5$): declared
  order costs $541{,}342{,}069$ cycles against $570{,}736{,}989$ in reverse,
  a difference of $+5.43\%$, while \texttt{passed} is identical in both. \emph{The
  semantics of the decision do not depend on order; the cost does.}
  \item \textbf{Marginal probe} ($L = 10^5$, filling the one matrix cell the
  sweep could not reach---regular expressions with no keyword rule):
  $[\text{len},\text{pat}]$ over $[\text{len}]$ has a marginal unit price of
  $3{,}633.7$, whereas ``add one more keyword rule'' has two different answers:
  $+31.6$ alone, $+585.6$ in a policy that already contains a regular expression
  ($-153.7$ if inserted before it). The same action changes sign between
  $-154$ and $+586$.
\end{itemize}

The mechanism lies in zkVM memory allocation: every \texttt{ascii\_lower} and
every Pike VM match allocates, and allocator state depends on previous
allocations. We therefore quote \emph{per-rule-set slopes} (the table above,
directly citable) rather than any ``unit price $\times$ count'' expression. The
simple model $\text{cycles} \approx a\cdot|T|\cdot\text{rules} + b$ that this
work anticipated is refuted by the data ($R^2 = 0.872$, negative intercept), and
we report that too: \emph{a model refuted by its own data is itself an evaluation
result.}

\subsection{Matcher ablation: quadratic degradation on a pathological input}
\label{sec:eval-ablation}

\emph{Construction.} Pattern \texttt{a+b} against input
\texttt{a}$^{\times n}$ containing no \texttt{b}, so the match is necessarily
absent and both implementations must scan to the end. The naive implementation
reinvents and restarts at every start position, giving $\Theta(n^2)$; the Pike VM
decides in a single linear pass (\texttt{bench/bench\_ablation.py}).

\begin{table}[t]
\centering
\caption{Matcher ablation. Both paths decide identically; only the cost differs.}
\label{tab:ablation}
\tabfont
\begin{tabular}{@{}rrrrr@{}}
\toprule
$n$ & Pike (cycles) & Naive (cycles) & Ratio & Naive growth/char \\
\midrule
100 & 643{,}894   & 12{,}382{,}720  & 19.2$\times$ & --- \\
200 & 1{,}145{,}209 & 48{,}638{,}946 & 42.5$\times$ & $+96.4\%$ \\
400 & 2{,}143{,}617 & 193{,}207{,}130 & 90.1$\times$ & $+98.6\%$ \\
\bottomrule
\end{tabular}
\end{table}

Read per-character cost rather than absolute figures. Doubling $n$ leaves the
Pike VM's per-character cost essentially unchanged ($-11.1\%$, $-6.4\%$) while
the naive path's nearly doubles ($+96.4\%$, $+98.6\%$). At $n = 400$ the naive
path is already $90\times$ slower, and that ratio doubles again with each
doubling of $n$. \emph{This is an attack surface}: the input length for policy
compliance is chosen externally---an agent's reply, a concatenated session---so
an adversary who needs only to write a long reply drives a naive implementation's
cost quadratically. The comparison is meaningful only because the two paths
decide identically, which \texttt{tests/test\_ablation.py} pins down (Python and
Rust \texttt{match\_search} $\equiv$ \texttt{match\_search\_naive}). To our
knowledge this is the first ablation evidence specific to policy-compliance
matching.

\subsection{Real proofs: time, size, memory---and the ceiling}
\label{sec:eval-proofs}

Sampling points share the corpus of \S\ref{sec:eval-cycles}. Each point runs in
its own subprocess (SP1's memory accumulates within a process, so two proofs in
one process do not price one proof), and failures are reported in the table
rather than omitted. Rows marked OOM are a \emph{result}, not a malfunction: they
measure where this machine cannot go.

\begin{table}[t]
\centering
\caption{Real proofs on a 24-core / 12\,GB CPU machine. ``Config'' is
characters $\times$ rules; $k$ abbreviates $10^3$. Every succeeded
point is at the \texttt{core} (non-zero-knowledge) tier.}
\label{tab:proofs}
\tabfont
\begin{tabular}{@{}lrrrl@{}}
\toprule
Config & Time & Size & Peak RSS & Result \\
\midrule
$200 \times 1$  & 123.5\,s & 2716.4\,KiB & 10{,}389\,MB & \checkmark \\
$200 \times 2$  & 118.9\,s & 2716.9\,KiB & 10{,}438\,MB & \checkmark \\
$2k \times 1$   & 138.1\,s & 2718.6\,KiB & 10{,}449\,MB & \checkmark \\
$10k \times 1$  & 172.5\,s & 2726.6\,KiB & 10{,}506\,MB & \checkmark \\
$200 \times 3$  & ---      & ---          & ---          & OOM \\
$20k \times 1$  & ---      & ---          & ---          & OOM \\
\bottomrule
\end{tabular}
\end{table}

\begin{itemize}
  \item \textbf{A fixed floor of $\approx$10.15\,GiB.} The smallest
  configuration---200 characters, one rule---already peaks at 10{,}389\,MB: the
  prover's own overhead (the trace and permutation structures of a \texttt{core}
  proof) is almost independent of the trace. All four successes fall inside the
  narrow band 10{,}389--10{,}506\,MB. Length and rule count have little marginal
  effect on memory---\emph{but the headroom is too thin for that to help}.
  \item \textbf{A two-step staircase, not a slope.} One rule is provable up to
  10k characters; two rules at about 200 characters; \textbf{three or more do not
  come out at all}. The third rule in the sweep is a \texttt{pattern\_block}: a
  regular-expression match activates another family of AIR chips and lifts the
  trace area in one step, so the cliff is at $2 \to 3$ rules; in length it is at
  $10\text{k} \to 20\text{k}$.
  \item \textbf{Memory binds, not CPU.} The cycle sweep of
  \S\ref{sec:eval-cycles} reaches 100k characters with 6 rules because it only
  executes in the zkVM and produces no proof; at 100k characters the cycle count
  is still on the order of $10^7$, which any CPU handles. \emph{The proving-side
  ceiling and the execution-side ceiling are in different places.} ``How large a
  policy can be proved'' is, on this machine, a fixed floor plus two steps rather
  than a cost curve.
  \item \textbf{Artefact size is nearly independent of configuration}: from 200
  to 10{,}000 characters the proof grows only from 2716.4 to 2726.6\,KiB. Pure
  verification takes $\approx$90\,ms (the \texttt{vkey} setup is 1.5\,s, once).
  \item \textbf{Reproducibility, stated precisely.} The \emph{position} of the
  wall reproduces (each failing point failed twice); the \emph{timings} at the
  wall do not. Re-running one group in isolation shows $\pm 10\%$ jitter in time
  and $\pm 1\%$ in memory, so we quote time only to an order of magnitude.
  \item \textbf{Session and multiparty.} The session aggregation proof (three
  certificates) measured 152.5\,s; multiparty verification needs one proof
  \emph{per non-empty slice}, so its cost grows linearly in the number of slices
  (a two-proof example policy). Neither introduces a new proving system---a slice
  and a session are ordinary inputs to the guest---but both are subject to the
  same floor: every additional proof adds another $\approx$10\,GiB peak, so
  multiparty proving must run in separate processes, serially.
\end{itemize}

\need{⚠️ 本节的 groth16 档与 ZK/非 ZK 交叉点属 WP1，尚未打通。骨架 :665--668 的
矩阵规格为规则 $\{1,2,3,4,6,8\}$ $\times$ 长度 $\{200,2k,10k,50k\}$ $\times$
模式 $\{$transparent, private, private-policy$\}$ $\times$ 证明档 $\{$core, groth16$\}$
（\texttt{bench/bench\_scale.py}）。表 \ref{tab:proofs} 只覆盖 core 档。
groth16 列落地前不得填入任何 ZK 数字。}

\subsection{The non-ZK baseline: what a signature cannot do}
\label{sec:eval-baseline}

A signed Merkle log is the obvious alternative to a proof, and on cost it is not
close: everything is milliseconds. We measured it so that the comparison is not
rhetorical. Policy package \texttt{eu\_ai\_act\_v1.json}, 103-character response,
one challenge nonce, a 1024-leaf log; each timing is the median of $n = 500$
repetitions.

\begin{table}[t]
\centering
\caption{Signed-Merkle baseline (WP9). No proofs, no zkVM: this arm is pure
Python, which is why the figures are milliseconds.}
\label{tab:baseline}
\tabfont
\begin{tabular}{@{}lrl@{}}
\toprule
Metric & Median & Note \\
\midrule
Sign certificate (Ed25519 + DSSE) & 0.132\,ms & $n=500$ \\
Merkle append (leaf + root recompute) & 2.131\,ms & $n=500$ \\
Public verification (sig + inclusion) & 0.322\,ms & $n=500$ \\
Re-run verification ($+J(\pi,T)$) & 0.772\,ms & $n=500$ \\
Certificate size & 648\,B & \\
Inclusion proof (10 siblings) & 761\,B & \\
\bottomrule
\end{tabular}
\end{table}

\need{⚠️ 引用 2.131 ms 时必须写明它是 \texttt{proofs/session.py} 的\emph{静态树重算}
$\Theta(n)$，\emph{不是} RFC-6962 增量日志的摊销 $\Theta(\log n)$。
⚠️ 共享机器有 $\pm 20\%$ 漂移。数据源 \texttt{bench/results/baseline.md}。}

The asymmetry is the point, and it is what this arm contributes. The verification
chain is \emph{verify signature} $\to$ \emph{verify log inclusion} $\to$
\emph{re-run $J(\pi,T)$}. The first two steps need only public data; the third
needs the plaintext of both $\pi$ and $T$. That is not an implementation defect
but a structural consequence of the signature paradigm: a signature binds
\emph{who said it} to the bytes, never \emph{whether it is true}. We measured the
consequence rather than asserting it---\textbf{the same key signs a certificate
whose declared content is false, and the public path passes green}, because what
it sees is a certificate with a valid signature and a valid inclusion proof. Only
re-executing $J(\pi,T)$ recovers the truth, and taking that step away from the
verifier is exactly what the zero-knowledge design buys. The machine-checked
assertions are enumerated in \texttt{bench/results/baseline.md}; one row of that
file is worth repeating because it is the honest one: on
\emph{unforgeability} the baseline is \emph{not weaker} than the proof path, both
being Ed25519 and vkey binding, and we do not score a point for ZK there.

\subsection{End-to-end and auditing}
\label{sec:eval-e2e}

The end-to-end demo produces 14 certificates (streaming chain with early stop,
MCP argument and response checks, one zk proof) and anchors them; a third party
verifies with public artefacts alone, and the ledger chain, signature,
\texttt{policy\_hash}, anchoring, stream chains and zk proof all pass. In the
on-chain variant (a local Anvil node, \texttt{Anchor.sol} deployed, 14
certificate digests registered) the verifier reads back each certificate's
\texttt{anchoredAt} and cross-checks the local ledger, passing 14/14; the
negative control (an unregistered digest reads back as zero) confirms the check
is not vacuously true. Cross-validation submits 19 vectors to the Python
reference decision and to both the SP1 host and proving paths, comparing
\texttt{passed} and the violated rule set row by row; a full re-run with real
\texttt{core} proofs gave 19/19 on both paths.

\subsection{Measuring \texorpdfstring{$k$}{k}}
\label{sec:eval-k}

\need{对一组真实 prompt、同一模型多次采样，估 $\mathsf{Guess}(T\mid Z)$，
代入 Thm \ref{thm:t1}(H) 的阈值 $k \ge \log q + \log(1/\varepsilon)$（承接
Rem \ref{rem:kexp}）。据先例扫查无人测过这个量，本身是一块贡献；
同时它也是 Prop \ref{prop:semantic} 的测量（发布 $\mathsf{encode}(T)$ 后
$T$ 的可恢复比例）。\textbf{⚠️ 需真实模型多次采样，当前阻塞（外部资源）。}}

\need{⚠️ 「Why zero-knowledge?」决策图（骨架 :670--675）\emph{未}放入本节：
它讲的是设计结论（ZK 只在 $\{$私有响应, 私有策略$\}$ 两角有正收益），是结论不是测量，
按 docs/tdsc-writing-plan.md 的建议移至 §7 Discussion。与骨架现状不同，\textbf{待作者认可}。}

% =====================================================================
\section{Related Work}
\label{sec:related}

We position PoP against five lineages. For each, Table~\ref{tab:related} records
the slot it occupies---which is also, read the other way, what this paper does
not have to claim.

\begin{table}[t]
\centering
\caption{Positioning against the five prior-art lineages. The right-hand column
names what each lineage already occupies, and therefore what we do not claim.}
\label{tab:related}
\tabfont
\setlength{\tabcolsep}{3pt}
\begin{tabular}{@{}L{0.26\columnwidth}L{0.42\columnwidth}L{0.26\columnwidth}@{}}
\toprule
Lineage & Representative work & Slot already occupied \\
\midrule
Policy $\to$ circuit & zkGen \cite{zkgen}, Circuitree, Prezta & DSL-to-circuit
translation \\
Agent/LLM compliance & NiyamAI, Zaguk, PCO \cite{pco}, WITNESS & policy-level
private-mode kernel \\
Model computation & DeepProve, zkGPT, zkAgent \cite{zkagent} & inference
integrity \\
Attacks / negative results & Hollow-LLM \cite{hollowllm}, zk-attach
\cite{zkattach} & naming the gaps \\
Non-ZK industry baseline & \texttt{@brivora/verify}, VAOL & signed audit logs \\
\bottomrule
\end{tabular}
\end{table}

One of the five is the lineage this paper is actually in dialogue with. The
attacks-and-negative-results row names two gaps---Hollow-LLM's \emph{effort gap}
and zk-attach's \emph{source-integrity gap}---and both concern \emph{what the
proof system was asked to prove}. Proposition~\ref{thm:gap} concerns a third
gap on a different axis, \emph{where the proved string is delivered}, and we are
not aware of prior work that quantifies the binding/hiding tradeoff in the terms
of \S\ref{sec:binding}. That is why we state it as a bound with a stated
hypothesis rather than cite it.

\paragraph{zkAgent: orthogonal, not competing.} zkAgent \cite{zkagent} is the
closest contemporaneous system, and the comparison is easy to get wrong in both
directions. It proves that a provider executed the declared model and tools; PoP
proves that a response satisfies a public policy. PoP proves nothing about model
computation, and zkAgent proves nothing about policy compliance, so the two
compose: their conjunction covers ``inference was executed \emph{and} the policy
held''. Table~\ref{tab:zkagent} places the reported figures side by side.

\begin{table}[t]
\centering
\caption{PoP against zkAgent. The two prove different obligations, so the rows
are a statement of orthogonality, not a race.}
\label{tab:zkagent}
\tabfont
\setlength{\tabcolsep}{3pt}
\begin{tabular}{@{}L{0.22\columnwidth}L{0.38\columnwidth}L{0.30\columnwidth}@{}}
\toprule
& \textbf{zkAgent} \cite{zkagent} & \textbf{PoP} \\
\midrule
Obligation & provider executed the declared model and tools (inference integrity)
& response/trace satisfies policy $\pi$ (compliance) \\
Depends on the model & yes (weights, arithmetisation) & \textbf{no} \\
Cost grows with & model size, sequence length & response length $\times$ rule
count \\
Content privacy & none & committed, selective disclosure \\
Hardware reported & 32 cores / 512\,GB & 24 cores / 12\,GB \\
Proving & 194.4\,s (Shout, $T{=}512$) & 118.9--172.5\,s (200--10k chars) \\
Proof size & 96\,MiB (Shout); 3.1\,MiB (LogUp) & $\approx$2.7\,MiB \\
Verification & 0.42\,s (Shout); 0.038\,s (LogUp) & $\approx$0.090\,s \\
\bottomrule
\end{tabular}
\end{table}

We deliberately do not claim that PoP is faster, and the temptation to do so is
worth naming because the raw numbers invite it. The correct statement is
narrower, and we make only it: \emph{for the compliance obligation}, PoP
completes in a proving time of the same order as zkAgent's, with a smaller memory
footprint and a comparable proof size, without touching the model---and it adds
committed content privacy that zkAgent does not provide. Two caveats must travel
with those figures. The hardware is not matched (24 cores and 12\,GB against 32
cores and 512\,GB), and PoP's ceiling is memory-bound rather than cycle-bound, so
its feasible region is a two-step staircase rather than a curve.
\S\ref{sec:eval} reports the measurements and their limits.

\need{⚠️ tab:zkagent 里 PoP 那一列全部取自 SP1 \texttt{core}（非零知识）档，
机器为 24 核 / 12\,GB WSL2；zkAgent 一列为 ePrint 2026/199 的报告值（32 核 / 512\,GB）。
数据源：\texttt{bench/comparison\_zkagent.md}。该文件自身警告：
\emph{不得}宣称「PoP 比 zkAgent 快」——正确口径是「对『合规』这一义务…
同量级证明时间 / 更小内存」。groth16 档（WP1）未打通前，本表不得出现 ZK 档数字。}

The one influence we should name plainly, rather than leave to resemblance:
zk-regex's ``match offline, verify the path in circuit'' pattern is the shape of
our NFA contract, and verifiable document redaction \cite{vdr} inspired the bit
selector behind \texttt{mask\_covered}. The borrowing stops at the circuit
technique and does not carry the security argument across: VDR assumes a
\emph{signed original document} exists, whereas in PoP's private mode both the
original and the redacted version are produced by the operator, so there is no
trusted original to appeal to. We use the mechanism, not the proof.

\fixme{引文纪律：除 Hollow-LLM \cite{hollowllm} / zkGen \cite{zkgen} /
zkAgent \cite{zkagent} / PCO \cite{pco} 四条外，其余引文必须逐条打开原文
核实后再进 bibliography（docs/tdsc-fixlist.md \S2、WP8）。
\textbf{骨架原先把「ePrint 2024/813」列进已核实四条，那是错的}，已剔除：
\textbf{绝不能再引 IACR ePrint 2024/813 作为 VDR}——那是
Kiraz/Larraia/Vaughan, ``How to Redact the Bitcoin Backbone Protocol''，
与 Verifiable Document Redacting 无关。见 fixlist \S A。
⚠️ 口径：WITNESS 参考实现为 Apache-2.0（并非闭源），故「开源」不构成差异化。}

% =====================================================================
\section{Discussion and Limitations}
\label{sec:discussion}

Everything above was kept deliberately narrow, and this section says how narrow.
It covers the four places where a reader is most likely to round our claims up:
the caliber of the privacy statement (\S\ref{sec:disc-privacy}), where
zero-knowledge actually pays (\S\ref{sec:disc-zkmap}), availability against an
adversarial response (\S\ref{sec:disc-dos}), coverage
(\S\ref{sec:disc-coverage}), and composition with an inference proof
(\S\ref{sec:disc-composition}). It closes with the engineering limitations that
remain (\S\ref{sec:disc-rest}).

\subsection{The caliber of the privacy claims}
\label{sec:disc-privacy}

\textbf{``Content privacy'' means \emph{no plaintext of $T$ in the public
values}, not \emph{content is unrecoverable}.} Two independent reasons force
this reading---one at the proof-system level, one at the design level---and the
weaker reading is the one our own measurements support.

\emph{Proof-system level.} The default proving tier is SP1's \texttt{core}, and
SP1's published security model states explicitly that individual STARK proofs do
not currently satisfy the zero-knowledge property; the source agrees, with no
blinding anywhere in the proof stack and the shard proof carrying the trace
Merkle root and trace openings in the clear.%
\need{⚠️ 引文债：「SP1 公开安全模型承认其 STARK 证明当前不满足零知识」是
\emph{外部}断言，出处（SP1 官方安全模型文档的具体章节/URL）尚未核实到可引的
一级来源，故此处只陈述、不给 \cite{}（WP8）。若核不到原文，须撤回该句，
只保留我们自己的源码审计结论那半句。}
Soundness (G1) therefore holds
while the proof artefact itself does not hide the witness. The
\texttt{gr\-oth16} and \texttt{plonk} tiers wrap the inner STARK proof as a
private witness---the only mode in which hiding is even possible---but that is a
wrapper-level claim, it is not post-quantum, it needs $\ge 16$\,GB, and we did
not measure it (\S\ref{sec:eval-proofs}).

\emph{Design level, and the more fundamental one.} The public values
\texttt{response\_binding} and \texttt{response\_commitment} are both
\emph{public, offline-recomputable functions of $T$}. This is not an oversight:
\texttt{response\_binding} \emph{must} be publicly recomputable, because that is
how a verifier independently checks that the delivered $T'$ is the $T$ that was
proved---the swap-response defence. But the same property is a guess-and-verify
oracle for $T$.

It is worth being precise about the attack, because the obvious version of the
argument is wrong. ``A one-line response carries only $10^2$--$10^3$ bits of
entropy, far below 256, so enumeration is feasible'' does \emph{not} follow:
$2^{100}$ is as unenumerable as $2^{256}$. The real attack is
\emph{re-execution} (Corollary \ref{cor:reexec}). The adversary's side
information is not the response space but the tuple (prompt, model, decoding
parameters) (Assumption \ref{as:Z}); under greedy or low-temperature decoding
the response is nearly determined, so $k = -\log_2 \mathsf{Guess}(T \mid Z)
\approx 0$ (Definition \ref{def:k}) and the attack costs \emph{one inference},
independent of $|T|$. Under the two premises that a verifier independently
recomputes the binding and that $k < \log q + \log(1/\varepsilon)$, response
binding and response-content hiding are mutually exclusive.

We therefore restate (G2) as what is actually established: \emph{in private mode
the public values contain no plaintext of $T$}---which the leak experiment
confirms directly---and not as content irrecoverability. The consequence for the
comparison in \S\ref{sec:related} is that the ``content privacy'' cells on
\emph{both} sides must not be read as ``content is unrecoverable.'' This is a
bound on our own system first.

\subsection{Where zero-knowledge pays, and where it does not}
\label{sec:disc-zkmap}

Section \ref{sec:eval-baseline} measured the non-ZK alternative---a signed
Merkle log---and everything in it is milliseconds. A proof therefore has to earn
its cost against what a signature cannot do, and our measurements locate that
boundary instead of asserting it.

The map has two axes: is the response secret, and is the policy secret?
Zero-knowledge has positive value only in the corners where at least one of them
is unknown to the verifier. In the transparent corner---public response, public
policy---a signed log certifies exactly the same thing at three to four orders
of magnitude less cost, and adding a proof is redundancy rather than rigour. The
familiar criticism that ``a signature would have done'' is correct in that one
corner and wrong in the other three; the decision map is what turns the
criticism into a measured design conclusion.

One dependency belongs in the open. The zero-knowledge property of a certificate
is a property of the \emph{wrapper}, not of the system as a whole, and the
default tier does not have it (\S\ref{sec:disc-privacy}). If \texttt{gr\-oth16}
does not become a delivery tier, we do not claim the artefact is
zero-knowledge, and the word must leave the title.

\subsection{Availability: the unanchored-pattern cost attack}
\label{sec:disc-dos}

A regular expression with an unanchored pattern---one with no left anchor, such
as \texttt{.*foo}---gives the matcher no way to reject a prefix of the response
early: every position stays a live start position until the whole string has
been read. The prover's work therefore scales with the response, and a party who
\emph{chooses} the response can drive it quadratically. In the agent setting
that party is often the agent itself, or whoever controls its output.

Two things are worth separating. This is a denial of service against the
\emph{prover}, not an attack on soundness: no false certificate is produced, and
the cost is simply borne by the honest side. And it is specific to proving
rather than to matching, because the trace---not merely the verdict---is what
gets proved. Mitigations exist and are partly in place: a Pike VM (already
implemented) simulates the NFA in linear time, committing to the response length
up front removes the adversary's freedom to inflate it, and chunked matching
bounds the per-chunk work. Prior work on zero-knowledge regular expressions
imposes the DFA restriction but does not analyse denial of service in a zkVM
setting, which is where this differs.

\need{⚠️ 作者决策未定：骨架 :649--651 把这条写成一条定理（$\Omega(n^2)$ 证明
工作量），但证明没写。需定：写成定理（并补证明），还是降为观察（Observation）。
本稿先按观察陈述、不给定量断言，以免留下一条无证明的定理。}

\subsection{Coverage: a provider can simply withhold certificates}
\label{sec:disc-coverage}

Nothing in the cryptographic layer can compel a provider to issue a certificate
at all. A provider that never produces one is indistinguishable, from the
verifier's side, from a provider with nothing to certify. We treat this as an
open problem rather than a defect of the construction: without trusted hardware
there is no anchor that forces issuance, and the mechanisms that could---%
sampling audits, economic penalties, disclosure obligations---live above the
cryptographic layer and outside the scope of this paper. Naming it as an open
problem is more useful than proposing a mitigation we cannot evaluate.

\subsection{Why composition alone does not close the gap}
\label{sec:disc-composition}

The natural next step is to compose: prove that the inference produced $T$ from
the request, conjoin that with a proof that $\pi(T)$ holds, and conclude that
the agent \emph{sent} a compliant response. We describe the mechanism, and we
record why the conclusion does not follow for free.

\emph{Composition by itself does not close the gap.} A model proof bound to
\emph{some} input shows only that \emph{some} input yields a compliant output;
the divergent-delivery gap of Proposition \ref{thm:gap} closes only when the
model proof is bound to the \emph{input the user submitted}, $x$, so that
$M(x) = T \wedge \pi(T)$ becomes a statement about the bytes actually sent. That
requirement is not an extra condition bolted on: it is exactly $f > 0$ in
Theorem \ref{thm:t3}, with the user-submitted $x$ playing the role of the $f$
fresh bits. The two halves of this paper meet here---the request-side theorem
names the condition the composition would have to satisfy, and the response-side
theorem is what makes that condition non-vacuous.

The composition contract is correspondingly short to state: the inference proof
must expose $T$ in its public values, and it must be bound to $x$.

\need{⚠️ 阻塞：真模型规模（$\ge 10^6$ 参数）与组合成本数据属 WP4 子问题② /
WP3 的 DeepProve 对照臂，需外部资源。代理规模下「组合成本由推理证明主导」
\emph{未被验证}（\texttt{bench/results/compose.md}），不得写入。
骨架 :641--642 的组合契约定义可在此处补。}

\subsection{Remaining limitations}
\label{sec:disc-rest}

\begin{itemize}
  \item \textbf{Three-party collusion is undetectable at the signature layer.}
    If the policy author, the operator, and the prover all collude, the
    certificate is indistinguishable from an honest one. The multi-prover
    construction partitions responsibility; it does not prevent agreement
    (\texttt{docs/security-model.md}).
  \item \textbf{Signing and key custody are not production-grade.} The on-chain
    path currently takes plaintext private-key parameters, the demo uses a
    public Anvil test key, and deployment is a local Anvil instance. Production
    signing (Ed25519/HSM) and key custody remain future work.
  \item \textbf{Semantic-rule model quality.} The delegated semantic rule works
    end to end, but the bundled model is a small demonstration proxy rather than
    a usable learned guardrail: ``$P(\text{harmful})$ is low'' is a statement
    about \emph{this} model. The companion proof's public instance also contains
    $\mathsf{encode}(T)$, so semantic rules are \emph{public-mode only}.
  \item \textbf{The inference proof is faithful, not good.} The inference half
    is a deterministic fixed-point MLP proxy whose weights are generated from a
    compile-time seed and baked into the program, hence committed by the
    verification key. It proves that \emph{this} graph computed \emph{this}
    output on \emph{this} response---not that the graph is any good. Swapping the
    proxy for a real prover touches only the prover; the composition layer does
    not move.
  \item \textbf{Verification remains $O(N)$.} There is no recursive aggregation,
    so a verifier pays per certificate. The supported regex subset and its ASCII
    semantics are bounded in the same way, and proof-cost optimisation (lookup
    arguments, parallelism, precomputation) is unexplored.
  \item \textbf{Slices share one guest.} Multiple provers hold distinct keys,
    but that is a division of responsibility rather than a per-role circuit: the
    repository does not build one circuit per role.
\end{itemize}

% =====================================================================
\section{Conclusion}
\label{sec:conclusion}

This paper began from a small question---what does a compliance certificate
actually bind?---and the answer came out the same on both sides of the
interaction. A proof binds only the bits that were submitted into it \emph{and}
that the adversary could not select in advance. On the response side those bits
are the entropy $k$ of $T$ given everything the adversary already knows, and no
amount of computation buys more of them; on the request side they are the $f$
bits of freshness the receiver contributes, and a certificate without them is
replayable without limit. Theorem \ref{thm:t1} and Theorem \ref{thm:t3} are two
faces of that one statement, and Proposition \ref{thm:gap}---a gap we state
against our own design---is what makes it bite.

PoP is the constructive instance of that analysis, not a claim to have invented
a new proof system. It composes an existing zkVM, an existing policy compiler,
and a signature layer into a certificate that a verifier holding only public
values can check without trusting hardware, and it is explicit about what such a
certificate does not do: it does not make the response unrecoverable, it does not
hide the witness at the default proving tier, and it cannot compel a reluctant
provider to issue anything at all. Those are limits of the problem, stated as
limits, rather than of one implementation.

The design conclusions are accordingly modest, and we hope reusable: publish the
leak profile of a certificate and audit it like a disclosure policy; let a
privacy claim be only as strong as its arithmetic supports; and put the
freshness where the binding is. A certificate is worth exactly as much as the
bits that went into it.

% =====================================================================
% ⚠️ IEEEtran 的 \appendix 是「单个附录」命令：它把 \section 重定义成
% \@IEEEdestroythesectionargument，直接丢弃 \section{...} 的参数（只打 warning）。
% 多个附录必须用 \appendices（复数），它把 section 计数器置 -1，
% 之后每个 \section{X} 才依次成为 A、B、……  （2026-10-12 实测：
% 用 \appendix 时 A/B 两个附录的 \ref 都解析成 "A"，且 title 字段退化成 \appendixname，
% 锚点退化成 section*.1；「0 undefined ref」检查抓不到这种错。）
\appendices
\section{Proofs}
\label{app:proofs}

\begin{proof}[Proof of Theorem \ref{thm:t1}]
\emph{(H).} Fix $z$ and let the adversary's query set be
$Q = \{q_1, \dots, q_q\}$. Decompose by first hit,
\[
\Pr[\exists i : q_i = T] \;=\; \sum_i \Pr\!\left[q_i = T \wedge \neg E_{<i}\right],
\]
where $E_{<i} \triangleq \{\exists j < i : q_j = T\}$.
Conditioned on $\neg E_{<i}$, every oracle answer $a_1,\dots,a_{i-1}$ already
revealed is uniform and independent of $T$ (no queried point equals $T$, and $H$
is a fresh random function), and $\beta(T)$ itself is uniform and independent of
$T$. Hence $q_i$---a deterministic function of $(z, \beta(T), a_{<i})$---is
\emph{independent of $T$}, so
\begin{multline*}
\Pr[q_i = T \mid Z = z]
= \sum_t \Pr[q_i = t]\,\Pr[T = t \mid Z = z] \\
\le \max_t \Pr[T = t \mid Z = z] = 2^{-k_z},
\end{multline*}
giving $\Pr[\exists i : q_i = T \mid z] \le q\,2^{-k_z}$. Taking the expectation
over $z$ and using $\mathbb{E}_z[2^{-k_z}] = \mathsf{Guess}(T \mid Z) = 2^{-k}$
from Definition \ref{def:k} yields $q \cdot 2^{-k}$. An adversary that outputs
an unqueried value wins with probability at most $2^{-k_z} \le q\,2^{-k_z}$,
which the same bound absorbs. The term $q \cdot 2^{-\lambda}$ charges the
probability that a query hits the uniform value $\beta(T)$ by chance.

\emph{(L).} Immediate from Definition \ref{def:k} after conditioning on
$\beta(T)$; no hypothesis on $\beta$ is used.
\end{proof}

\need{Thm \ref{thm:t1}(L) 无闭式界，只给测量口径（\S\ref{sec:eval}）。
三分法已降级为分类 Rem \ref{rem:taxonomy}，无需证明。
原 T2「分离定理」已\emph{停用}（见 Rem \ref{rem:designated}）—— 不要为它写附录。}

\need{\textbf{待补}：把 Thm \ref{thm:t1}(H) 的随机预言机模型换成标准模型，
或明确声明本文只在 ROM 下主张。当前证明依赖 $H$ 是新鲜随机函数这一条；
实践中的 SHA-256 只给计算意义下的近似，故正文措辞必须写成
\emph{heuristic}，不得写成信息论界。}

\section{Reproduction}
\label{app:repro}

The full walkthrough---environment, a single compliance proof, cross-checking
against the reference layer, private mode, certificate issuance, the one-command
demo, and third-party verification---is maintained in \texttt{docs/reproduce.md}
together with the exact commands and their expected output. This appendix records
only what a reader needs in order to judge the claims of \S\ref{sec:eval}.

\paragraph{Environment.} Ubuntu or WSL2 on x86-64, no GPU. The Rust/SP1 side
needs Rust stable, \texttt{cargo-prove}~v6.7.0, Go $\ge$\,1.24, and the usual
native toolchain; the reference layer is pure Python~3.10+ and runs on its own.
Memory, not time, is the binding constraint: the CPU prover carries a fixed
floor of approximately 10.15\,GiB regardless of the trace
(\S\ref{sec:eval-proofs}), and 12\,GB suffices only for small policies.

\paragraph{What reproduces without the prover.} The reference layer, the policy
compiler, and the whole test suite run in seconds and need no Rust. This is
enough to reproduce the non-ZK baseline of \S\ref{sec:eval-baseline}, to inspect
the certificate and inclusion proofs, and to check the cross-layer contract---%
that the Python reference implementation and the zkVM guest consume the same
compiled policy bytes.

\paragraph{What needs the prover.} Every timing, size, and memory figure in
\S\ref{sec:eval-cycles}--\S\ref{sec:eval-proofs}, and the end-to-end result of
\S\ref{sec:eval-e2e}, requires a real proof. The measurements were taken on a
24-core CPU machine with 12\,GB of RAM, which is what places the ceiling reported
there: the successes and the two out-of-memory configurations in
Table~\ref{tab:proofs} are both properties of that machine.

\paragraph{Where each number comes from.} Every measurement in \S\ref{sec:eval}
is the content of a checked-in result file with a machine-readable companion:
\texttt{bench/results/} holds \texttt{cycles}, \texttt{ablation},
\texttt{proofs}, \texttt{baseline}, \texttt{verify}, \texttt{semantic},
\texttt{streaming}, \texttt{compose}, and \texttt{prover\_knobs}, each as a
\texttt{.json} alongside its \texttt{.md}. The reproduction guide records the
commit, the hardware, and the prover binary digest for each run, so that a
reported figure can be traced to one specific execution rather than to a state
of the tree.

\need{⚠️ 数字漂移（本轮实测发现，需作者定夺）：正文 §2.2 写「801 tests / 15 skipped」，
其依据是实测（\texttt{docs/dev-plan.md:506} 的落地记录）\emph{以及}本轮对
\texttt{tests/} 的静态清点（\texttt{def test\_} 计 801 个）——\textbf{801 是对的}。
但 \texttt{docs/reproduce.md:31,428}、\texttt{docs/security-model.md:829}、
\texttt{docs/development.md:93,290}、\texttt{docs/modules/08-tests-bench.md:11,58,504,526}、
\texttt{docs/modules/README.md:94,161,201} 仍写 \textbf{783}（陈旧值，
\texttt{docs/report/advisor-report.md:586} 已把这条「文档数字漂移」登记在案）。
本附录既然指向 \texttt{reproduce.md}，照抄就会让复现者跑出 801 却以为环境坏了——
\textbf{这是真实的复现缺陷}。按仓库自己的「过期就标注、不静默改」纪律，
本稿不擅自改那些文档，只在此记债。作者需定：把五处文档回填为 801
（\texttt{advisor-report.md:599} 的建议），还是在附录里并列两套口径。}

% =====================================================================
\begin{thebibliography}{99}

\bibitem{hollowllm}
C.~Gong, B.~Liu, and M.~Li, ``Hollow-LLM Attack: Computationally Trivial
Weights in Zero-Knowledge Verification of LLM Inference,'' in \emph{IEEE S\&P},
2026, pp.~4339--4355, DOI 10.1109/SP63933.2026.00258.

\bibitem{zkattach}
\need{zk-attach, arXiv:2608.30083 --- 打开原文核实作者后再引}

\bibitem{zkgen}
J.~Lauinger, J.~Ernstberger, and S.~Steinhorst, ``zkGen: Policy-to-Circuit
Transpiler,'' in \emph{IEEE ICBC}, 2024, pp.~620--624,
DOI 10.1109/ICBC59979.2024.10634440.

\bibitem{pco}
\need{PCO, IACR ePrint 2026/994 --- 打开原文核实作者（Dolev/Yung 等）后再引}

\bibitem{zkagent}
L.~Wang, H.~Lou, C.~Li, Y.~Yu, and Y.~Hu, ``zkAgent: Verifiable LLM Agent
Execution via One-Shot Transcript Proofs,'' IACR ePrint 2026/199.

\bibitem{hiddencred}
D.~Balfanz and J.~Schweitzer, ``Hidden credentials,''
\need{核实出处：2002 年 unpublished manuscript / 2003 年会议版本，二者择一并核准页码}.

\bibitem{handshake}
\need{核实作者列表与页码（首发 S&P 2003, pp.~139--150）}D.~Balfanz, G.~Durfee,
N.~Shankar, D.~K.~Smetters, J.~Staddon, and H.-C.~Wong, ``Secret handshakes from
pairing-based key agreements,'' in \emph{IEEE S\&P}, 2003.

\bibitem{xacmlpriv}
\need{核实作者与出处。骨架点名的方向是「XACML 策略的隐私保护访问控制」，
候选：C.~A.~Ardagna 等关于 XACML 隐私的系列工作。\textbf{未核实前不得写成
具体条目}。}

\bibitem{sw05}
\need{核实页码（LNCS 3494, pp.~457--473）}A.~Sahai and B.~Waters, ``Fuzzy
identity-based encryption,'' in \emph{EUROCRYPT}, 2005.

\bibitem{ksw08}
\need{核实页码（LNCS 4965, pp.~515--534）}J.~Katz, A.~Sahai, and B.~Waters,
``Predicate encryption supporting disjunctions, polynomial equations, and inner
products,'' in \emph{EUROCRYPT}, 2008.

\bibitem{vdr}
\need{Verifiable Document Redacting, Springer, DOI 10.1007/978-3-319-66402-6_20
--- 核实作者与年份。\textbf{绝不能再引 IACR ePrint 2024/813}（那是
Kiraz/Larraia/Vaughan, "How to Redact the Bitcoin Backbone Protocol"）。
见 docs/tdsc-fixlist.md \S A。}

\end{thebibliography}

\end{document}
"   % 两遍，第二遍收交叉引用
%  \SKELQUIET 让所有中文标记宏吞掉参数；不带它则需要能排版中文的引擎
%  （xelatex + CJK 字体）。正文本身全英文，标记里才有中文。
%
%  ⚠️ IEEEtran.cls 需要装在 TeX 发行版里。WSL 侧已装进
%     /home/dong/texlive/texmf-local/tex/latex/ieeetran/（2026-10-12）。
% =====================================================================
\documentclass[journal]{IEEEtran}

\usepackage{amsmath,amssymb,amsthm}
\usepackage{booktabs}
\usepackage{array}
\usepackage{graphicx}
\usepackage{xcolor}
\usepackage{hyperref}

% 正文用 Times（ptm，窄），\texttt 用 Courier（pcr，等宽且宽）。
% 一行里只要有一个长 Courier 标识符（\texttt{normalized\_keyword\_block}
% 这类）就几乎没有可拉伸的胶，TeX 只能拉出 badness 10000 的大空隙。
% emergencystretch 允许第三遍在这些行上额外拉伸，把最坏情况吸收掉；
% 代价是那些行略松，远好于一道 10000 的豁口。
% （2026-10-12：装上真正的 Times 后此类行从 3 涨到 12 才暴露。）
\setlength{\emergencystretch}{2em}

% microtype：字形微伸缩（pdftex 下 ±2%）+ 标点悬挂。
% emergencystretch 只能吸收「可拉伸胶不足」，但 badness 在 10000 处封顶，
% 段首被后面长片段（consequences---customer-facing、normalized\_keyword\_block）
% 挤空的行照样报 10000。microtype 从字形宽度上补，是这一类的标准解法。
\usepackage{microtype}

% 窄栏里用 p{} 会因两端对齐产生大片字距（编译时刷一堆 Underfull \hbox）。
% 表格正文列一律用 L{}：左对齐、允许自然参差。
\newcolumntype{L}[1]{>{\raggedright\arraybackslash}p{#1}}
\newcommand{\tabfont}{\footnotesize}

\newtheorem{theorem}{Theorem}
\newtheorem{lemma}[theorem]{Lemma}
\newtheorem{proposition}[theorem]{Proposition}
\newtheorem{corollary}[theorem]{Corollary}
\newtheorem{definition}[theorem]{Definition}
\newtheorem{assumption}[theorem]{Assumption}
\newtheorem{remark}[theorem]{Remark}

% --- 工作标记：提交前全部清掉 ---
% 注意：宏名不能叫 \wp —— \wp 已是 LaTeX/amssymb 的 Weierstrass ℘，
% \newcommand{\wp} 会直接报 "Command \wp already defined" 并使后文级联报错。
\ifdefined\SKELQUIET
  \newcommand{\wpmark}[1]{}
  \newcommand{\need}[1]{}
  \newcommand{\fixme}[1]{}
\else
  \newcommand{\wpmark}[1]{\textcolor{red}{\textsuperscript{[WP#1]}}}
  \newcommand{\need}[1]{\textcolor{blue}{[[ NEED: #1 ]]}}
  \newcommand{\fixme}[1]{\textcolor{orange}{[[ FIX: #1 ]]}}
\fi

\begin{document}

% ---------------------------------------------------------------------
% 标题：决策 D5 —— 保留 "Zero-Knowledge" 的前提是 groth16 成为交付档 (WP1)。
%       若 WP1 未打通 groth16，改用下面被注释掉的那个标题，
%       并同步删掉摘要与 §8 里的 zero-knowledge 措辞。
% ---------------------------------------------------------------------
\title{Proof-of-Policy: Zero-Knowledge Compliance Certificates\\ for LLM Agent Outputs}
% 备选标题（若 groth16 未打通）：
% \title{Proof-of-Policy: Succinct Compliance Certificates for LLM Agent Outputs}

\author{\need{authors}}
\maketitle

% ---------------------------------------------------------------------
\begin{abstract}
% 决策 E2：内容隐私上界必须出现在摘要，不能埋进 §Limitations。
% 决策 D1：贡献陈述按"分析 x 系统各半"写 —— 分析在前，系统是分析所需的那一实例。
Proof-of-Policy (PoP) lets an auditor check that an LLM agent's response satisfies
a published policy, without trusted hardware and without the operator disclosing
the response in the clear. Our contribution is analytical before it is
constructive. We give a quantitative bound on what such a certificate binds---a
guessing-entropy term that no additional computation can improve---together with
the structural assumption the bound requires, and a proof that the assumption is
necessary rather than convenient. The same accounting governs the request side,
where a certificate without receiver-supplied freshness is replayable without
limit. PoP is then presented as the instance these bounds require---a policy
compiler, a
% 决策 D5 条件措辞：若 groth16 未成为交付档，把下行的 "zero-knowledge" 换成 "succinct"，
% 并同步改标题与 §8。
zero-knowledge circuit, and a signed certificate layer---measured end to end
rather than asserted, with the leak profile of the certificate reported as a
result in its own right.

\smallskip
\noindent\textbf{Scope of privacy.} PoP provides \emph{policy}-zero-knowledge:
the certificate reveals neither the policy evaluation trace nor, in
policy-private mode, the policy itself. It does \emph{not} guarantee that $T$
is unrecoverable: as we show in \S\ref{sec:binding}, any publicly recomputable
binding to $T$ admits a re-execution attack whenever the adversary knows the
prompt and the model. The same bit-budget governs the request side: a
certificate is replayable without limit unless the receiver contributes fresh
randomness (Theorem \ref{thm:t3}).
\end{abstract}

\begin{IEEEkeywords}
zero-knowledge proofs, LLM agents, policy compliance, verifiable computation,
accountability
\end{IEEEkeywords}

% =====================================================================
\section{Introduction}
\label{sec:intro}

LLM agents increasingly take actions that carry con\-se\-quences---\allowbreak{}
customer-facing replies, tool invocations, financial and medical content. Regulation such as the
EU AI Act (Articles~12 and~13) and enterprise governance both require that
compliance be auditable. In practice it can only be taken on trust: logs can be
forged, guardrail claims cannot be checked, and content privacy forbids retaining
the plaintext that would settle the question.

Four lines of work each occupy one corner of this problem without spanning it.
\emph{TEE execution proofs}---Proof-of-Guardrail and its relatives---show that
guardrail code was executed and that inputs and outputs were signed; as that work
itself observes, \emph{execution is not compliance}, and the approach carries
roughly $34\%$ latency overhead together with a hardware dependency.
\emph{TEE governance stacks} (Agentrust, cMCP, cA2A) inherit both the hardware
and the vendor facility. \emph{Formal verification} (Lean-Agent and similar)
requires the entire policy to be formalised by hand, and does not apply to
black-box responses. \emph{ZK regex and PII primitives} (zk-regex~V2, verifiable
document redaction \cite{vdr}, ZK Email) are building blocks; none of them is a
policy-level system.

The nearest work, WITNESS, supports two policy primitives, has no peer-reviewed
paper, and---by its own project page---ships an Apache-2.0 reference
implementation, so ``open source'' is not a differentiator; it also offers no
content privacy. \texttt{@brivora/verify} is not zero-knowledge, and DeepProve
proves the \emph{model computation} layer rather than the \emph{policy
compliance} layer. The combination of a general policy DSL compiled to ZK
constraints, \emph{without} a TEE, with selective disclosure of response content
and a formal security model, is an explicit vacancy. This paper occupies it with
PoP---but the system is not the contribution we lead with. Leading with a system
invites the objection that a signature would do; leading with the analysis is
what makes the system necessary, and it is the analysis that this paper is
organised around.

\need{⚠️ 上面点名的 Proof-of-Guardrail / Agentrust,cMCP,cA2A / Lean-Agent /
zk-regex V2 / ZK Email / WITNESS / \texttt{@brivora/verify} / DeepProve，
以及 EU AI Act Art.~12/13，\emph{均未给 \cite{}}：条目尚未逐条核实
（docs/tdsc-fixlist.md \S2、WP8）。核实前不得凭空写 bibitem，故此处只留名。}

\smallskip
\noindent\textbf{What we do \emph{not} claim.} We do not claim that policy proofs
bind nothing: our own system already binds the response through a verifier-issued
nonce (Definition~\ref{def:system}). We do not claim to be the first to observe
that a proof can fail to bind the artifact it purports to describe: the
\emph{effort gap} for ZK inference proofs is formalised by Hollow-LLM
\cite{hollowllm}, and the \emph{source-integrity gap} for predicate proofs is
named by zk-attach \cite{zkattach}. And we do not claim the separation theorem
that our designated-verifier interface makes tempting (``public verification is
impossible when $k$ is small''); Remark~\ref{rem:designated} gives the reason,
and Corollary~\ref{cor:reexec} is the short version of it.

\smallskip
\noindent\textbf{Contributions.}
\begin{enumerate}
  \item We isolate the gap that survives in every design we surveyed---\emph{divergent
  delivery}---and show that neither hardware nor any proof closes it, while having
  the receiver contribute to the proven statement does
  (Proposition~\ref{thm:gap}).
  \item We give a quantitative bound on what a public binding reveals, together
  with the \emph{structural hypothesis} that bound needs, and a demonstration
  that the hypothesis is necessary: our own two public values sit on opposite
  sides of it (Theorem~\ref{thm:t1}, Remark~\ref{rem:betanecessary}). The
  governing quantity $k$ is measurable rather than assumed.
  \item We price the fix. A certificate is replayable without limit unless the
  receiver contributes fresh randomness, and with an empty nonce the adversary's
  advantage is $1$ regardless of every other parameter of the system
  (Theorem~\ref{thm:t3}).
  \item We describe PoP, a system that realises this design without a TEE
  (\S\ref{sec:system}), and we report what its certificates leak about the policy
  they certify (\S\ref{sec:private-policy}).
\end{enumerate}

\smallskip
\noindent\textbf{One spine, two sides.} Theorems~\ref{thm:t1} and~\ref{thm:t3}
state the same fact about two different inputs: \emph{a proof binds only as much
as the committed, adversary-non-selectable quantity has bits for}. The response
side is governed by entropy, $k = -\log_2 \mathsf{Guess}(T \mid Z)$
(Definition~\ref{def:k}); the request side by freshness, $f$. Neither can be
bought with more computation.

\fixme{全文禁止出现"比 Proof-of-Guardrail 更强"或任何强弱对比。
Proof-of-Guardrail 与本文正交：它证"计算发生过"，本文证"性质成立"。}

% =====================================================================
\section{System and Threat Model}
\label{sec:system}

The object of study is a \emph{policy-proof system}: a pipeline that compiles a
policy written in a small DSL into a circuit, evaluates an agent's response
against that circuit inside a zkVM, and emits a certificate binding the verdict
to the response. This section fixes the object and its adversary before any
theorem is stated.

\subsection{The system}
\label{sec:system-overview}

\paragraph{Pipeline and the cross-layer contract.} A policy package is compiled
to a \texttt{ConstraintSpec}---a JSON artefact carrying the compiled NFA---which
is the \emph{only} object crossing the layer boundary: the Python reference
implementation and the Rust/zkVM guest consume the same bytes, so cross-layer
drift is structurally impossible rather than merely tested against. The guest
runs a deterministic evaluator and commits its outcome. Of eight rule classes,
seven are decided \emph{inside} the circuit
(\texttt{keyword\_\allowbreak block},
\texttt{normalized\_\allowbreak keyword\_\allowbreak block},
\texttt{length\_\allowbreak bound}, \texttt{pattern\_\allowbreak block},
\texttt{format\_\allowbreak check}, \texttt{tool\_\allowbreak arg\_\allowbreak guard},
\texttt{budget\_\allowbreak bound}). The eighth,
\texttt{semantic\_bound}---a rule of the form ``harm probability $\le$
threshold''---is precisely the kind for which no formal proof can be learned, and
the circuit does not decide it. The circuit registers only \emph{that the rule
was delegated}, in the public value \texttt{delegated}, and a verifier must
\emph{conjunctively} check a companion \texttt{ezkl}/halo2 proof. Verifying the
zkVM side alone does not establish compliance, and we say so here because it is
the most likely misreading of the certificate.

\paragraph{What the verdict is bound to.} Rules that reason about the agent's
\emph{tool calls} cannot be judged against a list the prover fills in itself.
Receipts are therefore issued by a tool gateway after each call and chained
(\texttt{seq}, \texttt{prev}, signature); the circuit checks the chain's
\emph{structure}---entries consecutive, mutually consistent, each digest
recomputed from its own content---and fails closed to \texttt{trace\_unbound}
otherwise. Signature verification happens \emph{off circuit} (Ed25519 inside the
zkVM is prohibitively expensive), a boundary we state rather than gloss: the
circuit is blind to a forgery in the chain's final element, and two independent
off-circuit gates---digest comparison and signature verification---carry that
obligation, neither substituting for the other. Truncating the chain's tail is a
separate and weaker opening, closed by a session-end \texttt{ToolSeal}
commitment from the gateway, and it remains conditional on the verifier holding
the gateway's public key. \S\ref{sec:discussion} develops these residual
boundaries; they are not incidental to the design.

\paragraph{Dual privacy modes.} In \emph{transparent} mode the response and the
verdict are public. In \emph{private} mode the certificate publishes a response
commitment, a commitment per violation, and a redaction proof
(\texttt{redaction\_ok} together with \texttt{mask\_covered}, the latter
requiring the masked positions to fall inside a span the circuit verified to be a
genuine complete match). The guarantee is that \emph{public values contain no
plaintext}---not that $T$ cannot be recovered. A commitment is a public function
of $T$ against which anyone can guess and check, while the response binding must
be publicly recomputable for a verifier to check the string it was delivered.
Those two requirements are in direct tension, and \S\ref{sec:binding} quantifies
exactly when.

\paragraph{Certificates and anchoring.} A certificate is signed with a DSSE
envelope under Ed25519 with the key held by the issuing party, and its digest is
appended to a hash-chain anchor ledger. Two backends share one interface: a file
ledger, verifiable offline, and an on-chain registration storing only a 32-byte
digest. The transaction hash and on-chain timestamp are written back into the
local entry, so the offline chain and the public timestamp cross-check each
other and tampering with either is detectable.

\paragraph{Framework integration.} One gateway serves LangChain callbacks,
LangGraph (\texttt{guard\_node} plus full event streams), and MCP tool argument
and response checks, so trace certificates and proofs share a single chain of
evidence. The implementation passes 801 tests with 15 skipped; the skipped set is
enumerated in Appendix~\ref{app:repro}.

\begin{definition}[Policy-proof system]
\label{def:system}
A policy-proof system is a tuple $(\mathsf{Prove}, \mathsf{Ver}, \pi, \mathcal{M})$.
On challenge $n$ from a verifier, $\mathsf{Prove}(\pi, T, n)$ outputs a
certificate $c$ whose public values include
\begin{multline*}
b \;=\; H(\texttt{pop-bind-v1} \,\|\, \mathsf{u32be}(|n|) \,\|\, n \,\|\, T),\\
C = H(T).
\end{multline*}
$\mathsf{Ver}(c, n, T')$ accepts iff the policy relation holds \emph{and}
$b$ recomputed from the received $T'$ matches.
\end{definition}

\fixme{\textbf{H3/H4}：早先写的 ``the relation contains no term referring to the
delivered bytes'' 与仓库事实冲突——公开值里\emph{已有} $b$，验证时由送达的
$T'$ 现场重算（\texttt{circuits/types/src/lib.rs:444}、
\texttt{docs/security-model.md:232}）。命题已按上面的定义改写。见
\texttt{docs/tdsc-fixlist.md} \S H。}

\subsection{Threat model}
\label{sec:threat}

The adversary is the \emph{agent operator}: the party that holds the policy, runs
the model, and produces the certificate. The verifier is a regulator, auditor, or
third party holding only public values. The policy author issues $\pi$; in this
section we treat it as honest, and \S\ref{sec:wp4-tension} makes that assumption
explicit and scopes it, because in policy-private mode its honesty is doing real
work and cannot be waved through. We state the goals, and more importantly the
non-goals, before the theorems.

\begin{assumption}[Adversary side information]
\label{as:Z}
The adversary knows the prompt, the model (possibly its weights), and the
decoding parameters. \emph{This is the assumption that makes the theorem
non-trivial}: under greedy decoding, $k \approx 0$.
\end{assumption}

\paragraph{Goals.}
\begin{description}
  \item[(G1)] \emph{Soundness}---a response $T$ violating $\pi$ cannot yield an
  accepted certificate.
  \item[(G2)] \emph{Committed privacy}---in private mode the public values
  contain no plaintext of $T$. This is explicitly \emph{not} the claim that $T$
  is unrecoverable: once the guessing entropy $k$ of Definition~\ref{def:k} falls
  below $\log q + \log(1/\varepsilon)$, a publicly recomputable binding and
  content hiding are mutually exclusive (Theorem~\ref{thm:t1}).
  \item[(G3)] \emph{Traceable binding}---policy version, circuit, and proof
  artefact are bound to one another and recomputed at verification.
  \item[(G4)] \emph{Selective disclosure}---redaction and evidence opening are
  themselves provable.
  \item[(G5)] \emph{Record integrity}---tampering with, or reordering, the audit
  ledger is detectable.
\end{description}

\paragraph{Non-goals.} We do not prove the model's inference, and we do not cover
training data or model cards. (G1) is modulo the soundness of the zkVM and of the
hash function; we use those primitives and do not re-derive them.

\paragraph{Where the policy-private case is developed.} The three-party model,
the issuing prerequisite (P-PP), and the table of what each verifier may and may
not conclude are specific to policy-private mode, where the policy itself is
secret. They are stated once, in \S\ref{sec:wp4-tension}, and are deliberately
\emph{not} duplicated here---at a twelve-page limit they cannot be. The general
adversary of this section is the same operator, viewed without the
policy-secrecy refinement.

% =====================================================================
\section{Binding of Policy Proofs}
\label{sec:binding}

A policy proof is meant to bind a response: a verifier who accepts a
certificate should conclude that \emph{this} string was checked against the
policy. This section asks what that binding is actually worth. We first isolate
the gap that survives in every design we surveyed (\S\ref{sec:gap}), then give
a quantitative bound on what a public value reveals, together with the
structural hypothesis that bound needs (\S\ref{sec:tradeoff}), and finally
state the same fact on the request side, where the governing resource is
freshness rather than entropy (\S\ref{sec:t3}).

\subsection{The gap}
\label{sec:gap}

\begin{proposition}[Divergent delivery]
\label{thm:gap}
Definition \ref{def:system} binds $T'$ \emph{as received by the verifying party}. If the
provider delivers a different string $T'' \ne T'$ to the end user, verification
still succeeds on $T'$. Neither hardware nor a proof can prevent this; only the
end user holding $n$ and recomputing $b$ can.
\end{proposition}
\begin{proof}
$b$ is a function of $(n,T)$ only; the delivery channel to the end user carries
no authenticated term. One line.
\end{proof}

Two remarks scope this proposition. First, it describes a \emph{residual} gap,
not the claim that policy proofs bind nothing. Our own system already binds $T$
through the verifier-issued nonce $n$ of Definition \ref{def:system}, and a
verifier who receives $T'$ and recomputes $b$ does learn that $T'$ was the
proved string; the gap is that the binding is anchored to the
\emph{verifying party's} copy. Second, we do not claim to be the first to
observe that a proof can fail to bind the artifact it purports to describe. The
\emph{effort gap} for zero-knowledge inference proofs is formalized by
Hollow-LLM \cite{hollowllm}, and the \emph{source-integrity gap} for predicate
proofs is named by zk-attach \cite{zkattach}; both concern what the proof system
was asked to prove, whereas Proposition \ref{thm:gap} concerns where the proved
string is delivered. We state the naive version of our own claim because it
would be false against our own system.

\subsection{A quantitative binding--hiding tradeoff}
\label{sec:tradeoff}

The certificate must publish \emph{something} tying it to $T$; otherwise it
binds nothing at all. We therefore ask how much that public value leaks. The
answer depends on one measurable quantity --- the guessing entropy of the
response given everything the adversary already knows --- and on one structural
property of the public value itself.

\begin{definition}[Guessing entropy of a response]
\label{def:k}
For a response distribution $\mathcal{D}$ and adversary side information $Z$,
let
\[
\mathsf{Guess}(T \mid Z) \;\triangleq\; \max_{\mathcal{A}} \Pr[\mathcal{A}(Z) = T],
\]
which equals $\mathbb{E}_z[\max_t \Pr[T = t \mid Z = z]]$, and
\[
k \;\triangleq\; -\log_2 \mathsf{Guess}(T \mid Z),
\]
the maximum ranging over all time-unbounded adversaries.
\end{definition}

\begin{remark}[The definition is operational, and it is the average]
\label{rem:kmeas}
$k$ is \emph{measured}, not assumed: sampling one model on one prompt many
times estimates $\mathsf{Guess}(T \mid Z)$ directly
(\S\ref{sec:eval}), so Theorem \ref{thm:t1} is instantiated with a
number we obtain rather than a number we posit. The definition must average
over $z$ rather than take the pointwise worst case
$k_z = -\log \max_t \Pr[T = t \mid Z = z]$: the proof takes an expectation over
$z$, and $\mathbb{E}_z[2^{-k_z}] = 2^{-k}$ holds only for the average form.
\end{remark}

\begin{theorem}[Binding, and the hypothesis it needs]
\label{thm:t1}
Let $\beta : \mathcal{M} \to \{0,1\}^\lambda$ be public, deterministic and
efficiently computable, and let the certificate publish $\beta(T)$. Let
$\mathcal{A}$ be any adversary, unbounded in time, that is given $Z$ and
evaluates $\beta$ at most $q$ times.
\begin{enumerate}
\item[(H)] \emph{(Randomised public value.)} Suppose
  $\beta(T) = H(\mathsf{dom} \,\|\, \mathsf{u32be}(|n|) \,\|\, n \,\|\, T)$ for a
  random oracle $H$ and a nonce $n$ drawn uniformly and \emph{independently of
  $T$}. Then
  \[
  \Pr\!\left[\mathcal{A}^\beta\!\left(Z, \beta(T)\right) = T\right]
  \;\le\; q \cdot 2^{-k} \;+\; q \cdot 2^{-\lambda}.
  \]
  Consequently hiding with advantage $\varepsilon$ requires
  $k \ge \log q + \log(1/\varepsilon)$.
\item[(L)] \emph{(Lossy public value.)} For an arbitrary deterministic public
  $\beta$ the correct bound is over the \emph{preimage class}:
  \begin{multline*}
  \Pr\!\left[\mathcal{A}^\beta\!\left(Z, \beta(T)\right) = T\right]
  \;\le\; \mathsf{Guess}\!\left(T \mid Z, \beta(T)\right) \\
  \triangleq\; \mathbb{E}_y\!\left[\max_t \Pr[T = t \mid Z, \beta(T) = y]\right].
  \end{multline*}
\end{enumerate}
\end{theorem}

\begin{proof}[Proof sketch]
Part (L) is immediate after conditioning on $\beta(T)$; no hypothesis on
$\beta$ is used. For (H), fix $z$ and decompose the adversary's success by the
first query that hits $T$. The oracle answers it has already seen are uniform
and independent of $T$, so each query is itself independent of $T$ and succeeds
with probability at most $2^{-k_z}$. Taking the expectation over $z$ and
applying Definition \ref{def:k} yields $q \cdot 2^{-k}$; the term
$q \cdot 2^{-\lambda}$ charges a query that hits the uniform value $\beta(T)$
by chance. Appendix~\ref{app:proofs} gives the full argument.
\end{proof}

\fixme{$\lambda$ 比特随机与随机预言机只给 \emph{heuristic} 保证：实践中的 SHA-256
不是真的随机预言机（骨架 :712--715）。正文措辞必须写成 heuristic，
\textbf{不得}写成信息论界。}

\begin{remark}[The hypothesis in (H) is necessary, and our two instances straddle it]
\label{rem:betanecessary}
Without the hypothesis, part (H) is not merely unproven but \emph{false}: take
$\beta = \mathrm{id}$, or $\beta = \mathsf{encode}$ as in
Proposition \ref{prop:semantic}. The adversary outputs $\beta(T)$ and succeeds
with probability $1$, independently of $q$ and $k$. Our own two public values
therefore sit on opposite sides of the hypothesis---$b =
H(\mathsf{dom}\|n\|T)$ satisfies it, $\mathsf{encode}(T)$ does not---and this is
what makes the pair a theorem about \emph{our design} rather than a citation of
a standard bound. An earlier draft of this work stated (H) without its
hypothesis; the counterexample above is why the hypothesis is now explicit.
\end{remark}

\begin{corollary}[Re-execution, not enumeration]
\label{cor:reexec}
Under Assumption \ref{as:Z}, when the decoding is greedy or low-temperature,
$k \approx 0$ and the attack of Theorem \ref{thm:t1} is realised by
\emph{re-running the model}, at the cost of one inference---independent of
$|T|$ and requiring \emph{no} evaluation of $\beta$. Enumeration over the string
space is \emph{not} the relevant attack.
\end{corollary}

\fixme{原 docs/sp1-zk-audit.md §4 写的是"熵(~10^2--10^3 bit)远低于 256 bit，
因此枚举可行"——这是非因果推论（2^100 与 2^256 一样枚举不了；真正的攻击是\textbf{重执行}，
对手的已知信息是 (prompt, 模型, 解码参数)）。
2026-10-12 已全库改正，docs/tdsc-fixlist.md §G 的 G1--G6 逐条结案。}

\begin{remark}[A design observation, not a separation]
\label{rem:designated}
Our verification interface is \emph{designated}: binding is checked against a
value $b = H(\mathsf{dom}\|n\|T)$ in which a verifier-issued nonce $n$ appears.
It is tempting to promote this into a separation (``public verification is
impossible when $k$ is small''). \textbf{We claim no such theorem, and
Corollary \ref{cor:reexec} is why.} When $k \approx 0$, hiding fails for
\emph{every} scheme, designated or not, because the adversary re-runs the
model; when $k$ is large, a zero-knowledge proof of ``$\exists\, T : b =
H(\mathsf{dom}\|n\|T) \wedge J(\pi, T) = 1$'' is simultaneously publicly
verifiable and hiding. The dichotomy is thus governed by $k$ alone and carries
no independent content. What does survive---and what \S\ref{sec:discussion}
develops---is that \emph{third-party auditing of delivery} forces disclosure of
$T$, which is an informational remark rather than a theorem. We record the
demotion here because an earlier draft asserted the separation as a theorem.
\end{remark}

We stress that the designated-verifier shape is a \emph{retrospective
justification, not a novel construction}. The system publishes
$b = H(\mathsf{dom}\|n\|T)$ with a verifier-issued nonce $n$, and
Remark \ref{rem:designated} explains the shape of that choice after the fact.
The security model does not pin down whether $n$ is public: making it public
forfeits hiding, while holding it at the verifier is what makes verification
designated.

\begin{remark}[A hypothesis that can be settled by experiment]
\label{rem:kexp}
Theorem \ref{thm:t1} is only as useful as the number put into it. Because
Definition \ref{def:k} is operational, $k$ can be measured rather than assumed:
sampling one model on one prompt many times estimates
$\mathsf{Guess}(T \mid Z)$ directly, and the estimate substitutes into the
threshold $k \ge \log q + \log(1/\varepsilon)$. To our knowledge no prior work
measures this quantity for agent responses, which makes the measurement a
contribution in its own right; \S\ref{sec:eval} describes the protocol and
reports the result.
\end{remark}

\begin{proposition}[The lossy instance of Theorem \ref{thm:t1}(L)]
\label{prop:semantic}
For the semantic rule class, the public values contain $\beta(T) =
\mathsf{encode}(T)$, a content-preserving feature vector of the response. This
$\beta$ is deterministic and injective on the features but not on strings, so
it violates the hypothesis of Theorem \ref{thm:t1}(H) and instantiates part
(L): the binding force is exactly
$\mathsf{Guess}(T \mid Z, \mathsf{encode}(T))$, i.e.\ the guessing probability
\emph{within the preimage class} of the published vector. Here $k$ need not be
argued from perplexity at all: the adversary's search space is that preimage,
which is dramatically smaller than the string space. The relevant asymmetry is
not ``low entropy'' but a \emph{lossy public encoding}.
\end{proposition}
\begin{proof}
Immediate. The measurement is direct: publish $\mathsf{encode}(T)$ for real
responses and measure how much of $T$ is recoverable.
\end{proof}

Proposition \ref{prop:semantic} is the cheapest experiment in this paper and
the strongest anchor for Theorem \ref{thm:t1}: it is measurable in a day, and it
makes the bound concrete rather than cited. It also shows that the relevant
asymmetry is not entropy at all --- for the semantic rule class there is no
entropy argument available, because $\mathsf{encode}(T)$ simply \emph{is} a
content-preserving function of the response.

\subsection{Quantitative request binding}
\label{sec:t3}

\begin{theorem}[Selective-presentation advantage]
\label{thm:t3}
Fix a policy $\pi$ and a request distribution. A selective-presentation
adversary $\mathcal{B}$ may, \emph{before} seeing the receiver's challenge,
compute a set $S$ of accepted transcripts. The receiver then contributes $f$
fresh, uniformly random bits $r$ that enter the proven statement. Then
\[
\Pr\!\left[\mathcal{B}\ \text{is accepted on } r\right] \;\le\; |S| \cdot 2^{-f}.
\]
If instead the receiver contributes no fresh randomness ($f = 0$; e.g.\ an
empty nonce), the advantage is $1$, and \emph{does not decay in any parameter
of the system}.
\end{theorem}
\begin{proof}
The receiver's freshness $r$ enters the proven statement through
$b = H(\mathsf{dom} \,\|\, r \,\|\, T)$. For a certificate in $S$ to be accepted
on $r$, the value $r$ must fall inside the challenge set that certificate
already covers --- an event of probability at most $2^{-f}$ for $f$ fresh
uniform bits --- and a union bound over $S$ gives $|S| \cdot 2^{-f}$. When
$f = 0$ every element of $S$ is a permanently valid replay, so the advantage is
$1$ and does not decay in any parameter of the system.

The $f = 0$ branch is not hypothetical.
\texttt{circuits/\allowbreak types/\allowbreak src/\allowbreak lib.rs}
(lines 950--995) documents an empty nonce as a legal input that ``provides no
replay protection''; it is therefore a \emph{valid input to the current API},
not a straw man.
\end{proof}

Theorem \ref{thm:t3} is \emph{isomorphic} to Theorem \ref{thm:t1}, and that is
what earns it a place here. Both state the same fact about different inputs: a
proof binds only as much as the committed quantity the adversary cannot select
in advance has bits for. On the response side that quantity is the response and
the resource is entropy $k$; on the request side it is freshness $f$. Neither
can be bought with more computation. This is the spine announced in
\S\ref{sec:intro}, and the two theorems are its two sides.

\begin{remark}[A taxonomy, not a trichotomy]
\label{rem:taxonomy}
Every policy-proof deployment we surveyed closes the divergent-delivery gap of
Proposition \ref{thm:gap} in one of three ways:
(a) root trust in hardware (attested enclaves);
(b) have the \emph{end user} contribute a secret to the proven statement (the
user issues the nonce); or
(c) gate delivery on successful verification (\emph{prove-before-deliver}).
\textbf{We offer this as a classification, not as an exhaustiveness theorem}:
we neither prove that a fourth family cannot exist, nor that the three are
disjoint. zk-attach \cite{zkattach} walks (a); NiyamAI and ZKAP walk (c); PoP
walks (b).
\end{remark}

\paragraph{Posture, stated plainly.} Our system already carries
\texttt{response\_binding}, but $n$ is issued by the \emph{verifier}. Moving it
to the \emph{end user}, so that $f > 0$ becomes the receiver's contribution,
closes the divergent-delivery branch of Proposition \ref{thm:gap}. We are
explicit that this is a small protocol amendment rather than a new system: the
mechanism is already implemented, and only the issuer of the nonce changes. What
Theorem \ref{thm:t3} buys is the \emph{price}: with $f = 0$ there is no closure
at any cost.

% =====================================================================
\section{Policy-Private Compliance}
\label{sec:private-policy}

For a policy whose rules are deterministic and a response that is public
anyway, a signed hash chain over the response suffices: the signature already
binds the bytes, and anyone can re-evaluate the rules. The one scenario a
signature cannot cover is a policy that must itself stay secret---where the
operator has to certify compliance to a third party without handing over the
rules it is being judged against. That is the mode this section is about, and we
lead with what the mode \emph{cannot} buy, because that determines whether it
means anything at all.

\subsection{The observation that reframes this work package}
\label{sec:wp4-finding}

\texttt{spec\_canonical} is already a \emph{guest input}---a witness---rather
than a public value; only its hash,
\[
\mathsf{policy\_hash} \;=\; \mathsf{SHA256}(\texttt{spec\_canonical}),
\]
enters the public values (\texttt{circuits/\allowbreak types/\allowbreak
src/\allowbreak lib.rs}:1--16, 415--416, 918, 1003--1004). The instruction
``do not reveal the policy text'' is therefore
\emph{already true of the architecture}; what leaks is everything \emph{derived}
from the policy. Policy privacy is thus not a new proof system to be built but a
\emph{disclosure policy to be audited}, and we write it that way---as a
systematic account of what a certificate reveals about the policy it certifies.

\begin{table}[t]
\centering
\caption{What a policy-private certificate leaks today, and what closing each
leak costs.}
\label{tab:wp4}
\tabfont
\begin{tabular}{@{}L{0.04\linewidth}L{0.46\linewidth}L{0.36\linewidth}@{}}
\toprule
\# & Leak & Difficulty / dependency \\
\midrule
1 & The policy text itself & \textbf{Already closed} (it is a witness). \\
2 & Witness hiding in the proof system & Medium; \emph{inherited from the
    proving tier}: only \texttt{groth16}/\texttt{plonk} is witness-hiding, per
    \texttt{docs/\allowbreak sp1-zk-audit.md}. \\
3 & \texttt{PrivateViolation}\allowbreak\{\texttt{rule},\allowbreak\texttt{kind}\}
    names the rule that fired & Medium; needs a coarsened violation class plus a
    commitment to it. \\
4 & Evidence opening reveals the rule
    (\texttt{policydsl/\allowbreak privacy/\allowbreak commit.py}:319) & Medium;
    disclose the \emph{existence} of a violation, not its identity. \\
5 & The \texttt{delegated} flag exposes policy structure (which rules were
    routed to the \texttt{ezkl} path) & Medium; easiest to overlook, cheapest to
    fix. \\
6 & Semantic rules publish $\beta(T)=\mathsf{encode}(T)$
    (Proposition~\ref{prop:semantic}) & Medium--high; either bound the
    preimage, or disable this rule class in policy-private mode. \\
\bottomrule
\end{tabular}
\end{table}

The honest structure of Table~\ref{tab:wp4} is worth stating, because it sets
the reader's expectations for the rest of the section. Sub-problem~1 is a
one-paragraph re-reading of existing code, not a contribution. Sub-problem~2 is
inherited wholesale from the proving tier and is not specific to this work.
Sub-problems~3--5 are engineering with a single formal core---the granularity of
disclosure. Sub-problem~6 is where the work actually is, and it collides with
\S\ref{sec:tradeoff}: one cannot hide the policy while publishing the lossy
encoding of the response that the semantic rule class requires
(Proposition~\ref{prop:semantic}).

\subsection{The tension worth writing up: policy privacy vs.\ accountability}
\label{sec:wp4-tension}

\paragraph{The prerequisite, stated as a prerequisite.} Let
$H = \mathsf{SHA256}(\pi)$. On its own $H$ is \emph{inert}. A prover can always
emit a self-consistent certificate for any policy $\pi'$ it prefers, simply by
placing $H' = \mathsf{SHA256}(\pi')$ in that certificate. \emph{Secrecy of
$\pi$ is what creates the gap}, and no cryptographic mechanism can close it:
closing it would require showing $\pi$ to the verifier, which is precisely what
the mode exists to avoid. What closes it is an \emph{issuing act}, so we state
it as an explicit prerequisite of the mode.

\begin{description}
  \item[(P-PP)] An \emph{issuing authority} signs $\pi$ and registers $H$ in a
  record readable by third parties. A policy is \emph{accountable} iff its
  \texttt{policy\_hash} can be traced to such a record.
\end{description}

Without (P-PP) the mode is not merely weaker---the verifier's conclusion is
empty. This is the same assumption as the ``honest policy author'' of
\texttt{docs/security-model.md}~\S0.2; what is new is scoping it to the
verification interface, which had not been done, and which a referee would
otherwise use to dismiss the whole mode in one sentence.

\paragraph{Three parties.} ``Three'' means \emph{authority / operator / third
party}; an auditor is a named instance of the third party, not a fourth trust
anchor.

\begin{table}[t]
\centering
\caption{The three-party model of policy-private mode.}
\label{tab:wp4-parties}
\tabfont
\setlength{\tabcolsep}{3pt}
\begin{tabular}{@{}L{0.16\columnwidth}L{0.14\columnwidth}L{0.62\columnwidth}@{}}
\toprule
Party & Honest? & Holds; wants \\
\midrule
Authority $R$ & \textbf{Yes}, by (P-PP) & Holds $\pi$ and the issuing record;
wants $\pi$ actually enforced, and wants to be able to check that it was. \\
Operator $O$ & \textbf{No} (the adversary) & Holds $\pi$, $T$, the proving key;
wants $\pi$ hidden from third parties while still being able to certify. \\
Third party $V$ & Yes & Holds only public values---\texttt{policy\_\allowbreak hash},
\texttt{circuit\_\allowbreak hash}, \texttt{response\_\allowbreak commitment},
\texttt{passed}, the proof; wants to confirm that \emph{some issued} policy
passed, and must \emph{not} learn $\pi$. \\
\emph{(Auditor $V^{*}$)} & Yes & As $V$, plus $\pi$ once $O$ discloses it;
wants ex-post accountability. \\
\bottomrule
\end{tabular}
\end{table}

Two boundaries must be stated with the model. First, secrecy of $\pi$ comes from
a \emph{data-flow fact}---$\pi$ never enters the public values---and not from
zero-knowledge; any rule that writes a fragment of $\pi$ into the public values
(the leaks in Table~\ref{tab:wp4}, items~3 and~5) breaks secrecy on that
fragment. Second, accountability for $V^{*}$ is \emph{voluntary}: $V^{*}$ can
check only when $O$ chooses to disclose $\pi$. The mode offers no forced
disclosure; that would require the release-gating family discussed in
\texttt{docs/tdsc-plan.md}~\S1.6(b), which we do not claim.

\paragraph{What each verifier may conclude.} Table~\ref{tab:wp4-capability} is
the contribution of this subsection: the cryptography is downstream of it, and
the table survives even if the cryptography only partially lands.

\begin{table*}[t]
\centering
\caption{What each verifier may and may not conclude in policy-private mode.}
\label{tab:wp4-capability}
\tabfont
\setlength{\tabcolsep}{4pt}
\begin{tabular}{@{}L{0.13\textwidth}L{0.30\textwidth}L{0.30\textwidth}L{0.19\textwidth}@{}}
\toprule
Verifier & \textbf{May} conclude & \textbf{May not} conclude & Rests on \\
\midrule
Third party $V$ & (i) the commitment binds $T$
(Definition~\ref{def:system}); (ii) \emph{some} policy with hash $H$ passed;
(iii) $H$ appears in the issuing record & (iv) the \emph{content} of $\pi$;
(v) how many rules $\pi$ has, or of which classes; (vi) ``the operator enforces
the regulation \emph{I} care about''---only ``some issued regulation'' &
(P-PP); Assumption~\ref{as:Z}; the $k$ of Theorem~\ref{thm:t1} \\
End user (also receiver) & As $V$, and additionally: the proved $T$ is the
string it received, since its nonce is bound (Theorem~\ref{thm:t3}, requires
$f>0$) & Only the string it received; \emph{not} what was delivered to another
downstream agent & $f>0$ \\
Authority $R$ & Everything; it holds $\pi$ and can recompute $H$ & --- & --- \\
Auditor $V^{*}$ (after disclosure) & (vii) $\mathsf{SHA256}(\pi) =
\texttt{policy\_\allowbreak hash}$, i.e.\ what ran was the issued $\pi$;
(viii) the disclosed violation record is consistent & (ix) the undisclosed part
of $T$ & voluntary disclosure \\
On-chain verifier & (i)--(iii) only; it hashes and verifies signatures, and does
not interpret policy semantics & (iv)--(ix) & (P-PP) \\
\bottomrule
\end{tabular}
\end{table*}

The $V$ row is the reason this subsection exists. A third party does not obtain
``compliance''; it obtains ``some issued policy was satisfied''. Conflating the
two is the sentence a referee would use to reject the mode, and preventing that
sentence is what the table is for.

\paragraph{Prior art.} \emph{Policy hiding is not a new primitive}, and we claim
no primacy for it. Hidden credentials \cite{hiddencred}, secret handshakes
\cite{handshake}, privacy-preserving access control over XACML policies
\cite{xacmlpriv}, and attribute-hiding and fuzzy predicate encryption
\cite{sw05,ksw08} are one family, and policy-private mode is an instance of it.
What is new here is narrow: (a) the \emph{application} to LLM output compliance,
together with selective disclosure of violation type and location, and (b) the
accountability tension of this subsection.

\need{⚠️ 骨架 §4.3「Expected results, stated in advance」（子问题置信度百分比表）
\emph{未}迁入正文。那是 WP 规划的产物，不是论文内容——期刊稿里出现
「我们对子问题 4 有 70\% 把握」是硬伤。若作者要保留这份校准，应改写为
\S\ref{sec:discussion} 里一段关于适用范围与未完成项的讨论，而不是概率表。}

\need{(P-PP) 的落地形式尚未定：链上已签发策略哈希注册表，还是签名清单？
属作者决策。先例引文已按骨架 :623--630 的要求入 bibliography，
但 \cite{hiddencred} / \cite{handshake} / \cite{xacmlpriv} / \cite{sw05} /
\cite{ksw08} 五条\emph{均待逐条打开原文核实}（外部资源阻塞）。}

% =====================================================================
\section{Evaluation}
\label{sec:eval}

\subsection{Cycle count: length versus rule count}
\label{sec:eval-cycles}

Response length was pushed from 200 to 100{,}000 characters and the rule count
from 1 to 6, giving 40 measurement points (\texttt{bench/bench\_cycles.py}; the
full matrix is in \texttt{bench/results/cycles.md}). The response text comes from
an evaluation artefact the repository's own end-to-end demo produced (372
characters, three segments deduplicated by content); longer points are tiled from
it. What is measured is therefore the \emph{curve} of scanning cost, not ``a real
100{,}000-character document''. Every point was first confirmed to match no rule:
a matching rule short-circuits, and the cost of a short-circuit says nothing.

\begin{table}[t]
\centering
\caption{Cycle count against length, per rule set. Five rule sets, six length
points each; $R^2 \approx 1.0000$ throughout.}
\label{tab:cycles}
\tabfont
\begin{tabular}{@{}rrrrr@{}}
\toprule
Rules & \texttt{keyword} & \texttt{pattern} & Slope (cyc/char) & Intercept \\
\midrule
1 & 0 & 0 & 73.9      & 50{,}137 \\
2 & 1 & 0 & 105.3     & 65{,}942 \\
3 & 1 & 1 & 4{,}289.8 & 387{,}501 \\
4 & 2 & 1 & 4{,}048.8 & 405{,}785 \\
6 & 3 & 2 & 5{,}403.5 & 993{,}257 \\
\bottomrule
\end{tabular}
\end{table}

Three findings, in decreasing order of how much they change what may be claimed.

\emph{(1) With the rule set fixed, cycle count is exactly linear in length.}
Five rule sets, six length points each, $R^2 \approx 1.0000$ throughout
(Table~\ref{tab:cycles}).

\emph{(2) Scanning cost is dominated by the NFA, not by the rule count.} One
keyword rule costs about $+31.6$ cycles per character at the margin; one NFA
regular expression (an email pattern) costs about $+3{,}634$---roughly
$115\times$ more. Equivalently, a rule set containing a regular expression has a
per-character slope ($4{,}289.8$) $58\times$ that of a rule set with only string
rules ($73.9$).

\emph{(3) Unit prices are not additive: a unit price is a property of a rule
\emph{in context}, not of the rule.} Fitting the slopes as
$(\text{keyword},\text{pattern})$ unit prices gives $R^2 = 0.9169$ and a
\emph{negative} keyword unit price ($-477.2$ cycles per character). This is not
noise; two probes give the mechanism directly:

\begin{itemize}
  \item \textbf{Ordering probe} (same rules, same text, $L = 10^5$): declared
  order costs $541{,}342{,}069$ cycles against $570{,}736{,}989$ in reverse,
  a difference of $+5.43\%$, while \texttt{passed} is identical in both. \emph{The
  semantics of the decision do not depend on order; the cost does.}
  \item \textbf{Marginal probe} ($L = 10^5$, filling the one matrix cell the
  sweep could not reach---regular expressions with no keyword rule):
  $[\text{len},\text{pat}]$ over $[\text{len}]$ has a marginal unit price of
  $3{,}633.7$, whereas ``add one more keyword rule'' has two different answers:
  $+31.6$ alone, $+585.6$ in a policy that already contains a regular expression
  ($-153.7$ if inserted before it). The same action changes sign between
  $-154$ and $+586$.
\end{itemize}

The mechanism lies in zkVM memory allocation: every \texttt{ascii\_lower} and
every Pike VM match allocates, and allocator state depends on previous
allocations. We therefore quote \emph{per-rule-set slopes} (the table above,
directly citable) rather than any ``unit price $\times$ count'' expression. The
simple model $\text{cycles} \approx a\cdot|T|\cdot\text{rules} + b$ that this
work anticipated is refuted by the data ($R^2 = 0.872$, negative intercept), and
we report that too: \emph{a model refuted by its own data is itself an evaluation
result.}

\subsection{Matcher ablation: quadratic degradation on a pathological input}
\label{sec:eval-ablation}

\emph{Construction.} Pattern \texttt{a+b} against input
\texttt{a}$^{\times n}$ containing no \texttt{b}, so the match is necessarily
absent and both implementations must scan to the end. The naive implementation
reinvents and restarts at every start position, giving $\Theta(n^2)$; the Pike VM
decides in a single linear pass (\texttt{bench/bench\_ablation.py}).

\begin{table}[t]
\centering
\caption{Matcher ablation. Both paths decide identically; only the cost differs.}
\label{tab:ablation}
\tabfont
\begin{tabular}{@{}rrrrr@{}}
\toprule
$n$ & Pike (cycles) & Naive (cycles) & Ratio & Naive growth/char \\
\midrule
100 & 643{,}894   & 12{,}382{,}720  & 19.2$\times$ & --- \\
200 & 1{,}145{,}209 & 48{,}638{,}946 & 42.5$\times$ & $+96.4\%$ \\
400 & 2{,}143{,}617 & 193{,}207{,}130 & 90.1$\times$ & $+98.6\%$ \\
\bottomrule
\end{tabular}
\end{table}

Read per-character cost rather than absolute figures. Doubling $n$ leaves the
Pike VM's per-character cost essentially unchanged ($-11.1\%$, $-6.4\%$) while
the naive path's nearly doubles ($+96.4\%$, $+98.6\%$). At $n = 400$ the naive
path is already $90\times$ slower, and that ratio doubles again with each
doubling of $n$. \emph{This is an attack surface}: the input length for policy
compliance is chosen externally---an agent's reply, a concatenated session---so
an adversary who needs only to write a long reply drives a naive implementation's
cost quadratically. The comparison is meaningful only because the two paths
decide identically, which \texttt{tests/test\_ablation.py} pins down (Python and
Rust \texttt{match\_search} $\equiv$ \texttt{match\_search\_naive}). To our
knowledge this is the first ablation evidence specific to policy-compliance
matching.

\subsection{Real proofs: time, size, memory---and the ceiling}
\label{sec:eval-proofs}

Sampling points share the corpus of \S\ref{sec:eval-cycles}. Each point runs in
its own subprocess (SP1's memory accumulates within a process, so two proofs in
one process do not price one proof), and failures are reported in the table
rather than omitted. Rows marked OOM are a \emph{result}, not a malfunction: they
measure where this machine cannot go.

\begin{table}[t]
\centering
\caption{Real proofs on a 24-core / 12\,GB CPU machine. ``Config'' is
characters $\times$ rules; $k$ abbreviates $10^3$. Every succeeded
point is at the \texttt{core} (non-zero-knowledge) tier.}
\label{tab:proofs}
\tabfont
\begin{tabular}{@{}lrrrl@{}}
\toprule
Config & Time & Size & Peak RSS & Result \\
\midrule
$200 \times 1$  & 123.5\,s & 2716.4\,KiB & 10{,}389\,MB & \checkmark \\
$200 \times 2$  & 118.9\,s & 2716.9\,KiB & 10{,}438\,MB & \checkmark \\
$2k \times 1$   & 138.1\,s & 2718.6\,KiB & 10{,}449\,MB & \checkmark \\
$10k \times 1$  & 172.5\,s & 2726.6\,KiB & 10{,}506\,MB & \checkmark \\
$200 \times 3$  & ---      & ---          & ---          & OOM \\
$20k \times 1$  & ---      & ---          & ---          & OOM \\
\bottomrule
\end{tabular}
\end{table}

\begin{itemize}
  \item \textbf{A fixed floor of $\approx$10.15\,GiB.} The smallest
  configuration---200 characters, one rule---already peaks at 10{,}389\,MB: the
  prover's own overhead (the trace and permutation structures of a \texttt{core}
  proof) is almost independent of the trace. All four successes fall inside the
  narrow band 10{,}389--10{,}506\,MB. Length and rule count have little marginal
  effect on memory---\emph{but the headroom is too thin for that to help}.
  \item \textbf{A two-step staircase, not a slope.} One rule is provable up to
  10k characters; two rules at about 200 characters; \textbf{three or more do not
  come out at all}. The third rule in the sweep is a \texttt{pattern\_block}: a
  regular-expression match activates another family of AIR chips and lifts the
  trace area in one step, so the cliff is at $2 \to 3$ rules; in length it is at
  $10\text{k} \to 20\text{k}$.
  \item \textbf{Memory binds, not CPU.} The cycle sweep of
  \S\ref{sec:eval-cycles} reaches 100k characters with 6 rules because it only
  executes in the zkVM and produces no proof; at 100k characters the cycle count
  is still on the order of $10^7$, which any CPU handles. \emph{The proving-side
  ceiling and the execution-side ceiling are in different places.} ``How large a
  policy can be proved'' is, on this machine, a fixed floor plus two steps rather
  than a cost curve.
  \item \textbf{Artefact size is nearly independent of configuration}: from 200
  to 10{,}000 characters the proof grows only from 2716.4 to 2726.6\,KiB. Pure
  verification takes $\approx$90\,ms (the \texttt{vkey} setup is 1.5\,s, once).
  \item \textbf{Reproducibility, stated precisely.} The \emph{position} of the
  wall reproduces (each failing point failed twice); the \emph{timings} at the
  wall do not. Re-running one group in isolation shows $\pm 10\%$ jitter in time
  and $\pm 1\%$ in memory, so we quote time only to an order of magnitude.
  \item \textbf{Session and multiparty.} The session aggregation proof (three
  certificates) measured 152.5\,s; multiparty verification needs one proof
  \emph{per non-empty slice}, so its cost grows linearly in the number of slices
  (a two-proof example policy). Neither introduces a new proving system---a slice
  and a session are ordinary inputs to the guest---but both are subject to the
  same floor: every additional proof adds another $\approx$10\,GiB peak, so
  multiparty proving must run in separate processes, serially.
\end{itemize}

\need{⚠️ 本节的 groth16 档与 ZK/非 ZK 交叉点属 WP1，尚未打通。骨架 :665--668 的
矩阵规格为规则 $\{1,2,3,4,6,8\}$ $\times$ 长度 $\{200,2k,10k,50k\}$ $\times$
模式 $\{$transparent, private, private-policy$\}$ $\times$ 证明档 $\{$core, groth16$\}$
（\texttt{bench/bench\_scale.py}）。表 \ref{tab:proofs} 只覆盖 core 档。
groth16 列落地前不得填入任何 ZK 数字。}

\subsection{The non-ZK baseline: what a signature cannot do}
\label{sec:eval-baseline}

A signed Merkle log is the obvious alternative to a proof, and on cost it is not
close: everything is milliseconds. We measured it so that the comparison is not
rhetorical. Policy package \texttt{eu\_ai\_act\_v1.json}, 103-character response,
one challenge nonce, a 1024-leaf log; each timing is the median of $n = 500$
repetitions.

\begin{table}[t]
\centering
\caption{Signed-Merkle baseline (WP9). No proofs, no zkVM: this arm is pure
Python, which is why the figures are milliseconds.}
\label{tab:baseline}
\tabfont
\begin{tabular}{@{}lrl@{}}
\toprule
Metric & Median & Note \\
\midrule
Sign certificate (Ed25519 + DSSE) & 0.132\,ms & $n=500$ \\
Merkle append (leaf + root recompute) & 2.131\,ms & $n=500$ \\
Public verification (sig + inclusion) & 0.322\,ms & $n=500$ \\
Re-run verification ($+J(\pi,T)$) & 0.772\,ms & $n=500$ \\
Certificate size & 648\,B & \\
Inclusion proof (10 siblings) & 761\,B & \\
\bottomrule
\end{tabular}
\end{table}

\need{⚠️ 引用 2.131 ms 时必须写明它是 \texttt{proofs/session.py} 的\emph{静态树重算}
$\Theta(n)$，\emph{不是} RFC-6962 增量日志的摊销 $\Theta(\log n)$。
⚠️ 共享机器有 $\pm 20\%$ 漂移。数据源 \texttt{bench/results/baseline.md}。}

The asymmetry is the point, and it is what this arm contributes. The verification
chain is \emph{verify signature} $\to$ \emph{verify log inclusion} $\to$
\emph{re-run $J(\pi,T)$}. The first two steps need only public data; the third
needs the plaintext of both $\pi$ and $T$. That is not an implementation defect
but a structural consequence of the signature paradigm: a signature binds
\emph{who said it} to the bytes, never \emph{whether it is true}. We measured the
consequence rather than asserting it---\textbf{the same key signs a certificate
whose declared content is false, and the public path passes green}, because what
it sees is a certificate with a valid signature and a valid inclusion proof. Only
re-executing $J(\pi,T)$ recovers the truth, and taking that step away from the
verifier is exactly what the zero-knowledge design buys. The machine-checked
assertions are enumerated in \texttt{bench/results/baseline.md}; one row of that
file is worth repeating because it is the honest one: on
\emph{unforgeability} the baseline is \emph{not weaker} than the proof path, both
being Ed25519 and vkey binding, and we do not score a point for ZK there.

\subsection{End-to-end and auditing}
\label{sec:eval-e2e}

The end-to-end demo produces 14 certificates (streaming chain with early stop,
MCP argument and response checks, one zk proof) and anchors them; a third party
verifies with public artefacts alone, and the ledger chain, signature,
\texttt{policy\_hash}, anchoring, stream chains and zk proof all pass. In the
on-chain variant (a local Anvil node, \texttt{Anchor.sol} deployed, 14
certificate digests registered) the verifier reads back each certificate's
\texttt{anchoredAt} and cross-checks the local ledger, passing 14/14; the
negative control (an unregistered digest reads back as zero) confirms the check
is not vacuously true. Cross-validation submits 19 vectors to the Python
reference decision and to both the SP1 host and proving paths, comparing
\texttt{passed} and the violated rule set row by row; a full re-run with real
\texttt{core} proofs gave 19/19 on both paths.

\subsection{Measuring \texorpdfstring{$k$}{k}}
\label{sec:eval-k}

\need{对一组真实 prompt、同一模型多次采样，估 $\mathsf{Guess}(T\mid Z)$，
代入 Thm \ref{thm:t1}(H) 的阈值 $k \ge \log q + \log(1/\varepsilon)$（承接
Rem \ref{rem:kexp}）。据先例扫查无人测过这个量，本身是一块贡献；
同时它也是 Prop \ref{prop:semantic} 的测量（发布 $\mathsf{encode}(T)$ 后
$T$ 的可恢复比例）。\textbf{⚠️ 需真实模型多次采样，当前阻塞（外部资源）。}}

\need{⚠️ 「Why zero-knowledge?」决策图（骨架 :670--675）\emph{未}放入本节：
它讲的是设计结论（ZK 只在 $\{$私有响应, 私有策略$\}$ 两角有正收益），是结论不是测量，
按 docs/tdsc-writing-plan.md 的建议移至 §7 Discussion。与骨架现状不同，\textbf{待作者认可}。}

% =====================================================================
\section{Related Work}
\label{sec:related}

We position PoP against five lineages. For each, Table~\ref{tab:related} records
the slot it occupies---which is also, read the other way, what this paper does
not have to claim.

\begin{table}[t]
\centering
\caption{Positioning against the five prior-art lineages. The right-hand column
names what each lineage already occupies, and therefore what we do not claim.}
\label{tab:related}
\tabfont
\setlength{\tabcolsep}{3pt}
\begin{tabular}{@{}L{0.26\columnwidth}L{0.42\columnwidth}L{0.26\columnwidth}@{}}
\toprule
Lineage & Representative work & Slot already occupied \\
\midrule
Policy $\to$ circuit & zkGen \cite{zkgen}, Circuitree, Prezta & DSL-to-circuit
translation \\
Agent/LLM compliance & NiyamAI, Zaguk, PCO \cite{pco}, WITNESS & policy-level
private-mode kernel \\
Model computation & DeepProve, zkGPT, zkAgent \cite{zkagent} & inference
integrity \\
Attacks / negative results & Hollow-LLM \cite{hollowllm}, zk-attach
\cite{zkattach} & naming the gaps \\
Non-ZK industry baseline & \texttt{@brivora/verify}, VAOL & signed audit logs \\
\bottomrule
\end{tabular}
\end{table}

One of the five is the lineage this paper is actually in dialogue with. The
attacks-and-negative-results row names two gaps---Hollow-LLM's \emph{effort gap}
and zk-attach's \emph{source-integrity gap}---and both concern \emph{what the
proof system was asked to prove}. Proposition~\ref{thm:gap} concerns a third
gap on a different axis, \emph{where the proved string is delivered}, and we are
not aware of prior work that quantifies the binding/hiding tradeoff in the terms
of \S\ref{sec:binding}. That is why we state it as a bound with a stated
hypothesis rather than cite it.

\paragraph{zkAgent: orthogonal, not competing.} zkAgent \cite{zkagent} is the
closest contemporaneous system, and the comparison is easy to get wrong in both
directions. It proves that a provider executed the declared model and tools; PoP
proves that a response satisfies a public policy. PoP proves nothing about model
computation, and zkAgent proves nothing about policy compliance, so the two
compose: their conjunction covers ``inference was executed \emph{and} the policy
held''. Table~\ref{tab:zkagent} places the reported figures side by side.

\begin{table}[t]
\centering
\caption{PoP against zkAgent. The two prove different obligations, so the rows
are a statement of orthogonality, not a race.}
\label{tab:zkagent}
\tabfont
\setlength{\tabcolsep}{3pt}
\begin{tabular}{@{}L{0.22\columnwidth}L{0.38\columnwidth}L{0.30\columnwidth}@{}}
\toprule
& \textbf{zkAgent} \cite{zkagent} & \textbf{PoP} \\
\midrule
Obligation & provider executed the declared model and tools (inference integrity)
& response/trace satisfies policy $\pi$ (compliance) \\
Depends on the model & yes (weights, arithmetisation) & \textbf{no} \\
Cost grows with & model size, sequence length & response length $\times$ rule
count \\
Content privacy & none & committed, selective disclosure \\
Hardware reported & 32 cores / 512\,GB & 24 cores / 12\,GB \\
Proving & 194.4\,s (Shout, $T{=}512$) & 118.9--172.5\,s (200--10k chars) \\
Proof size & 96\,MiB (Shout); 3.1\,MiB (LogUp) & $\approx$2.7\,MiB \\
Verification & 0.42\,s (Shout); 0.038\,s (LogUp) & $\approx$0.090\,s \\
\bottomrule
\end{tabular}
\end{table}

We deliberately do not claim that PoP is faster, and the temptation to do so is
worth naming because the raw numbers invite it. The correct statement is
narrower, and we make only it: \emph{for the compliance obligation}, PoP
completes in a proving time of the same order as zkAgent's, with a smaller memory
footprint and a comparable proof size, without touching the model---and it adds
committed content privacy that zkAgent does not provide. Two caveats must travel
with those figures. The hardware is not matched (24 cores and 12\,GB against 32
cores and 512\,GB), and PoP's ceiling is memory-bound rather than cycle-bound, so
its feasible region is a two-step staircase rather than a curve.
\S\ref{sec:eval} reports the measurements and their limits.

\need{⚠️ tab:zkagent 里 PoP 那一列全部取自 SP1 \texttt{core}（非零知识）档，
机器为 24 核 / 12\,GB WSL2；zkAgent 一列为 ePrint 2026/199 的报告值（32 核 / 512\,GB）。
数据源：\texttt{bench/comparison\_zkagent.md}。该文件自身警告：
\emph{不得}宣称「PoP 比 zkAgent 快」——正确口径是「对『合规』这一义务…
同量级证明时间 / 更小内存」。groth16 档（WP1）未打通前，本表不得出现 ZK 档数字。}

The one influence we should name plainly, rather than leave to resemblance:
zk-regex's ``match offline, verify the path in circuit'' pattern is the shape of
our NFA contract, and verifiable document redaction \cite{vdr} inspired the bit
selector behind \texttt{mask\_covered}. The borrowing stops at the circuit
technique and does not carry the security argument across: VDR assumes a
\emph{signed original document} exists, whereas in PoP's private mode both the
original and the redacted version are produced by the operator, so there is no
trusted original to appeal to. We use the mechanism, not the proof.

\fixme{引文纪律：除 Hollow-LLM \cite{hollowllm} / zkGen \cite{zkgen} /
zkAgent \cite{zkagent} / PCO \cite{pco} 四条外，其余引文必须逐条打开原文
核实后再进 bibliography（docs/tdsc-fixlist.md \S2、WP8）。
\textbf{骨架原先把「ePrint 2024/813」列进已核实四条，那是错的}，已剔除：
\textbf{绝不能再引 IACR ePrint 2024/813 作为 VDR}——那是
Kiraz/Larraia/Vaughan, ``How to Redact the Bitcoin Backbone Protocol''，
与 Verifiable Document Redacting 无关。见 fixlist \S A。
⚠️ 口径：WITNESS 参考实现为 Apache-2.0（并非闭源），故「开源」不构成差异化。}

% =====================================================================
\section{Discussion and Limitations}
\label{sec:discussion}

Everything above was kept deliberately narrow, and this section says how narrow.
It covers the four places where a reader is most likely to round our claims up:
the caliber of the privacy statement (\S\ref{sec:disc-privacy}), where
zero-knowledge actually pays (\S\ref{sec:disc-zkmap}), availability against an
adversarial response (\S\ref{sec:disc-dos}), coverage
(\S\ref{sec:disc-coverage}), and composition with an inference proof
(\S\ref{sec:disc-composition}). It closes with the engineering limitations that
remain (\S\ref{sec:disc-rest}).

\subsection{The caliber of the privacy claims}
\label{sec:disc-privacy}

\textbf{``Content privacy'' means \emph{no plaintext of $T$ in the public
values}, not \emph{content is unrecoverable}.} Two independent reasons force
this reading---one at the proof-system level, one at the design level---and the
weaker reading is the one our own measurements support.

\emph{Proof-system level.} The default proving tier is SP1's \texttt{core}, and
SP1's published security model states explicitly that individual STARK proofs do
not currently satisfy the zero-knowledge property; the source agrees, with no
blinding anywhere in the proof stack and the shard proof carrying the trace
Merkle root and trace openings in the clear.%
\need{⚠️ 引文债：「SP1 公开安全模型承认其 STARK 证明当前不满足零知识」是
\emph{外部}断言，出处（SP1 官方安全模型文档的具体章节/URL）尚未核实到可引的
一级来源，故此处只陈述、不给 \cite{}（WP8）。若核不到原文，须撤回该句，
只保留我们自己的源码审计结论那半句。}
Soundness (G1) therefore holds
while the proof artefact itself does not hide the witness. The
\texttt{gr\-oth16} and \texttt{plonk} tiers wrap the inner STARK proof as a
private witness---the only mode in which hiding is even possible---but that is a
wrapper-level claim, it is not post-quantum, it needs $\ge 16$\,GB, and we did
not measure it (\S\ref{sec:eval-proofs}).

\emph{Design level, and the more fundamental one.} The public values
\texttt{response\_binding} and \texttt{response\_commitment} are both
\emph{public, offline-recomputable functions of $T$}. This is not an oversight:
\texttt{response\_binding} \emph{must} be publicly recomputable, because that is
how a verifier independently checks that the delivered $T'$ is the $T$ that was
proved---the swap-response defence. But the same property is a guess-and-verify
oracle for $T$.

It is worth being precise about the attack, because the obvious version of the
argument is wrong. ``A one-line response carries only $10^2$--$10^3$ bits of
entropy, far below 256, so enumeration is feasible'' does \emph{not} follow:
$2^{100}$ is as unenumerable as $2^{256}$. The real attack is
\emph{re-execution} (Corollary \ref{cor:reexec}). The adversary's side
information is not the response space but the tuple (prompt, model, decoding
parameters) (Assumption \ref{as:Z}); under greedy or low-temperature decoding
the response is nearly determined, so $k = -\log_2 \mathsf{Guess}(T \mid Z)
\approx 0$ (Definition \ref{def:k}) and the attack costs \emph{one inference},
independent of $|T|$. Under the two premises that a verifier independently
recomputes the binding and that $k < \log q + \log(1/\varepsilon)$, response
binding and response-content hiding are mutually exclusive.

We therefore restate (G2) as what is actually established: \emph{in private mode
the public values contain no plaintext of $T$}---which the leak experiment
confirms directly---and not as content irrecoverability. The consequence for the
comparison in \S\ref{sec:related} is that the ``content privacy'' cells on
\emph{both} sides must not be read as ``content is unrecoverable.'' This is a
bound on our own system first.

\subsection{Where zero-knowledge pays, and where it does not}
\label{sec:disc-zkmap}

Section \ref{sec:eval-baseline} measured the non-ZK alternative---a signed
Merkle log---and everything in it is milliseconds. A proof therefore has to earn
its cost against what a signature cannot do, and our measurements locate that
boundary instead of asserting it.

The map has two axes: is the response secret, and is the policy secret?
Zero-knowledge has positive value only in the corners where at least one of them
is unknown to the verifier. In the transparent corner---public response, public
policy---a signed log certifies exactly the same thing at three to four orders
of magnitude less cost, and adding a proof is redundancy rather than rigour. The
familiar criticism that ``a signature would have done'' is correct in that one
corner and wrong in the other three; the decision map is what turns the
criticism into a measured design conclusion.

One dependency belongs in the open. The zero-knowledge property of a certificate
is a property of the \emph{wrapper}, not of the system as a whole, and the
default tier does not have it (\S\ref{sec:disc-privacy}). If \texttt{gr\-oth16}
does not become a delivery tier, we do not claim the artefact is
zero-knowledge, and the word must leave the title.

\subsection{Availability: the unanchored-pattern cost attack}
\label{sec:disc-dos}

A regular expression with an unanchored pattern---one with no left anchor, such
as \texttt{.*foo}---gives the matcher no way to reject a prefix of the response
early: every position stays a live start position until the whole string has
been read. The prover's work therefore scales with the response, and a party who
\emph{chooses} the response can drive it quadratically. In the agent setting
that party is often the agent itself, or whoever controls its output.

Two things are worth separating. This is a denial of service against the
\emph{prover}, not an attack on soundness: no false certificate is produced, and
the cost is simply borne by the honest side. And it is specific to proving
rather than to matching, because the trace---not merely the verdict---is what
gets proved. Mitigations exist and are partly in place: a Pike VM (already
implemented) simulates the NFA in linear time, committing to the response length
up front removes the adversary's freedom to inflate it, and chunked matching
bounds the per-chunk work. Prior work on zero-knowledge regular expressions
imposes the DFA restriction but does not analyse denial of service in a zkVM
setting, which is where this differs.

\need{⚠️ 作者决策未定：骨架 :649--651 把这条写成一条定理（$\Omega(n^2)$ 证明
工作量），但证明没写。需定：写成定理（并补证明），还是降为观察（Observation）。
本稿先按观察陈述、不给定量断言，以免留下一条无证明的定理。}

\subsection{Coverage: a provider can simply withhold certificates}
\label{sec:disc-coverage}

Nothing in the cryptographic layer can compel a provider to issue a certificate
at all. A provider that never produces one is indistinguishable, from the
verifier's side, from a provider with nothing to certify. We treat this as an
open problem rather than a defect of the construction: without trusted hardware
there is no anchor that forces issuance, and the mechanisms that could---%
sampling audits, economic penalties, disclosure obligations---live above the
cryptographic layer and outside the scope of this paper. Naming it as an open
problem is more useful than proposing a mitigation we cannot evaluate.

\subsection{Why composition alone does not close the gap}
\label{sec:disc-composition}

The natural next step is to compose: prove that the inference produced $T$ from
the request, conjoin that with a proof that $\pi(T)$ holds, and conclude that
the agent \emph{sent} a compliant response. We describe the mechanism, and we
record why the conclusion does not follow for free.

\emph{Composition by itself does not close the gap.} A model proof bound to
\emph{some} input shows only that \emph{some} input yields a compliant output;
the divergent-delivery gap of Proposition \ref{thm:gap} closes only when the
model proof is bound to the \emph{input the user submitted}, $x$, so that
$M(x) = T \wedge \pi(T)$ becomes a statement about the bytes actually sent. That
requirement is not an extra condition bolted on: it is exactly $f > 0$ in
Theorem \ref{thm:t3}, with the user-submitted $x$ playing the role of the $f$
fresh bits. The two halves of this paper meet here---the request-side theorem
names the condition the composition would have to satisfy, and the response-side
theorem is what makes that condition non-vacuous.

The composition contract is correspondingly short to state: the inference proof
must expose $T$ in its public values, and it must be bound to $x$.

\need{⚠️ 阻塞：真模型规模（$\ge 10^6$ 参数）与组合成本数据属 WP4 子问题② /
WP3 的 DeepProve 对照臂，需外部资源。代理规模下「组合成本由推理证明主导」
\emph{未被验证}（\texttt{bench/results/compose.md}），不得写入。
骨架 :641--642 的组合契约定义可在此处补。}

\subsection{Remaining limitations}
\label{sec:disc-rest}

\begin{itemize}
  \item \textbf{Three-party collusion is undetectable at the signature layer.}
    If the policy author, the operator, and the prover all collude, the
    certificate is indistinguishable from an honest one. The multi-prover
    construction partitions responsibility; it does not prevent agreement
    (\texttt{docs/security-model.md}).
  \item \textbf{Signing and key custody are not production-grade.} The on-chain
    path currently takes plaintext private-key parameters, the demo uses a
    public Anvil test key, and deployment is a local Anvil instance. Production
    signing (Ed25519/HSM) and key custody remain future work.
  \item \textbf{Semantic-rule model quality.} The delegated semantic rule works
    end to end, but the bundled model is a small demonstration proxy rather than
    a usable learned guardrail: ``$P(\text{harmful})$ is low'' is a statement
    about \emph{this} model. The companion proof's public instance also contains
    $\mathsf{encode}(T)$, so semantic rules are \emph{public-mode only}.
  \item \textbf{The inference proof is faithful, not good.} The inference half
    is a deterministic fixed-point MLP proxy whose weights are generated from a
    compile-time seed and baked into the program, hence committed by the
    verification key. It proves that \emph{this} graph computed \emph{this}
    output on \emph{this} response---not that the graph is any good. Swapping the
    proxy for a real prover touches only the prover; the composition layer does
    not move.
  \item \textbf{Verification remains $O(N)$.} There is no recursive aggregation,
    so a verifier pays per certificate. The supported regex subset and its ASCII
    semantics are bounded in the same way, and proof-cost optimisation (lookup
    arguments, parallelism, precomputation) is unexplored.
  \item \textbf{Slices share one guest.} Multiple provers hold distinct keys,
    but that is a division of responsibility rather than a per-role circuit: the
    repository does not build one circuit per role.
\end{itemize}

% =====================================================================
\section{Conclusion}
\label{sec:conclusion}

This paper began from a small question---what does a compliance certificate
actually bind?---and the answer came out the same on both sides of the
interaction. A proof binds only the bits that were submitted into it \emph{and}
that the adversary could not select in advance. On the response side those bits
are the entropy $k$ of $T$ given everything the adversary already knows, and no
amount of computation buys more of them; on the request side they are the $f$
bits of freshness the receiver contributes, and a certificate without them is
replayable without limit. Theorem \ref{thm:t1} and Theorem \ref{thm:t3} are two
faces of that one statement, and Proposition \ref{thm:gap}---a gap we state
against our own design---is what makes it bite.

PoP is the constructive instance of that analysis, not a claim to have invented
a new proof system. It composes an existing zkVM, an existing policy compiler,
and a signature layer into a certificate that a verifier holding only public
values can check without trusting hardware, and it is explicit about what such a
certificate does not do: it does not make the response unrecoverable, it does not
hide the witness at the default proving tier, and it cannot compel a reluctant
provider to issue anything at all. Those are limits of the problem, stated as
limits, rather than of one implementation.

The design conclusions are accordingly modest, and we hope reusable: publish the
leak profile of a certificate and audit it like a disclosure policy; let a
privacy claim be only as strong as its arithmetic supports; and put the
freshness where the binding is. A certificate is worth exactly as much as the
bits that went into it.

% =====================================================================
% ⚠️ IEEEtran 的 \appendix 是「单个附录」命令：它把 \section 重定义成
% \@IEEEdestroythesectionargument，直接丢弃 \section{...} 的参数（只打 warning）。
% 多个附录必须用 \appendices（复数），它把 section 计数器置 -1，
% 之后每个 \section{X} 才依次成为 A、B、……  （2026-10-12 实测：
% 用 \appendix 时 A/B 两个附录的 \ref 都解析成 "A"，且 title 字段退化成 \appendixname，
% 锚点退化成 section*.1；「0 undefined ref」检查抓不到这种错。）
\appendices
\section{Proofs}
\label{app:proofs}

\begin{proof}[Proof of Theorem \ref{thm:t1}]
\emph{(H).} Fix $z$ and let the adversary's query set be
$Q = \{q_1, \dots, q_q\}$. Decompose by first hit,
\[
\Pr[\exists i : q_i = T] \;=\; \sum_i \Pr\!\left[q_i = T \wedge \neg E_{<i}\right],
\]
where $E_{<i} \triangleq \{\exists j < i : q_j = T\}$.
Conditioned on $\neg E_{<i}$, every oracle answer $a_1,\dots,a_{i-1}$ already
revealed is uniform and independent of $T$ (no queried point equals $T$, and $H$
is a fresh random function), and $\beta(T)$ itself is uniform and independent of
$T$. Hence $q_i$---a deterministic function of $(z, \beta(T), a_{<i})$---is
\emph{independent of $T$}, so
\begin{multline*}
\Pr[q_i = T \mid Z = z]
= \sum_t \Pr[q_i = t]\,\Pr[T = t \mid Z = z] \\
\le \max_t \Pr[T = t \mid Z = z] = 2^{-k_z},
\end{multline*}
giving $\Pr[\exists i : q_i = T \mid z] \le q\,2^{-k_z}$. Taking the expectation
over $z$ and using $\mathbb{E}_z[2^{-k_z}] = \mathsf{Guess}(T \mid Z) = 2^{-k}$
from Definition \ref{def:k} yields $q \cdot 2^{-k}$. An adversary that outputs
an unqueried value wins with probability at most $2^{-k_z} \le q\,2^{-k_z}$,
which the same bound absorbs. The term $q \cdot 2^{-\lambda}$ charges the
probability that a query hits the uniform value $\beta(T)$ by chance.

\emph{(L).} Immediate from Definition \ref{def:k} after conditioning on
$\beta(T)$; no hypothesis on $\beta$ is used.
\end{proof}

\need{Thm \ref{thm:t1}(L) 无闭式界，只给测量口径（\S\ref{sec:eval}）。
三分法已降级为分类 Rem \ref{rem:taxonomy}，无需证明。
原 T2「分离定理」已\emph{停用}（见 Rem \ref{rem:designated}）—— 不要为它写附录。}

\need{\textbf{待补}：把 Thm \ref{thm:t1}(H) 的随机预言机模型换成标准模型，
或明确声明本文只在 ROM 下主张。当前证明依赖 $H$ 是新鲜随机函数这一条；
实践中的 SHA-256 只给计算意义下的近似，故正文措辞必须写成
\emph{heuristic}，不得写成信息论界。}

\section{Reproduction}
\label{app:repro}

The full walkthrough---environment, a single compliance proof, cross-checking
against the reference layer, private mode, certificate issuance, the one-command
demo, and third-party verification---is maintained in \texttt{docs/reproduce.md}
together with the exact commands and their expected output. This appendix records
only what a reader needs in order to judge the claims of \S\ref{sec:eval}.

\paragraph{Environment.} Ubuntu or WSL2 on x86-64, no GPU. The Rust/SP1 side
needs Rust stable, \texttt{cargo-prove}~v6.7.0, Go $\ge$\,1.24, and the usual
native toolchain; the reference layer is pure Python~3.10+ and runs on its own.
Memory, not time, is the binding constraint: the CPU prover carries a fixed
floor of approximately 10.15\,GiB regardless of the trace
(\S\ref{sec:eval-proofs}), and 12\,GB suffices only for small policies.

\paragraph{What reproduces without the prover.} The reference layer, the policy
compiler, and the whole test suite run in seconds and need no Rust. This is
enough to reproduce the non-ZK baseline of \S\ref{sec:eval-baseline}, to inspect
the certificate and inclusion proofs, and to check the cross-layer contract---%
that the Python reference implementation and the zkVM guest consume the same
compiled policy bytes.

\paragraph{What needs the prover.} Every timing, size, and memory figure in
\S\ref{sec:eval-cycles}--\S\ref{sec:eval-proofs}, and the end-to-end result of
\S\ref{sec:eval-e2e}, requires a real proof. The measurements were taken on a
24-core CPU machine with 12\,GB of RAM, which is what places the ceiling reported
there: the successes and the two out-of-memory configurations in
Table~\ref{tab:proofs} are both properties of that machine.

\paragraph{Where each number comes from.} Every measurement in \S\ref{sec:eval}
is the content of a checked-in result file with a machine-readable companion:
\texttt{bench/results/} holds \texttt{cycles}, \texttt{ablation},
\texttt{proofs}, \texttt{baseline}, \texttt{verify}, \texttt{semantic},
\texttt{streaming}, \texttt{compose}, and \texttt{prover\_knobs}, each as a
\texttt{.json} alongside its \texttt{.md}. The reproduction guide records the
commit, the hardware, and the prover binary digest for each run, so that a
reported figure can be traced to one specific execution rather than to a state
of the tree.

\need{⚠️ 数字漂移（本轮实测发现，需作者定夺）：正文 §2.2 写「801 tests / 15 skipped」，
其依据是实测（\texttt{docs/dev-plan.md:506} 的落地记录）\emph{以及}本轮对
\texttt{tests/} 的静态清点（\texttt{def test\_} 计 801 个）——\textbf{801 是对的}。
但 \texttt{docs/reproduce.md:31,428}、\texttt{docs/security-model.md:829}、
\texttt{docs/development.md:93,290}、\texttt{docs/modules/08-tests-bench.md:11,58,504,526}、
\texttt{docs/modules/README.md:94,161,201} 仍写 \textbf{783}（陈旧值，
\texttt{docs/report/advisor-report.md:586} 已把这条「文档数字漂移」登记在案）。
本附录既然指向 \texttt{reproduce.md}，照抄就会让复现者跑出 801 却以为环境坏了——
\textbf{这是真实的复现缺陷}。按仓库自己的「过期就标注、不静默改」纪律，
本稿不擅自改那些文档，只在此记债。作者需定：把五处文档回填为 801
（\texttt{advisor-report.md:599} 的建议），还是在附录里并列两套口径。}

% =====================================================================
\begin{thebibliography}{99}

\bibitem{hollowllm}
C.~Gong, B.~Liu, and M.~Li, ``Hollow-LLM Attack: Computationally Trivial
Weights in Zero-Knowledge Verification of LLM Inference,'' in \emph{IEEE S\&P},
2026, pp.~4339--4355, DOI 10.1109/SP63933.2026.00258.

\bibitem{zkattach}
\need{zk-attach, arXiv:2608.30083 --- 打开原文核实作者后再引}

\bibitem{zkgen}
J.~Lauinger, J.~Ernstberger, and S.~Steinhorst, ``zkGen: Policy-to-Circuit
Transpiler,'' in \emph{IEEE ICBC}, 2024, pp.~620--624,
DOI 10.1109/ICBC59979.2024.10634440.

\bibitem{pco}
\need{PCO, IACR ePrint 2026/994 --- 打开原文核实作者（Dolev/Yung 等）后再引}

\bibitem{zkagent}
L.~Wang, H.~Lou, C.~Li, Y.~Yu, and Y.~Hu, ``zkAgent: Verifiable LLM Agent
Execution via One-Shot Transcript Proofs,'' IACR ePrint 2026/199.

\bibitem{hiddencred}
D.~Balfanz and J.~Schweitzer, ``Hidden credentials,''
\need{核实出处：2002 年 unpublished manuscript / 2003 年会议版本，二者择一并核准页码}.

\bibitem{handshake}
\need{核实作者列表与页码（首发 S&P 2003, pp.~139--150）}D.~Balfanz, G.~Durfee,
N.~Shankar, D.~K.~Smetters, J.~Staddon, and H.-C.~Wong, ``Secret handshakes from
pairing-based key agreements,'' in \emph{IEEE S\&P}, 2003.

\bibitem{xacmlpriv}
\need{核实作者与出处。骨架点名的方向是「XACML 策略的隐私保护访问控制」，
候选：C.~A.~Ardagna 等关于 XACML 隐私的系列工作。\textbf{未核实前不得写成
具体条目}。}

\bibitem{sw05}
\need{核实页码（LNCS 3494, pp.~457--473）}A.~Sahai and B.~Waters, ``Fuzzy
identity-based encryption,'' in \emph{EUROCRYPT}, 2005.

\bibitem{ksw08}
\need{核实页码（LNCS 4965, pp.~515--534）}J.~Katz, A.~Sahai, and B.~Waters,
``Predicate encryption supporting disjunctions, polynomial equations, and inner
products,'' in \emph{EUROCRYPT}, 2008.

\bibitem{vdr}
\need{Verifiable Document Redacting, Springer, DOI 10.1007/978-3-319-66402-6_20
--- 核实作者与年份。\textbf{绝不能再引 IACR ePrint 2024/813}（那是
Kiraz/Larraia/Vaughan, "How to Redact the Bitcoin Backbone Protocol"）。
见 docs/tdsc-fixlist.md \S A。}

\end{thebibliography}

\end{document}
